% !TeX program = pdflatex
% !BIB program = biber
\documentclass[12pt,a4paper]{article}

% =====================================================
% ROUTTECH - ENTREGA 3 (2A)
% PORTADA, CONTROL DOCUMENTAL E INTRODUCCION
% Compatible con pdfLaTeX + Biber en VS Code y Overleaf.
% Las citas se numeran por orden de primera aparicion mediante el estilo IEEE.
% =====================================================

\usepackage[utf8]{inputenc}
\DeclareUnicodeCharacter{011F}{{\fontfamily{cmss}\selectfont \u{g}}}
\usepackage[T1]{fontenc}
\usepackage{textcomp}
\usepackage[spanish,es-nodecimaldot]{babel}

% Tipografia Gill Sans libre compatible con pdfLaTeX
\usepackage[sfdefault]{gillius2}
\renewcommand{\familydefault}{\sfdefault}

\usepackage[
  a4paper,
  includeheadfoot,
  top=1.65cm,
  bottom=1.45cm,
  left=2.5cm,
  right=2.5cm,
  headheight=23pt,
  headsep=0.65cm,
  footskip=0.80cm
]{geometry}

\usepackage[table]{xcolor}
\usepackage{array,tabularx,booktabs,longtable}
\usepackage{enumitem}
\usepackage{fancyhdr}
\usepackage{microtype}
\usepackage{setspace}
\usepackage{titlesec}
\usepackage{hyperref}
\usepackage{xurl}
\usepackage{ragged2e}
\usepackage{graphicx}
\usepackage{caption}
\usepackage{csquotes}
\usepackage{float}
\usepackage{pdflscape}
\usepackage{seqsplit}
\usepackage[most]{tcolorbox}
\usepackage{tocloft}
\usepackage[backend=biber,style=ieee,sorting=none,maxbibnames=99]{biblatex}
\addbibresource{referencias.bib}

% =====================================================
% DATOS DEL DOCUMENTO
% =====================================================
\newcommand{\Proyecto}{RoutTech: Sistema de Movilidad Colaborativa para la Comunidad Universitaria de la UTEQ}
\newcommand{\Equipo}{Equipo CN}
\newcommand{\VersionDocumento}{2A-v1.0}
\newcommand{\FechaDocumento}{31 de julio de 2026}
\newcommand{\Repositorio}{https://github.com/AnderJM3/RoutTech_ISR401}
\newcommand{\Mayra}{Córdova Carreño Mayra Lucila}
\newcommand{\Anderson}{Naranjo Flores Anderson Jeampiere}
\newcommand{\Docente}{Ing. Gleiston Cicerón Guerrero Ulloa, PhD.}

% =====================================================
% PALETA SOBRIA Y PROFESIONAL
% =====================================================
\definecolor{DocNavy}{HTML}{20364D}
\definecolor{DocBlue}{HTML}{3F5F7A}
\definecolor{DocGray}{HTML}{59636E}
\definecolor{DocLight}{HTML}{F1F3F5}
\definecolor{DocLighter}{HTML}{F8F9FA}
\definecolor{DocRule}{HTML}{A9B1B9}
\definecolor{DocText}{HTML}{1F2328}

% =====================================================
% CONFIGURACION GENERAL
% =====================================================
\hypersetup{
  colorlinks=true,
  linkcolor=DocNavy,
  urlcolor=DocBlue,
  citecolor=DocBlue,
  pdftitle={ERS/SRS RoutTech - Entrega 3 (2A)},
  pdfauthor={Córdova Carreño Mayra Lucila; Naranjo Flores Anderson Jeampiere}
}
\urlstyle{same}

\pagestyle{fancy}
\fancyhf{}
\fancyhead[C]{%
  \footnotesize
  RoutTech (Movilidad Colaborativa) -- ERS/SRS
  \textbar\ UTEQ \textbar\ \Equipo
}
\fancyfoot[C]{\small\thepage}
\renewcommand{\headrulewidth}{0.45pt}
\renewcommand{\footrulewidth}{0pt}
\renewcommand{\headrule}{%
  \hbox to\headwidth{\color{DocRule}\leaders\hrule height \headrulewidth\hfill}
}

\fancypagestyle{plain}{%
  \fancyhf{}%
  \fancyhead[C]{\footnotesize RoutTech (Movilidad Colaborativa) -- ERS/SRS \textbar\ UTEQ \textbar\ \Equipo}%
  \fancyfoot[C]{\small\thepage}%
  \renewcommand{\headrulewidth}{0.45pt}%
  \renewcommand{\footrulewidth}{0pt}%
}

\titleformat{\section}
  {\large\bfseries\color{DocNavy}}
  {\thesection.}{0.55em}{}
\titleformat{\subsection}
  {\normalsize\bfseries\color{DocNavy}}
  {\thesubsection.}{0.50em}{}
\titleformat{\subsubsection}
  {\normalsize\bfseries\color{DocGray}}
  {\thesubsubsection.}{0.50em}{}
\titlespacing*{\section}{0pt}{1.05em}{0.45em}
\titlespacing*{\subsection}{0pt}{0.80em}{0.30em}
\titlespacing*{\subsubsection}{0pt}{0.70em}{0.25em}

% Sin sangria; separacion breve y uniforme entre parrafos
\setlength{\parindent}{0pt}
\setlength{\JustifyingParindent}{0pt}
\setlength{\parskip}{0.48em}
\setstretch{1.10}
\setlist[itemize]{leftmargin=1.45em,itemsep=1.5pt,topsep=2.5pt,parsep=0pt}
\setlist[enumerate]{leftmargin=1.55em,itemsep=1.5pt,topsep=2.5pt,parsep=0pt}

\renewcommand{\arraystretch}{1.22}
\setlength{\tabcolsep}{5pt}
\arrayrulecolor{DocRule}
\newcolumntype{Y}{>{\RaggedRight\arraybackslash}X}
\newcolumntype{C}[1]{>{\centering\arraybackslash}p{#1}}
\newcolumntype{L}[1]{>{\RaggedRight\arraybackslash}p{#1}}
\captionsetup{font=small,labelfont=bf,labelsep=colon,justification=centering}

% Inserta la figura si existe; mientras tanto muestra un espacio identificado.
\newcommand{\FiguraDocumento}[4]{%
  \begin{figure}[H]
    \centering
    \IfFileExists{#1}{%
      \includegraphics[width=#2\textwidth,height=0.78\textheight,keepaspectratio]{#1}%
    }{%
      \fcolorbox{DocRule}{DocLighter}{%
        \parbox[c][5.1cm][c]{0.88\textwidth}{%
          \centering
          \textbf{Imagen pendiente de insertar}\\[5pt]
          \scriptsize Archivo esperado:\\[-1pt]\nolinkurl{#1}%
        }%
      }%
    }%
    \caption{#3}
    \label{#4}
  \end{figure}
}

\addto\captionsspanish{%
  \renewcommand{\contentsname}{ÍNDICE GENERAL}%
  \renewcommand{\listtablename}{ÍNDICE DE TABLAS}%
  \renewcommand{\listfigurename}{ÍNDICE DE FIGURAS}%
  \renewcommand{\tablename}{Tabla}%
  \renewcommand{\figurename}{Figura}%
}
\providecommand*{\tableautorefname}{Tabla}
\providecommand*{\figureautorefname}{Figura}
\renewcommand{\cfttoctitlefont}{\hfill\large\bfseries\color{DocNavy}}
\renewcommand{\cftaftertoctitle}{\hfill}
\renewcommand{\cftlottitlefont}{\hfill\large\bfseries\color{DocNavy}}
\renewcommand{\cftafterlottitle}{\hfill}
\renewcommand{\cftloftitlefont}{\hfill\large\bfseries\color{DocNavy}}
\renewcommand{\cftafterloftitle}{\hfill}
\setlength{\cftbeforesecskip}{2pt}


% =====================================================
% FICHAS DE REQUISITOS FUNCIONALES
% =====================================================
\newtcolorbox{RFBox}[2][]{
  enhanced,
  breakable,
  colback=white,
  colframe=DocNavy,
  colbacktitle=DocNavy,
  coltitle=white,
  fonttitle=\bfseries,
  title={#2},
  boxrule=0.65pt,
  arc=1.2mm,
  left=2.0mm,
  right=2.0mm,
  top=1.4mm,
  bottom=1.4mm,
  before skip=5pt,
  after skip=6pt,
  #1
}

\newcommand{\CampoRF}[2]{\textbf{#1} & #2\\}

% =====================================================
% FICHAS DE REQUISITOS NO FUNCIONALES
% =====================================================
\newtcolorbox{RNFBox}[2][]{
  enhanced,
  breakable,
  colback=white,
  colframe=DocBlue,
  colbacktitle=DocBlue,
  coltitle=white,
  fonttitle=\bfseries,
  title={#2},
  boxrule=0.65pt,
  arc=1.2mm,
  left=2.0mm,
  right=2.0mm,
  top=1.4mm,
  bottom=1.4mm,
  before skip=5pt,
  after skip=6pt,
  #1
}

\newcommand{\CampoRNF}[2]{\textbf{#1} & #2\\}


% =====================================================
% HISTORIAS DE USUARIO Y CRITERIOS DE ACEPTACION
% =====================================================
\newtcolorbox{HUBox}[2][]{
  enhanced,
  breakable,
  colback=white,
  colframe=DocNavy,
  colbacktitle=DocNavy,
  coltitle=white,
  fonttitle=\bfseries,
  title={#2},
  boxrule=0.65pt,
  arc=1.2mm,
  left=2.0mm,
  right=2.0mm,
  top=1.4mm,
  bottom=1.4mm,
  before skip=5pt,
  after skip=6pt,
  #1
}

\newtcolorbox{GherkinBox}[2][]{
  enhanced,
  breakable,
  colback=DocLighter,
  colframe=DocBlue,
  colbacktitle=DocBlue,
  coltitle=white,
  fonttitle=\bfseries,
  title={#2},
  boxrule=0.55pt,
  arc=1.0mm,
  left=2.2mm,
  right=2.2mm,
  top=1.5mm,
  bottom=1.5mm,
  before skip=4pt,
  after skip=5pt,
  #1
}

\newcommand{\CampoHU}[2]{\textbf{#1} & #2\\}


% =====================================================
% FICHAS DE CASOS DE USO
% =====================================================
\newtcolorbox{CUBox}[2][]{
  enhanced,
  breakable,
  colback=white,
  colframe=DocNavy,
  colbacktitle=DocNavy,
  coltitle=white,
  fonttitle=\bfseries,
  title={#2},
  boxrule=0.65pt,
  arc=1.2mm,
  left=2.0mm,
  right=2.0mm,
  top=1.4mm,
  bottom=1.4mm,
  before skip=5pt,
  after skip=7pt,
  #1
}

\newcommand{\CampoCU}[2]{\textbf{#1} & #2\\}

\begin{document}
\color{DocText}
\justifying

% =====================================================
% PORTADA
% =====================================================
\thispagestyle{empty}

\begin{center}
  {\scriptsize
  RoutTech (Movilidad Colaborativa) -- ERS/SRS
  \textbar\ UTEQ \textbar\ \Equipo}
\end{center}

\vspace{0.85cm}

\begin{center}
  {\normalsize\bfseries
  UNIVERSIDAD TÉCNICA ESTATAL DE QUEVEDO\\
  FACULTAD DE CIENCIAS DE LA COMPUTACIÓN\\
  INGENIERÍA DE SOFTWARE}
\end{center}

\vspace{2.35cm}

\begin{center}
  {\LARGE\bfseries\color{DocNavy}
  RoutTech: Sistema de Movilidad Colaborativa para la\\
  Comunidad Universitaria de la UTEQ}

  \vspace{0.55cm}
  {\normalsize\bfseries\color{DocGray}
  ENTREGA 3 (2A): ESPECIFICACIÓN DE REQUISITOS\\
  COMPLETA + COMPONENTE EMPÍRICO}
\end{center}

\vspace{0.38cm}
\noindent\color{DocNavy}\rule{\textwidth}{0.75pt}\color{DocText}
\vspace{0.72cm}

\begin{center}
\small
\begin{tabularx}{\textwidth}{|>{\centering\arraybackslash\columncolor{DocLight}}p{0.37\textwidth}|>{\centering\arraybackslash}X|}
\hline
\textbf{Integrantes:} &
\begin{minipage}[c]{\linewidth}
\centering\vspace{1.8mm}
\Mayra\\[1mm]
\Anderson
\vspace{1.8mm}
\end{minipage}\\
\hline
\textbf{Docente:} & \Docente\\
\hline
\textbf{Asignatura:} & Ingeniería de Requerimientos\\
\hline
\textbf{Curso y paralelo:} & Cuarto nivel ``B'' -- Software\\
\hline
\textbf{Fecha:} & \FechaDocumento\\
\hline
\textbf{Versión:} & \VersionDocumento\ -- \Equipo\\
\hline
\textbf{Repositorio GitHub:} &
\begin{minipage}[c]{\linewidth}
\centering\vspace{1.3mm}
{\scriptsize
\href{\Repositorio}{\nolinkurl{\Repositorio}}}
\vspace{1.3mm}
\end{minipage}\\
\hline
\end{tabularx}
\end{center}

\vfill
\begin{center}
  {\normalsize\bfseries
  Quevedo -- Ecuador\\
  PPA 2026--2027}
\end{center}
\clearpage

% =====================================================
% CONTROL DOCUMENTAL
% =====================================================
\section*{Historial de versiones}
\addcontentsline{toc}{section}{Historial de versiones}

\begin{table}[H]
\centering
\caption{Historial de versiones del ERS/SRS}
\label{tab:historial-versiones}
\scriptsize
\rowcolors{2}{DocLight}{white}
\begin{tabularx}{\textwidth}{C{1.15cm} C{1.55cm} C{2.25cm} Y}
\rowcolor{DocNavy}
\textcolor{white}{\textbf{Versión}} &
\textcolor{white}{\textbf{Fecha}} &
\textcolor{white}{\textbf{Autoría}} &
\textcolor{white}{\textbf{Descripción del cambio}}\\
0.1 & 24/05/2026 & Equipo CN & Creación inicial del documento ERS/SRS del Proyecto Fin de Curso.\\
0.5 & 31/05/2026 & Equipo CN & Aplicación de técnicas de elicitación: entrevista estructurada y Design Thinking.\\
0.8 & 03/06/2026 & Equipo CN & Inclusión inicial de modelado UML, priorización MoSCoW y trazabilidad.\\
1.0 & 18/06/2026 & Equipo CN & Primera versión consolidada correspondiente a la Entrega 1B.\\
1.1 & 11/07/2026 & Equipo CN & Construcción de la estructura del ERS/SRS según ISO/IEC/IEEE 29148:2018.\\
2A-v1.0 & 31/07/2026 & Equipo CN & Nueva línea base de la Entrega 3 (2A), sustentada en 16 entrevistas recientes y 43 respuestas del cuestionario.\\
\end{tabularx}
\end{table}

\clearpage
\tableofcontents
\clearpage
\listoftables
\clearpage
\listoffigures
\clearpage

% =====================================================
% 1. INTRODUCCION
% =====================================================
\section{Introducción}

\subsection{Propósito del documento}

El presente documento constituye la Especificación de Requisitos de Software (ERS/SRS) de \textbf{RoutTech}, una propuesta de movilidad colaborativa para la comunidad universitaria de la Universidad Técnica Estatal de Quevedo. Su propósito es definir de forma completa, verificable y trazable las funciones, los atributos de calidad, las restricciones, las interfaces y los modelos que orientarán el diseño, la implementación, la validación y la evolución del producto. La estructura se alinea con ISO/IEC/IEEE 29148:2018 para la especificación y trazabilidad de requisitos, con ISO/IEC 25010:2023 para la calidad del producto, con SWEBOK v4.0 y con el material académico de la asignatura \cite{iso29148,iso25010,swebok2024,uteq2026}.

La Entrega 3 (2A) se desarrolla desde una nueva línea base empírica. Los requisitos no se copiarán automáticamente de versiones anteriores: cada afirmación sustantiva deberá rastrearse a una entrevista, una respuesta del cuestionario, una observación, un documento verificable, una obligación legal o una referencia bibliográfica. Este criterio permite controlar ambigüedades, evitar funciones sin respaldo y mantener una relación comprobable entre evidencia, requisito, historia de usuario, criterio de aceptación, caso de uso, componente y mockup.

\subsection{Alcance del producto}

RoutTech se plantea como una plataforma web y móvil destinada a facilitar la coordinación de viajes compartidos entre personas adultas y verificadas de la comunidad UTEQ. La movilidad compartida y el carpooling han sido estudiados como alternativas sostenibles para atender necesidades de traslado en comunidades universitarias \cite{alquhtani2022,lopezchila2025}. Otros estudios analizan sus efectos ambientales, las preferencias modales y el comportamiento de viaje de estudiantes universitarios \cite{arbelaez2023,bagdatli2022,etminani2023,tomas2021,zhu2022}. La definición progresiva del alcance, los requisitos y su validación seguirá además el enfoque sistemático de Ingeniería de Requisitos propuesto por Pohl \cite{pohl2010}.

El alcance funcional inicial comprende:

\begin{itemize}
  \item registro, autenticación y verificación de usuarios mediante información institucional;
  \item gestión de perfiles, roles, preferencias y vehículos;
  \item publicación, búsqueda, recomendación, reserva y cancelación de rutas;
  \item control de cupos y estados del viaje;
  \item coordinación de puntos de encuentro generales;
  \item comunicación y notificaciones relacionadas con reservas y cambios;
  \item calificaciones, reputación e incidencias;
  \item reportes e indicadores agregados para administración;
  \item explicaciones comprensibles sobre los factores que originan una recomendación de ruta.
\end{itemize}

La versión académica no procesará pagos, no administrará flotas institucionales, no sustituirá servicios de emergencia ni canales oficiales de la universidad y no publicará domicilios, coordenadas exactas o recorridos individuales identificables. El MVP se demostrará con datos ficticios o sintéticos. Cualquier tratamiento de información personal se sujetará a la Ley Orgánica de Protección de Datos Personales del Ecuador y a la documentación ética del proyecto \cite{lopdp2021}.

\subsection{Público destinatario}

El documento está dirigido al equipo investigador y de desarrollo, al docente responsable, a las personas que participen en la validación, a representantes institucionales y a quienes deban revisar o comprobar el comportamiento esperado del sistema. La ERS/SRS servirá como referencia común para:

\begin{itemize}
  \item comprender el límite del producto y sus actores;
  \item validar que los requisitos correspondan con necesidades reales;
  \item planificar el MVP y sus pruebas;
  \item controlar cambios y decisiones de diseño;
  \item verificar la cobertura de evidencias, requisitos, modelos y mockups.
\end{itemize}

\subsection{Fuentes de evidencia iniciales}

La línea base empírica disponible al inicio de esta versión está compuesta por 16 entrevistas recientes y 43 respuestas del cuestionario de movilidad universitaria. En esta sección se registra su existencia y su uso previsto; el catálogo detallado, los perfiles anonimizados, la codificación temática y la trazabilidad individual se incorporarán en los apartados correspondientes.

\begin{table}[H]
\centering
\caption{Fuentes de evidencia que sustentan la nueva línea base}
\label{tab:fuentes-evidencia}
\small
\rowcolors{2}{DocLight}{white}
\begin{tabularx}{\textwidth}{L{3.0cm} L{3.4cm} Y}
\rowcolor{DocNavy}
\textcolor{white}{\textbf{Identificador}} &
\textcolor{white}{\textbf{Fuente}} &
\textcolor{white}{\textbf{Uso previsto}}\\
\texttt{EV-ENT-001} a \texttt{EV-ENT-016} & 16 entrevistas semiestructuradas & Identificación cualitativa de problemas, necesidades, expectativas, riesgos, condiciones de confianza y funciones esperadas.\\
\texttt{EV-CUE-001} & Cuestionario con 43 respuestas & Descripción cuantitativa de medios de transporte, dificultades, percepción de seguridad, retrasos y disposición a utilizar una plataforma oficial.\\
\end{tabularx}
\end{table}

Las respuestas del cuestionario fueron registradas entre el 21 y el 27 de julio de 2026. De las 43 respuestas, 33 señalaron el bus público como medio habitual; 25 indicaron que sí utilizarían una plataforma oficial de viajes compartidos y 17 respondieron ``Tal vez''. La inseguridad fue identificada como el problema más grave por 27 participantes y 40 informaron que perdieron o llegaron tarde a una clase al menos una vez durante el último mes debido a dificultades de transporte.

Estos valores se utilizarán como estadística descriptiva de la muestra disponible y no como una estimación representativa de toda la población universitaria. Los requisitos finales se derivarán mediante el análisis conjunto de entrevistas, cuestionario, normativa y validación con stakeholders.

\subsection{Definiciones, acrónimos y abreviaturas}

\begin{longtable}{L{3.0cm} L{12.0cm}}
\caption{Definiciones, acrónimos y abreviaturas}\label{tab:glosario}\\
\rowcolor{DocNavy}
\textcolor{white}{\textbf{Término}} & \textcolor{white}{\textbf{Definición}}\\
\endfirsthead
\rowcolor{DocNavy}
\textcolor{white}{\textbf{Término}} & \textcolor{white}{\textbf{Definición}}\\
\endhead
ERS/SRS & Especificación de Requisitos de Software; documento que describe funciones, atributos de calidad, restricciones e interfaces del sistema.\\
RF & Requisito funcional; comportamiento o servicio que el sistema debe proporcionar.\\
RNF & Requisito no funcional; condición de calidad o propiedad medible que debe cumplir el sistema.\\
RD & Restricción de diseño; condición que limita las alternativas de implementación.\\
EV & Identificador de evidencia utilizado para rastrear el origen de una afirmación o requisito.\\
HU & Historia de usuario redactada desde la perspectiva de un rol y su beneficio esperado.\\
CA & Criterio de aceptación que permite comprobar una historia de usuario o requisito.\\
CU & Caso de uso que describe la interacción entre un actor y el sistema para alcanzar un objetivo.\\
MVP & Producto Mínimo Viable utilizado para demostrar y validar las funciones esenciales.\\
IA & Inteligencia artificial aplicada al componente de recomendación de rutas.\\
Explicabilidad & Capacidad del sistema para comunicar de forma comprensible los factores principales que sustentan una recomendación automática \cite{chazette2020,chazette2021,chazette2022}.\\
Stakeholder & Persona, grupo u organización que afecta al sistema o puede verse afectada por él.\\
MoSCoW & Técnica de priorización que clasifica requisitos como Must, Should, Could y Won't.\\
Kano & Modelo que clasifica características según su relación con la satisfacción del usuario.\\
WSJF & Método de priorización basado en la relación entre costo de retraso y tamaño del trabajo.\\
PICOC & Marco para formular preguntas de investigación considerando población, intervención, comparación, resultado y contexto.\\
LOPDP & Ley Orgánica de Protección de Datos Personales del Ecuador.\\
FAIR & Principios para que los datos sean localizables, accesibles, interoperables y reutilizables \cite{wilkinson2016}.\\
OSF & Open Science Framework; plataforma utilizada para registrar previamente el protocolo experimental.\\
\end{longtable}

\subsection{Referencias del documento}

Las fuentes bibliográficas se gestionan en el archivo \texttt{referencias.bib} y se citan automáticamente en estilo IEEE. La base incluye normas de Ingeniería de Requisitos y calidad de software, estudios de movilidad compartida y universitaria, protección de datos, explicabilidad, investigación empírica y gestión de datos FAIR. La versión final deberá citar al menos 25 fuentes primarias y un mínimo de 10 publicaciones de los últimos tres años, conforme a la guía de la Entrega 3 (2A).

\subsection{Visión general del documento}

La ERS/SRS completa se organizará de la siguiente manera:

\begin{itemize}
  \item \textbf{Sección 1:} introducción, alcance, evidencia inicial, glosario y referencias;
  \item \textbf{Sección 2:} perspectiva del producto, límite del sistema, usuarios, stakeholders, entorno operativo, suposiciones, dependencias y diagramas de dependencia y razón estratégica;
  \item \textbf{Sección 3:} interfaces, requisitos funcionales, requisitos no funcionales, requisitos legales, restricciones, historias de usuario, INVEST, Gherkin y mockups;
  \item \textbf{Sección 4:} modelado UML completo;
  \item \textbf{Sección 5:} priorización MoSCoW, Kano, WSJF y trazabilidad extendida;
  \item \textbf{Sección 6:} alcance, cobertura y verificación del MVP;
  \item \textbf{Sección 7:} componente empírico, protocolo experimental y registro OSF;
  \item \textbf{Sección 8:} preparación del conjunto de datos, manuscrito y análisis de revistas;
  \item \textbf{Sección 9:} conclusiones y trabajo futuro.
\end{itemize}

% =====================================================
% 2. DESCRIPCIÓN GENERAL
% =====================================================
\section{Descripción general}

\subsection{Perspectiva del producto}

RoutTech se concibe como un producto académico de software orientado a apoyar la coordinación de movilidad colaborativa entre integrantes adultos de la comunidad universitaria de la UTEQ. Su propósito no es operar como una empresa de transporte ni sustituir los servicios institucionales existentes, sino centralizar la publicación, consulta, reserva y seguimiento básico de viajes compartidos que actualmente pueden coordinarse mediante canales informales.

El producto tendrá una interfaz web adaptable a dispositivos móviles y organizará sus funciones alrededor de cuatro grupos principales: pasajeros universitarios, conductores universitarios, administradores del sistema y personal institucional autorizado para consultar información agregada. El sistema podrá relacionarse con servicios externos de autenticación, mapas y notificaciones; durante el MVP estas integraciones podrán ser simuladas cuando no exista autorización o disponibilidad técnica.

La frontera funcional del sistema se representa en la Figura~\ref{fig:contexto-routtech}. Dentro del límite de RoutTech se encuentran la gestión de identidad, perfiles, vehículos, rutas, reservas, cupos, reputación, incidencias, notificaciones, reportes y recomendaciones explicables. Permanecen fuera del sistema los pagos, la gestión de flotas institucionales, la atención de emergencias, el seguimiento público en tiempo real y la publicación de ubicaciones exactas o datos identificables.

\FiguraDocumento
  {imagenes/FIG-01_Diagrama_Contexto_RoutTech.jpg}
  {0.92}
  {Diagrama de contexto de RoutTech}
  {fig:contexto-routtech}

\subsection{Funciones generales del producto}

Las capacidades generales descritas en la Tabla~\ref{tab:funciones-generales} constituyen una organización preliminar del producto. No reemplazan los requisitos funcionales de la Sección 3; su finalidad es presentar el sistema de manera comprensible antes de formalizar cada requisito con sus ocho atributos y su evidencia específica.

\begin{table}[H]
\centering
\caption{Funciones generales previstas para RoutTech}
\label{tab:funciones-generales}
\small
\rowcolors{2}{DocLight}{white}
\begin{tabularx}{\textwidth}{L{3.2cm} Y L{3.0cm}}
\rowcolor{DocNavy}
\textcolor{white}{\textbf{Área funcional}} &
\textcolor{white}{\textbf{Descripción general}} &
\textcolor{white}{\textbf{Usuarios principales}}\\
Identidad y acceso & Registrar cuentas, validar la pertenencia institucional, iniciar y cerrar sesión y aplicar permisos según el rol. & Todos los usuarios; administrador.\\
Perfiles y vehículos & Gestionar datos de perfil, preferencias de viaje y la información mínima del vehículo utilizado para publicar rutas. & Pasajero; conductor.\\
Publicación de rutas & Registrar origen generalizado, destino, fecha, hora, puntos de encuentro, cupos y condiciones del viaje. & Conductor universitario.\\
Búsqueda y recomendación & Consultar rutas por filtros y mostrar coincidencias justificadas mediante factores comprensibles para el usuario. & Pasajero universitario.\\
Reservas y ciclo del viaje & Solicitar cupos, confirmar o rechazar solicitudes, cancelar reservas y controlar los estados del viaje. & Pasajero; conductor.\\
Comunicación y avisos & Informar confirmaciones, cambios, cancelaciones y recordatorios sin exponer información personal innecesaria. & Pasajero; conductor.\\
Confianza y seguridad & Gestionar calificaciones, reportes de incidencias, controles de privacidad y trazabilidad de acciones relevantes. & Todos los usuarios; administrador.\\
Administración y reportes & Supervisar cuentas, rutas, incidencias y estadísticas agregadas de uso, sin publicar trayectorias individuales. & Administrador; personal institucional autorizado.\\
\end{tabularx}
\end{table}

\subsection{Clases y características de usuarios}

Las clases de usuario de la Tabla~\ref{tab:clases-usuario} se definieron a partir del alcance de RoutTech y de los perfiles considerados en las entrevistas y el cuestionario. La clasificación se revisará durante los walkthroughs para confirmar responsabilidades, permisos y necesidades particulares.

\begin{table}[H]
\centering
\caption{Clases y características de usuarios de RoutTech}
\label{tab:clases-usuario}
\small
\rowcolors{2}{DocLight}{white}
\begin{tabularx}{\textwidth}{L{3.15cm} Y L{2.35cm} L{3.25cm}}
\rowcolor{DocNavy}
\textcolor{white}{\textbf{Clase de usuario}} &
\textcolor{white}{\textbf{Objetivo dentro del sistema}} &
\textcolor{white}{\textbf{Nivel técnico esperado}} &
\textcolor{white}{\textbf{Necesidades principales}}\\
Pasajero universitario & Encontrar una ruta compatible, solicitar un cupo y mantenerse informado sobre el estado del viaje. & Básico. & Búsqueda sencilla, información clara, privacidad, confianza y cancelación controlada.\\
Conductor universitario & Publicar rutas disponibles, administrar cupos y coordinar el encuentro con pasajeros verificados. & Básico. & Registro de ruta ágil, control de solicitudes, información suficiente y gestión de incidencias.\\
Administrador del sistema & Controlar usuarios, roles, incidencias, configuraciones y registros operativos. & Intermedio o técnico. & Trazabilidad, permisos, reportes, auditoría y herramientas de moderación.\\
Personal institucional autorizado & Consultar indicadores agregados para analizar tendencias generales de movilidad universitaria. & Básico o intermedio. & Reportes comprensibles, datos anonimizados y filtros por periodos o sectores generales.\\
\end{tabularx}
\end{table}

\subsection{Stakeholders y matriz poder--interés}

RoutTech involucra actores que utilizan directamente el producto y otros que influyen en su autorización, desarrollo, evaluación o sostenibilidad. La Tabla~\ref{tab:stakeholders} resume el análisis preliminar de poder e interés. Esta clasificación es una herramienta de gestión y deberá actualizarse cuando cambien las responsabilidades o aparezcan nuevos interesados.

\begin{table}[H]
\centering
\caption{Mapa preliminar de stakeholders de RoutTech}
\label{tab:stakeholders}
\small
\rowcolors{2}{DocLight}{white}
\begin{tabularx}{\textwidth}{L{3.5cm} C{1.55cm} C{1.55cm} Y}
\rowcolor{DocNavy}
\textcolor{white}{\textbf{Stakeholder}} &
\textcolor{white}{\textbf{Poder}} &
\textcolor{white}{\textbf{Interés}} &
\textcolor{white}{\textbf{Estrategia de relación}}\\
Estudiantes y pasajeros potenciales & Medio & Alto & Involucrar en la elicitación, priorización y validación de flujos de búsqueda y reserva.\\
Conductores universitarios potenciales & Medio & Alto & Validar publicación de rutas, cupos, puntos de encuentro, privacidad y condiciones de confianza.\\
Administrador del sistema & Alto & Alto & Gestionar de cerca; validar permisos, incidencias, reportes y trazabilidad.\\
Autoridades o unidades de la UTEQ & Alto & Medio & Mantener informadas; solicitar avales y presentar resultados agregados, no identificables.\\
Docente responsable & Alto & Alto & Mantener coordinación continua sobre requisitos, ética, evidencia y entregables.\\
Equipo CN & Medio & Alto & Coordinar análisis, documentación, modelado, MVP, trazabilidad y control de cambios.\\
Proveedores de servicios externos & Bajo & Medio & Supervisar disponibilidad, términos de uso, privacidad y alternativas de simulación.\\
\end{tabularx}
\end{table}

La ubicación visual de estos grupos se presenta en la Figura~\ref{fig:poder-interes}. En la versión definitiva, la figura deberá coincidir con los valores consignados en la Tabla~\ref{tab:stakeholders}.

\FiguraDocumento
  {imagenes/FIG-02_Matriz_Poder_Interes_RoutTech.jpg}
  {0.82}
  {Matriz poder--interés de los stakeholders de RoutTech}
  {fig:poder-interes}

\subsection{Entorno operativo}

El entorno descrito en la Tabla~\ref{tab:entorno-operativo} corresponde a la versión académica y al MVP. Las tecnologías concretas de implementación se documentarán cuando se consolide la arquitectura y se verifique su compatibilidad con los requisitos priorizados.

\begin{table}[H]
\centering
\caption{Entorno operativo previsto para RoutTech}
\label{tab:entorno-operativo}
\small
\rowcolors{2}{DocLight}{white}
\begin{tabularx}{\textwidth}{L{3.0cm} Y L{4.0cm}}
\rowcolor{DocNavy}
\textcolor{white}{\textbf{Elemento}} &
\textcolor{white}{\textbf{Entorno previsto}} &
\textcolor{white}{\textbf{Condición de operación}}\\
Cliente web & Navegadores modernos en equipos de escritorio, portátiles y teléfonos. & Diseño adaptable y compatibilidad verificada en al menos Chrome, Edge y Firefox.\\
Dispositivo móvil & Teléfono con navegador actualizado y acceso a Internet. & La versión académica no requiere seguimiento GPS permanente.\\
Servidor de aplicación & Entorno local o servicio de prueba con API para reglas de negocio. & Acceso controlado y configuración reproducible mediante instrucciones del repositorio.\\
Persistencia & Base de datos relacional para cuentas, rutas, reservas, estados e incidencias sintéticas. & Integridad referencial, control por roles y respaldo de datos de prueba.\\
Servicios externos & Autenticación, mapas y notificaciones, cuando exista autorización y disponibilidad. & En el MVP podrán sustituirse por adaptadores simulados o datos sintéticos.\\
Conectividad & Acceso a Internet para sincronización de operaciones. & Los errores de red deberán informarse sin provocar pérdida o duplicación de transacciones.\\
\end{tabularx}
\end{table}

\subsection{Suposiciones y dependencias}

Las suposiciones de la Tabla~\ref{tab:suposiciones} describen condiciones consideradas razonables para planificar el producto. Las dependencias de la Tabla~\ref{tab:dependencias} corresponden a elementos externos cuya ausencia puede afectar el alcance, el MVP o la validación.

\begin{table}[H]
\centering
\caption{Suposiciones del proyecto RoutTech}
\label{tab:suposiciones}
\small
\rowcolors{2}{DocLight}{white}
\begin{tabularx}{\textwidth}{C{1.35cm} Y Y}
\rowcolor{DocNavy}
\textcolor{white}{\textbf{ID}} &
\textcolor{white}{\textbf{Suposición}} &
\textcolor{white}{\textbf{Justificación}}\\
SUP-01 & Las personas usuarias pertenecen a la comunidad universitaria y son mayores de edad en esta fase. & La delimitación reduce riesgos éticos y corresponde al protocolo aprobado para el trabajo de campo.\\
SUP-02 & Los usuarios disponen de un dispositivo con navegador y conexión a Internet. & Las funciones principales se ejecutarán mediante una interfaz web adaptable.\\
SUP-03 & Los conductores interesados proporcionarán información veraz y mínima sobre su vehículo y disponibilidad. & La publicación y la reserva dependen de datos consistentes sobre cupos y horarios.\\
SUP-04 & Los pasajeros aceptarán usar puntos de encuentro generales en lugar de domicilios exactos. & La generalización de ubicación protege la privacidad y disminuye exposición innecesaria.\\
SUP-05 & Las personas participantes autorizarán el uso académico de información anonimizada. & La evidencia solo puede utilizarse conforme al consentimiento informado y al plan de gestión de datos.\\
SUP-06 & El MVP utilizará exclusivamente datos ficticios o sintéticos. & La demostración no debe exponer credenciales, placas, rutas ni datos personales reales.\\
\end{tabularx}
\end{table}

\begin{table}[H]
\centering
\caption{Dependencias externas de RoutTech}
\label{tab:dependencias}
\small
\rowcolors{2}{DocLight}{white}
\begin{tabularx}{\textwidth}{C{1.35cm} Y Y}
\rowcolor{DocNavy}
\textcolor{white}{\textbf{ID}} &
\textcolor{white}{\textbf{Dependencia}} &
\textcolor{white}{\textbf{Efecto sobre el proyecto}}\\
DEP-01 & Vigencia de la documentación ética y de los consentimientos aplicables. & Sin autorización válida no se podrán incorporar evidencias ni realizar nuevas sesiones con participantes.\\
DEP-02 & Disponibilidad de un mecanismo de verificación institucional o de su simulación controlada. & Afecta el registro y la diferenciación de usuarios autorizados.\\
DEP-03 & Servicio de mapas o conjunto sintético equivalente para representar rutas y puntos generales. & Afecta la visualización y el cálculo de coincidencias entre trayectos.\\
DEP-04 & Servicio de notificaciones o componente simulado durante el MVP. & Afecta la comunicación de cambios, reservas, cancelaciones y recordatorios.\\
DEP-05 & Disponibilidad de participantes para walkthroughs técnicos y no técnicos. & Afecta la validación de requisitos, mockups y explicaciones de recomendaciones.\\
DEP-06 & Repositorio GitHub y herramientas de compilación y despliegue disponibles. & Afecta la reproducibilidad del documento, el MVP y los artefactos de evaluación.\\
\end{tabularx}
\end{table}

\subsection{Diagramas de dependencia y razón estratégica}

El Diagrama de Dependencia Estratégica representará las relaciones de dependencia entre los actores de RoutTech. Mostrará, entre otras relaciones, que el pasajero depende del sistema para encontrar y reservar rutas, que el conductor depende de la plataforma para publicar viajes y administrar cupos, y que la administración depende de información agregada para supervisar la operación.

La Figura~\ref{fig:dependencia-estrategica} presenta el Diagrama de Dependencia Estratégica. La Figura~\ref{fig:razon-estrategica} amplía el análisis mediante el Diagrama de Razón Estratégica, en el que se descomponen objetivos como coordinar un viaje seguro, proteger los datos personales, comprender una recomendación y mantener la disponibilidad de cupos.

\FiguraDocumento
  {imagenes/FIG-03_Diagrama_Dependencia_Estrategica_RoutTech.jpg}
  {0.94}
  {Diagrama de Dependencia Estratégica de RoutTech}
  {fig:dependencia-estrategica}

\FiguraDocumento
  {imagenes/FIG-04_Diagrama_Razon_Estrategica_RoutTech.jpg}
  {0.94}
  {Diagrama de Razón Estratégica de RoutTech}
  {fig:razon-estrategica}


% =====================================================
% 3. REQUISITOS ESPECIFICOS
% =====================================================
\section{Requisitos específicos}

\subsection{Convención de especificación}

El catálogo funcional se redactó con lenguaje verificable y trazabilidad explícita, de acuerdo con la estructura recomendada por ISO/IEC/IEEE 29148:2018 \cite{iso29148}. Cada requisito incluye ocho atributos de especificación: descripción, actor principal, entradas, salidas, precondición, postcondición, prioridad y método de verificación. La evidencia se registra en un campo independiente mediante códigos \texttt{EV-XX}.

Los códigos \texttt{EV-ENT-001} a \texttt{EV-ENT-016} corresponden a las entrevistas semiestructuradas y \texttt{EV-CUE-001} al cuestionario con 43 respuestas. Para las obligaciones que provienen de la guía o de la documentación ética se incorporan los códigos documentales indicados en la Tabla~\ref{tab:evidencia-documental-rf}. La guía de la Entrega 3 exige un mínimo de 25 requisitos funcionales con atributos completos y evidencia verificable \cite{uteq2026}.

\begin{table}[H]
\centering
\caption{Evidencias documentales utilizadas en el catálogo funcional}
\label{tab:evidencia-documental-rf}
\small
\rowcolors{2}{DocLight}{white}
\begin{tabularx}{\textwidth}{L{2.8cm} Y Y}
\rowcolor{DocNavy}
\textcolor{white}{\textbf{Código}} &
\textcolor{white}{\textbf{Documento de origen}} &
\textcolor{white}{\textbf{Uso dentro del catálogo}}\\
\texttt{EV-DOC-001} & Guía y Rúbrica de la Entrega 3 (2A). & Obligatoriedad de tratar la explicabilidad en todo componente basado en IA.\\
\texttt{EV-DOC-002} & Política de geolocalización y análisis de riesgos de RoutTech. & Consentimiento, ubicación generalizada, revocación y explicación no obligatoria de recomendaciones.\\
\texttt{EV-DOC-003} & Protocolo de investigación y Plan de Gestión de Datos. & Uso de resultados agregados o disociados y generación de productos administrativos sin datos identificables.\\
\end{tabularx}
\end{table}

\subsection{Resumen del catálogo funcional}

La Tabla~\ref{tab:resumen-rf} presenta la identificación y prioridad de los 25 requisitos funcionales. La especificación detallada se desarrolla posteriormente en fichas individuales para evitar tablas horizontales de difícil lectura.

\begin{longtable}{C{1.45cm} L{10.5cm} C{2.2cm}}
\caption{Resumen de requisitos funcionales de RoutTech}\label{tab:resumen-rf}\\
\rowcolor{DocNavy}
\textcolor{white}{\textbf{ID}} & \textcolor{white}{\textbf{Nombre}} & \textcolor{white}{\textbf{MoSCoW}}\\
\endfirsthead
\rowcolor{DocNavy}
\textcolor{white}{\textbf{ID}} & \textcolor{white}{\textbf{Nombre}} & \textcolor{white}{\textbf{MoSCoW}}\\
\endhead
RF-01 & Autenticación de usuarios institucionales & Must\\
RF-02 & Gestión de roles y permisos & Must\\
RF-03 & Recomendación de rutas compatibles & Must\\
RF-04 & Publicación de rutas & Must\\
RF-05 & Reserva de cupos & Must\\
RF-06 & Cancelación de reserva o ruta & Should\\
RF-07 & Gestión del perfil de usuario & Should\\
RF-08 & Consulta del historial de viajes & Should\\
RF-09 & Calificación de viajes & Should\\
RF-10 & Análisis de patrones de movilidad & Could\\
RF-11 & Generación de reportes administrativos & Could\\
RF-12 & Notificaciones de eventos de viaje & Must\\
RF-13 & Gestión de incidencias & Should\\
RF-14 & Búsqueda avanzada de rutas & Should\\
RF-15 & Registro de vehículo & Must\\
RF-16 & Verificación de identidad institucional & Must\\
RF-17 & Compartir ubicación general del recorrido & Should\\
RF-18 & Confirmación de llegada segura & Should\\
RF-19 & Gestión de puntos de encuentro & Should\\
RF-20 & Comunicación mediante chat & Should\\
RF-21 & Recuperación ante cancelaciones & Should\\
RF-22 & Explicación de recomendaciones & Could\\
RF-23 & Configuración de preferencias de viaje & Should\\
RF-24 & Consulta del historial de reputación & Must\\
RF-25 & Panel estadístico institucional & Could\\
\end{longtable}

\subsection{Catálogo detallado de requisitos funcionales}

\begin{RFBox}{RF-01: Autenticación de usuarios institucionales}
\small
\begin{tabularx}{\linewidth}{@{}L{3.15cm}Y@{}}
\CampoRF{Descripción}{El sistema deberá autenticar a cada usuario mediante correo institucional y contraseña válidos antes de permitir el acceso a las funciones correspondientes a su perfil.}
\CampoRF{Actor principal}{Usuario registrado.}
\CampoRF{Entradas}{Correo institucional y contraseña.}
\CampoRF{Salidas}{Acceso autorizado o mensaje de autenticación rechazada.}
\CampoRF{Precondición}{La cuenta existe, se encuentra activa y no está bloqueada.}
\CampoRF{Postcondición}{La sesión queda iniciada y asociada al rol autorizado, o se registra el intento fallido.}
\CampoRF{Prioridad}{Must have.}
\CampoRF{Método de verificación}{Pruebas funcionales con credenciales válidas, inválidas, inactivas y bloqueadas.}
\CampoRF{Evidencia}{\texttt{EV-CUE-001; EV-ENT-001; EV-ENT-005.}}
\end{tabularx}
\end{RFBox}

\begin{RFBox}{RF-02: Gestión de roles y permisos}
\small
\begin{tabularx}{\linewidth}{@{}L{3.15cm}Y@{}}
\CampoRF{Descripción}{El sistema deberá permitir al administrador asignar y actualizar los roles de pasajero, conductor y administrador, restringiendo las funciones disponibles según los permisos definidos.}
\CampoRF{Actor principal}{Administrador del sistema.}
\CampoRF{Entradas}{Usuario seleccionado, rol y permisos autorizados.}
\CampoRF{Salidas}{Rol asignado y matriz de permisos actualizada.}
\CampoRF{Precondición}{El administrador inició sesión y el usuario objetivo existe.}
\CampoRF{Postcondición}{Los nuevos permisos quedan vigentes y el cambio se registra para auditoría.}
\CampoRF{Prioridad}{Must have.}
\CampoRF{Método de verificación}{Pruebas de autorización intentando acceder a funciones permitidas y restringidas para cada rol.}
\CampoRF{Evidencia}{\texttt{EV-ENT-001; EV-ENT-006.}}
\end{tabularx}
\end{RFBox}

\begin{RFBox}{RF-03: Recomendación de rutas compatibles}
\small
\begin{tabularx}{\linewidth}{@{}L{3.15cm}Y@{}}
\CampoRF{Descripción}{El sistema deberá generar una lista ordenada de rutas compatibles considerando origen y destino generalizados, fecha, horario, disponibilidad de cupos y preferencias configuradas por el usuario.}
\CampoRF{Actor principal}{Pasajero universitario y sistema de recomendación.}
\CampoRF{Entradas}{Origen, destino, fecha, intervalo horario y preferencias de viaje.}
\CampoRF{Salidas}{Lista priorizada de rutas compatibles.}
\CampoRF{Precondición}{El usuario está autenticado y existen rutas activas con cupos disponibles.}
\CampoRF{Postcondición}{Las recomendaciones se muestran sin crear automáticamente una reserva.}
\CampoRF{Prioridad}{Must have.}
\CampoRF{Método de verificación}{Pruebas con combinaciones controladas de rutas y comparación entre los criterios registrados y el orden obtenido.}
\CampoRF{Evidencia}{\texttt{EV-CUE-001; EV-ENT-003; EV-ENT-007; EV-ENT-010.}}
\end{tabularx}
\end{RFBox}

\begin{RFBox}{RF-04: Publicación de rutas}
\small
\begin{tabularx}{\linewidth}{@{}L{3.15cm}Y@{}}
\CampoRF{Descripción}{El sistema deberá permitir al conductor publicar una ruta indicando origen y destino generalizados, fecha, hora de salida, cupos disponibles, punto de encuentro y observaciones permitidas.}
\CampoRF{Actor principal}{Conductor universitario.}
\CampoRF{Entradas}{Datos del trayecto, horario, cupos y punto de encuentro.}
\CampoRF{Salidas}{Ruta publicada con identificador y estado activo.}
\CampoRF{Precondición}{El conductor está autenticado, verificado y tiene un vehículo habilitado.}
\CampoRF{Postcondición}{La ruta queda disponible para búsqueda y reserva mientras mantenga cupos y vigencia.}
\CampoRF{Prioridad}{Must have.}
\CampoRF{Método de verificación}{Pruebas de creación con datos válidos, campos obligatorios vacíos, fechas pasadas y cupos inválidos.}
\CampoRF{Evidencia}{\texttt{EV-ENT-002; EV-ENT-012; EV-ENT-015.}}
\end{tabularx}
\end{RFBox}

\begin{RFBox}{RF-05: Reserva de cupos}
\small
\begin{tabularx}{\linewidth}{@{}L{3.15cm}Y@{}}
\CampoRF{Descripción}{El sistema deberá permitir al pasajero solicitar y confirmar un cupo en una ruta disponible, evitando superar la capacidad declarada por el conductor.}
\CampoRF{Actor principal}{Pasajero universitario.}
\CampoRF{Entradas}{Ruta seleccionada y cantidad de cupos solicitados.}
\CampoRF{Salidas}{Reserva confirmada o mensaje de indisponibilidad.}
\CampoRF{Precondición}{El pasajero está autenticado y la ruta permanece activa con cupos suficientes.}
\CampoRF{Postcondición}{La reserva queda registrada y la disponibilidad de la ruta se actualiza de forma atómica.}
\CampoRF{Prioridad}{Must have.}
\CampoRF{Método de verificación}{Pruebas de reserva concurrente, último cupo, ruta llena, ruta cancelada y solicitud duplicada.}
\CampoRF{Evidencia}{\texttt{EV-CUE-001; EV-ENT-004; EV-ENT-011.}}
\end{tabularx}
\end{RFBox}

\begin{RFBox}{RF-06: Cancelación de reserva o ruta}
\small
\begin{tabularx}{\linewidth}{@{}L{3.15cm}Y@{}}
\CampoRF{Descripción}{El sistema deberá permitir cancelar una reserva o una ruta activa, solicitando el motivo y notificando a las personas afectadas por el cambio.}
\CampoRF{Actor principal}{Pasajero universitario o conductor universitario.}
\CampoRF{Entradas}{Identificador de reserva o ruta y motivo de cancelación.}
\CampoRF{Salidas}{Confirmación de cancelación y notificaciones asociadas.}
\CampoRF{Precondición}{Existe una reserva o ruta activa perteneciente al usuario que solicita la cancelación.}
\CampoRF{Postcondición}{El elemento queda en estado cancelado, se liberan los cupos correspondientes y se registra el evento.}
\CampoRF{Prioridad}{Should have.}
\CampoRF{Método de verificación}{Pruebas de cancelación por pasajero y conductor, incluyendo liberación de cupos y envío de avisos.}
\CampoRF{Evidencia}{\texttt{EV-ENT-008; EV-ENT-016.}}
\end{tabularx}
\end{RFBox}

\begin{RFBox}{RF-07: Gestión del perfil de usuario}
\small
\begin{tabularx}{\linewidth}{@{}L{3.15cm}Y@{}}
\CampoRF{Descripción}{El sistema deberá permitir al usuario consultar y actualizar los datos editables de su perfil, manteniendo protegidos los campos institucionales que no puedan modificarse directamente.}
\CampoRF{Actor principal}{Usuario autenticado.}
\CampoRF{Entradas}{Datos personales permitidos, preferencias y configuración de contacto.}
\CampoRF{Salidas}{Perfil actualizado o detalle de validaciones incumplidas.}
\CampoRF{Precondición}{El usuario inició sesión.}
\CampoRF{Postcondición}{Los cambios válidos quedan almacenados y disponibles para las funciones relacionadas.}
\CampoRF{Prioridad}{Should have.}
\CampoRF{Método de verificación}{Pruebas de actualización válida, datos incompletos, formato incorrecto y modificación de campos protegidos.}
\CampoRF{Evidencia}{\texttt{EV-CUE-001; EV-ENT-008.}}
\end{tabularx}
\end{RFBox}

\begin{RFBox}{RF-08: Consulta del historial de viajes}
\small
\begin{tabularx}{\linewidth}{@{}L{3.15cm}Y@{}}
\CampoRF{Descripción}{El sistema deberá almacenar y mostrar al usuario el historial de rutas publicadas, reservas y viajes finalizados que le correspondan.}
\CampoRF{Actor principal}{Pasajero universitario o conductor universitario.}
\CampoRF{Entradas}{Solicitud de consulta y filtros opcionales por fecha o estado.}
\CampoRF{Salidas}{Listado cronológico de viajes relacionados con el usuario.}
\CampoRF{Precondición}{El usuario está autenticado y posee registros asociados.}
\CampoRF{Postcondición}{El historial se presenta sin modificar la información almacenada.}
\CampoRF{Prioridad}{Should have.}
\CampoRF{Método de verificación}{Pruebas de consulta con diferentes estados, periodos, roles y usuarios sin historial.}
\CampoRF{Evidencia}{\texttt{EV-ENT-005; EV-ENT-009.}}
\end{tabularx}
\end{RFBox}

\begin{RFBox}{RF-09: Calificación de viajes}
\small
\begin{tabularx}{\linewidth}{@{}L{3.15cm}Y@{}}
\CampoRF{Descripción}{El sistema deberá permitir que pasajeros y conductores califiquen la experiencia una vez finalizado el viaje, mediante una puntuación y un comentario opcional moderado.}
\CampoRF{Actor principal}{Pasajero universitario y conductor universitario.}
\CampoRF{Entradas}{Puntuación, comentario opcional e identificador del viaje.}
\CampoRF{Salidas}{Calificación registrada y reputación recalculada.}
\CampoRF{Precondición}{El viaje finalizó y el usuario participó en él sin haberlo calificado previamente.}
\CampoRF{Postcondición}{La evaluación queda asociada al viaje y al usuario evaluado.}
\CampoRF{Prioridad}{Should have.}
\CampoRF{Método de verificación}{Pruebas de calificación válida, duplicada, fuera de rango y vinculada a un viaje ajeno.}
\CampoRF{Evidencia}{\texttt{EV-ENT-003; EV-ENT-010; EV-ENT-014.}}
\end{tabularx}
\end{RFBox}

\begin{RFBox}{RF-10: Análisis de patrones de movilidad}
\small
\begin{tabularx}{\linewidth}{@{}L{3.15cm}Y@{}}
\CampoRF{Descripción}{El sistema deberá procesar datos históricos anonimizados para identificar tendencias agregadas de horarios, zonas generales, demanda y disponibilidad de rutas.}
\CampoRF{Actor principal}{Sistema y administrador.}
\CampoRF{Entradas}{Registros históricos anonimizados y periodo de análisis.}
\CampoRF{Salidas}{Indicadores y patrones agregados de movilidad.}
\CampoRF{Precondición}{Existe un volumen mínimo de registros anonimizados y el administrador posee autorización.}
\CampoRF{Postcondición}{Los resultados se generan sin exponer recorridos individuales ni identificadores personales.}
\CampoRF{Prioridad}{Could have.}
\CampoRF{Método de verificación}{Pruebas con conjuntos sintéticos de datos y revisión de que los resultados coincidan con valores calculados de referencia.}
\CampoRF{Evidencia}{\texttt{EV-CUE-001; EV-ENT-011.}}
\end{tabularx}
\end{RFBox}

\begin{RFBox}{RF-11: Generación de reportes administrativos}
\small
\begin{tabularx}{\linewidth}{@{}L{3.15cm}Y@{}}
\CampoRF{Descripción}{El sistema deberá generar reportes agregados sobre usuarios activos, rutas, reservas, cancelaciones, incidencias y calificaciones, con filtros por periodo y zona general.}
\CampoRF{Actor principal}{Administrador del sistema.}
\CampoRF{Entradas}{Tipo de reporte, periodo y filtros autorizados.}
\CampoRF{Salidas}{Reporte visual y archivo exportable con información agregada.}
\CampoRF{Precondición}{El administrador está autenticado y existen datos procesables.}
\CampoRF{Postcondición}{El reporte queda disponible sin incluir datos personales identificables.}
\CampoRF{Prioridad}{Could have.}
\CampoRF{Método de verificación}{Comparación del reporte con un conjunto sintético conocido y revisión de anonimización de las salidas.}
\CampoRF{Evidencia}{\texttt{EV-DOC-003.}}
\end{tabularx}
\end{RFBox}

\begin{RFBox}{RF-12: Notificaciones de eventos de viaje}
\small
\begin{tabularx}{\linewidth}{@{}L{3.15cm}Y@{}}
\CampoRF{Descripción}{El sistema deberá emitir notificaciones por eventos relevantes, incluyendo aceptación o rechazo de solicitudes, confirmaciones, cambios de horario, cancelaciones y recordatorios.}
\CampoRF{Actor principal}{Sistema.}
\CampoRF{Entradas}{Evento de negocio y destinatarios relacionados.}
\CampoRF{Salidas}{Notificación enviada o intento de entrega registrado.}
\CampoRF{Precondición}{Existe un evento válido asociado con una ruta, reserva, viaje o incidencia.}
\CampoRF{Postcondición}{El aviso queda registrado con estado de entrega y marca temporal.}
\CampoRF{Prioridad}{Must have.}
\CampoRF{Método de verificación}{Pruebas por cada tipo de evento, destinatario y estado de entrega, evitando avisos duplicados.}
\CampoRF{Evidencia}{\texttt{EV-ENT-006; EV-ENT-011; EV-ENT-012.}}
\end{tabularx}
\end{RFBox}

\begin{RFBox}{RF-13: Gestión de incidencias}
\small
\begin{tabularx}{\linewidth}{@{}L{3.15cm}Y@{}}
\CampoRF{Descripción}{El sistema deberá permitir a los usuarios registrar incidencias relacionadas con el viaje y al administrador clasificarlas, darles seguimiento y cerrar su atención.}
\CampoRF{Actor principal}{Usuario y administrador del sistema.}
\CampoRF{Entradas}{Categoría, descripción, viaje relacionado y evidencia permitida.}
\CampoRF{Salidas}{Número de incidencia, estado y respuesta de seguimiento.}
\CampoRF{Precondición}{El usuario está autenticado y la incidencia corresponde a una interacción registrada.}
\CampoRF{Postcondición}{La incidencia queda almacenada con estado, responsable y trazabilidad de cambios.}
\CampoRF{Prioridad}{Should have.}
\CampoRF{Método de verificación}{Pruebas de creación, clasificación, cambio de estado, cierre y restricción de acceso por rol.}
\CampoRF{Evidencia}{\texttt{EV-ENT-009; EV-ENT-013.}}
\end{tabularx}
\end{RFBox}

\begin{RFBox}{RF-14: Búsqueda avanzada de rutas}
\small
\begin{tabularx}{\linewidth}{@{}L{3.15cm}Y@{}}
\CampoRF{Descripción}{El sistema deberá permitir buscar y filtrar rutas por origen y destino generalizados, fecha, intervalo horario, cupos, conductor y preferencias compatibles.}
\CampoRF{Actor principal}{Pasajero universitario.}
\CampoRF{Entradas}{Criterios de búsqueda y filtros seleccionados.}
\CampoRF{Salidas}{Resultados que cumplen los criterios indicados.}
\CampoRF{Precondición}{El usuario está autenticado y existen rutas registradas.}
\CampoRF{Postcondición}{Los resultados se muestran ordenados sin alterar las rutas publicadas.}
\CampoRF{Prioridad}{Should have.}
\CampoRF{Método de verificación}{Pruebas de filtrado individual, combinado, sin resultados y con datos límite.}
\CampoRF{Evidencia}{\texttt{EV-CUE-001; EV-ENT-014.}}
\end{tabularx}
\end{RFBox}

\begin{RFBox}{RF-15: Registro de vehículo}
\small
\begin{tabularx}{\linewidth}{@{}L{3.15cm}Y@{}}
\CampoRF{Descripción}{El sistema deberá permitir al conductor registrar y mantener la información mínima del vehículo utilizado en las rutas publicadas.}
\CampoRF{Actor principal}{Conductor universitario.}
\CampoRF{Entradas}{Tipo, marca, modelo, color, capacidad y datos de identificación autorizados.}
\CampoRF{Salidas}{Vehículo asociado al perfil del conductor.}
\CampoRF{Precondición}{El conductor está autenticado y verificado.}
\CampoRF{Postcondición}{El vehículo queda registrado para su selección al publicar una ruta.}
\CampoRF{Prioridad}{Must have.}
\CampoRF{Método de verificación}{Pruebas de alta, actualización, capacidad inválida, duplicados y eliminación de vehículo asociado a rutas activas.}
\CampoRF{Evidencia}{\texttt{EV-ENT-002; EV-ENT-015.}}
\end{tabularx}
\end{RFBox}

\begin{RFBox}{RF-16: Verificación de identidad institucional}
\small
\begin{tabularx}{\linewidth}{@{}L{3.15cm}Y@{}}
\CampoRF{Descripción}{El sistema deberá comprobar que pasajeros y conductores pertenecen a la comunidad UTEQ antes de habilitar las funciones de movilidad colaborativa.}
\CampoRF{Actor principal}{Sistema y servicio de autenticación institucional.}
\CampoRF{Entradas}{Correo institucional, token o respuesta del servicio de verificación.}
\CampoRF{Salidas}{Estado de usuario verificado, rechazado o pendiente.}
\CampoRF{Precondición}{El usuario completó el registro y aceptó las condiciones de uso aplicables.}
\CampoRF{Postcondición}{La cuenta queda habilitada según su estado de verificación o permanece restringida.}
\CampoRF{Prioridad}{Must have.}
\CampoRF{Método de verificación}{Pruebas con identidades válidas, inexistentes, expiradas y respuestas no disponibles del servicio externo.}
\CampoRF{Evidencia}{\texttt{EV-CUE-001; EV-ENT-001; EV-ENT-006; EV-ENT-009.}}
\end{tabularx}
\end{RFBox}

\begin{RFBox}{RF-17: Compartir ubicación general del recorrido}
\small
\begin{tabularx}{\linewidth}{@{}L{3.15cm}Y@{}}
\CampoRF{Descripción}{El sistema deberá permitir compartir durante un viaje una ubicación general o aproximada entre conductor y pasajero únicamente cuando exista autorización expresa y revocable de las partes.}
\CampoRF{Actor principal}{Conductor universitario y pasajero universitario.}
\CampoRF{Entradas}{Consentimiento, estado del viaje y ubicación generalizada.}
\CampoRF{Salidas}{Ubicación aproximada visible durante el periodo autorizado.}
\CampoRF{Precondición}{Existe un viaje activo y ambas partes autorizaron el intercambio.}
\CampoRF{Postcondición}{La ubicación deja de compartirse al finalizar el viaje, revocar el permiso o vencer el periodo autorizado.}
\CampoRF{Prioridad}{Should have.}
\CampoRF{Método de verificación}{Pruebas de activación, revocación, finalización automática y verificación de que no se muestran domicilios ni coordenadas exactas.}
\CampoRF{Evidencia}{\texttt{EV-ENT-003; EV-ENT-010; EV-ENT-014; EV-DOC-002.}}
\end{tabularx}
\end{RFBox}

\begin{RFBox}{RF-18: Confirmación de llegada segura}
\small
\begin{tabularx}{\linewidth}{@{}L{3.15cm}Y@{}}
\CampoRF{Descripción}{El sistema deberá solicitar la confirmación de llegada al destino al finalizar el recorrido y registrar el cierre del viaje para los participantes.}
\CampoRF{Actor principal}{Pasajero universitario, conductor universitario y sistema.}
\CampoRF{Entradas}{Estado de viaje finalizado y confirmación de llegada.}
\CampoRF{Salidas}{Viaje cerrado y confirmación registrada.}
\CampoRF{Precondición}{El viaje se encuentra en curso o pendiente de cierre.}
\CampoRF{Postcondición}{El viaje pasa a estado finalizado y se habilitan las calificaciones correspondientes.}
\CampoRF{Prioridad}{Should have.}
\CampoRF{Método de verificación}{Pruebas de confirmación por ambos roles, cierre pendiente, ausencia de respuesta y cierre duplicado.}
\CampoRF{Evidencia}{\texttt{EV-ENT-005; EV-ENT-011.}}
\end{tabularx}
\end{RFBox}

\begin{RFBox}{RF-19: Gestión de puntos de encuentro}
\small
\begin{tabularx}{\linewidth}{@{}L{3.15cm}Y@{}}
\CampoRF{Descripción}{El sistema deberá permitir definir, consultar y seleccionar puntos de encuentro generales para facilitar el inicio del viaje sin revelar domicilios exactos.}
\CampoRF{Actor principal}{Conductor universitario y pasajero universitario.}
\CampoRF{Entradas}{Zona o punto general seleccionado para la ruta o reserva.}
\CampoRF{Salidas}{Punto de encuentro asociado y visible para los participantes autorizados.}
\CampoRF{Precondición}{Existe una ruta activa y el punto cumple las restricciones de ubicación establecidas.}
\CampoRF{Postcondición}{El punto queda registrado y puede ser actualizado antes del inicio del viaje.}
\CampoRF{Prioridad}{Should have.}
\CampoRF{Método de verificación}{Pruebas de selección, actualización, acceso por participantes y rechazo de ubicaciones no permitidas.}
\CampoRF{Evidencia}{\texttt{EV-CUE-001; EV-ENT-007; EV-ENT-012.}}
\end{tabularx}
\end{RFBox}

\begin{RFBox}{RF-20: Comunicación mediante chat}
\small
\begin{tabularx}{\linewidth}{@{}L{3.15cm}Y@{}}
\CampoRF{Descripción}{El sistema deberá habilitar un canal de mensajería entre conductor y pasajero después de aceptar la reserva, limitado a la coordinación del viaje.}
\CampoRF{Actor principal}{Conductor universitario y pasajero universitario.}
\CampoRF{Entradas}{Mensaje de texto y reserva asociada.}
\CampoRF{Salidas}{Mensaje entregado y conversación actualizada.}
\CampoRF{Precondición}{La reserva fue aceptada y ambos usuarios participan en el mismo viaje.}
\CampoRF{Postcondición}{El mensaje queda disponible en la conversación y se registra su estado de entrega.}
\CampoRF{Prioridad}{Should have.}
\CampoRF{Método de verificación}{Pruebas de envío, recepción, acceso por terceros, conversación cerrada y contenido fuera de los límites permitidos.}
\CampoRF{Evidencia}{\texttt{EV-ENT-002; EV-ENT-008; EV-ENT-015.}}
\end{tabularx}
\end{RFBox}

\begin{RFBox}{RF-21: Recuperación ante cancelaciones}
\small
\begin{tabularx}{\linewidth}{@{}L{3.15cm}Y@{}}
\CampoRF{Descripción}{El sistema deberá buscar y sugerir alternativas compatibles cuando una ruta o reserva sea cancelada antes del viaje.}
\CampoRF{Actor principal}{Sistema y pasajero universitario.}
\CampoRF{Entradas}{Evento de cancelación y parámetros de la reserva afectada.}
\CampoRF{Salidas}{Lista de rutas alternativas o aviso de que no existen opciones.}
\CampoRF{Precondición}{La cancelación fue confirmada y conserva información suficiente para una nueva búsqueda.}
\CampoRF{Postcondición}{Las alternativas se muestran sin reservar automáticamente un cupo.}
\CampoRF{Prioridad}{Should have.}
\CampoRF{Método de verificación}{Pruebas con alternativas disponibles, ausencia de alternativas, múltiples coincidencias y rutas sin cupos.}
\CampoRF{Evidencia}{\texttt{EV-ENT-008; EV-ENT-016.}}
\end{tabularx}
\end{RFBox}

\begin{RFBox}{RF-22: Explicación de recomendaciones}
\small
\begin{tabularx}{\linewidth}{@{}L{3.15cm}Y@{}}
\CampoRF{Descripción}{El sistema deberá mostrar una explicación comprensible de los principales factores utilizados para recomendar una ruta, permitiendo al usuario revisar por qué fue presentada.}
\CampoRF{Actor principal}{Sistema de recomendación y usuario.}
\CampoRF{Entradas}{Recomendación seleccionada y factores que influyeron en su cálculo.}
\CampoRF{Salidas}{Explicación de la recomendación y advertencia de que puede ser rechazada.}
\CampoRF{Precondición}{El sistema generó al menos una recomendación de ruta.}
\CampoRF{Postcondición}{La explicación queda visible sin convertir la recomendación en una decisión obligatoria.}
\CampoRF{Prioridad}{Could have.}
\CampoRF{Método de verificación}{Walkthrough con usuarios técnicos y no técnicos, comprobando la correspondencia entre factores de entrada y explicación mostrada.}
\CampoRF{Evidencia}{\texttt{EV-DOC-001; EV-DOC-002; posteriormente EV-VAL-XXX.}}
\end{tabularx}
\end{RFBox}

\begin{RFBox}{RF-23: Configuración de preferencias de viaje}
\small
\begin{tabularx}{\linewidth}{@{}L{3.15cm}Y@{}}
\CampoRF{Descripción}{El sistema deberá permitir al usuario configurar horarios habituales, zonas generales, tolerancia de tiempo y otras preferencias admitidas para personalizar búsquedas y recomendaciones.}
\CampoRF{Actor principal}{Usuario autenticado.}
\CampoRF{Entradas}{Preferencias de horario, zonas, frecuencia y características permitidas del viaje.}
\CampoRF{Salidas}{Preferencias almacenadas en el perfil.}
\CampoRF{Precondición}{El usuario inició sesión.}
\CampoRF{Postcondición}{Las preferencias quedan disponibles para búsquedas y recomendaciones posteriores.}
\CampoRF{Prioridad}{Should have.}
\CampoRF{Método de verificación}{Pruebas de registro, actualización, eliminación y aplicación de preferencias en resultados controlados.}
\CampoRF{Evidencia}{\texttt{EV-CUE-001; EV-ENT-004; EV-ENT-010.}}
\end{tabularx}
\end{RFBox}

\begin{RFBox}{RF-24: Consulta del historial de reputación}
\small
\begin{tabularx}{\linewidth}{@{}L{3.15cm}Y@{}}
\CampoRF{Descripción}{El sistema deberá mostrar al pasajero el promedio de calificaciones, cantidad de viajes evaluados y valoraciones permitidas del conductor antes de confirmar una reserva.}
\CampoRF{Actor principal}{Pasajero universitario.}
\CampoRF{Entradas}{Consulta del perfil del conductor asociado a una ruta.}
\CampoRF{Salidas}{Resumen de reputación e historial de calificaciones.}
\CampoRF{Precondición}{El conductor posee un perfil visible y la ruta está disponible.}
\CampoRF{Postcondición}{La información se presenta sin revelar identidades de quienes emitieron las evaluaciones.}
\CampoRF{Prioridad}{Must have.}
\CampoRF{Método de verificación}{Pruebas con conductores sin calificaciones, con múltiples evaluaciones, comentarios moderados y acceso previo a la reserva.}
\CampoRF{Evidencia}{\texttt{EV-CUE-001; EV-ENT-003; EV-ENT-009; EV-ENT-014.}}
\end{tabularx}
\end{RFBox}

\begin{RFBox}{RF-25: Panel estadístico institucional}
\small
\begin{tabularx}{\linewidth}{@{}L{3.15cm}Y@{}}
\CampoRF{Descripción}{El sistema deberá proporcionar al administrador un panel con indicadores agregados de movilidad, sin exponer datos personales ni recorridos individuales identificables.}
\CampoRF{Actor principal}{Administrador del sistema y personal institucional autorizado.}
\CampoRF{Entradas}{Periodo, zona general, tipo de indicador y filtros autorizados.}
\CampoRF{Salidas}{Panel con métricas agregadas y opciones de exportación permitidas.}
\CampoRF{Precondición}{El usuario posee permisos administrativos y existen registros anonimizados suficientes.}
\CampoRF{Postcondición}{Los indicadores se muestran o exportan con controles de anonimización y sin datos identificables.}
\CampoRF{Prioridad}{Could have.}
\CampoRF{Método de verificación}{Pruebas con datos sintéticos, comparación con cálculos de referencia y revisión documental de anonimización.}
\CampoRF{Evidencia}{\texttt{EV-CUE-001; EV-DOC-003; posteriormente EV-VAL-XXX.}}
\end{tabularx}
\end{RFBox}

\subsection{Cobertura y control de trazabilidad}

El catálogo contiene 25 requisitos funcionales: nueve clasificados como \emph{Must have}, doce como \emph{Should have} y cuatro como \emph{Could have}. Todos poseen una fuente identificada. Los requisitos RF-22 y RF-25 incorporan inicialmente evidencia documental porque su origen corresponde a obligaciones metodológicas y éticas; antes de cerrar la validación deberán complementarse con actas \texttt{EV-VAL-XXX} obtenidas durante los walkthroughs con usuarios técnicos, usuarios no técnicos y representantes autorizados.

La relación entre cada RF, historia de usuario, criterio de aceptación, caso de uso, componente y mockup se consolidará posteriormente en la matriz extendida de trazabilidad. La existencia de un código de evidencia no sustituye la verificación física del archivo correspondiente en la zona pública o restringida del repositorio.


% =====================================================
% 3.5 REQUISITOS NO FUNCIONALES
% =====================================================
\subsection{Requisitos no funcionales}

Los requisitos no funcionales definen las condiciones medibles de calidad bajo las cuales RoutTech deberá prestar sus funciones. La clasificación se basa en las nueve características del modelo ISO/IEC 25010:2023 \cite{iso25010}: adecuación funcional, eficiencia de desempeño, compatibilidad, capacidad de interacción, fiabilidad, seguridad, mantenibilidad, portabilidad y flexibilidad. Además, se incorpora la explicabilidad como requisito obligatorio del componente de recomendación, de acuerdo con la guía de la Entrega 3 y la literatura especializada \cite{uteq2026,chazette2020,chazette2021,chazette2022}.

Los valores objetivo se consideran criterios de aceptación de la versión académica y deberán comprobarse con pruebas repetibles sobre el MVP, inspecciones técnicas o sesiones de validación. La evidencia \texttt{EV-DOC-004} corresponde al catálogo técnico de requisitos no funcionales elaborado por el equipo; las obligaciones de protección de datos y geolocalización se respaldan además mediante \texttt{EV-DOC-002} y \texttt{EV-DOC-003}.

\begin{table}[H]
\centering
\caption{Cobertura de características de calidad en los requisitos no funcionales}
\label{tab:cobertura-iso25010}
\small
\rowcolors{2}{DocLight}{white}
\begin{tabularx}{\textwidth}{L{4.2cm} L{3.4cm} Y}
\rowcolor{DocNavy}
\textcolor{white}{\textbf{Característica}} &
\textcolor{white}{\textbf{RNF asociados}} &
\textcolor{white}{\textbf{Enfoque de calidad}}\\
Adecuación funcional & RNF-01 & Cobertura y corrección de los flujos críticos.\\
Eficiencia de desempeño & RNF-02, RNF-03 & Tiempo de respuesta y entrega de notificaciones.\\
Compatibilidad & RNF-04 & Interoperabilidad con navegadores y dispositivos soportados.\\
Capacidad de interacción & RNF-05, RNF-06 & Usabilidad, aprendizaje y accesibilidad.\\
Fiabilidad & RNF-07, RNF-08 & Disponibilidad e integridad de operaciones críticas.\\
Seguridad & RNF-09, RNF-10, RNF-11 & Acceso, protección de datos y privacidad de ubicación.\\
Mantenibilidad & RNF-12 & Modularidad, documentación y pruebas.\\
Portabilidad & RNF-13 & Instalación reproducible en ambientes soportados.\\
Flexibilidad & RNF-14 & Incorporación controlada de nuevos factores o módulos.\\
Explicabilidad & RNF-15 & Comprensión de recomendaciones automáticas.\\
\end{tabularx}
\end{table}

\begin{RNFBox}{RNF-01: Cobertura de los flujos funcionales críticos}
\small
\begin{tabularx}{\linewidth}{@{}L{3.35cm}Y@{}}
\CampoRNF{Característica}{Adecuación funcional.}
\CampoRNF{Requisito}{El MVP deberá implementar correctamente los flujos críticos de autenticación, publicación de rutas, búsqueda, reserva, cancelación y consulta de reputación definidos como \emph{Must have}.}
\CampoRNF{Métrica y objetivo}{El 100\,\% de los casos de prueba de aceptación asociados con dichos flujos deberá finalizar con el resultado esperado y sin defectos bloqueantes.}
\CampoRNF{Prioridad}{Must have.}
\CampoRNF{Método de verificación}{Ejecución automatizada o documentada de la suite de aceptación y revisión de resultados por requisito.}
\CampoRNF{Evidencia}{\texttt{EV-DOC-001; EV-DOC-004; RF-01, RF-03, RF-04, RF-05, RF-06 y RF-24.}}
\end{tabularx}
\end{RNFBox}

\begin{RNFBox}{RNF-02: Tiempo de respuesta de búsqueda y recomendación}
\small
\begin{tabularx}{\linewidth}{@{}L{3.35cm}Y@{}}
\CampoRNF{Característica}{Eficiencia de desempeño.}
\CampoRNF{Requisito}{El sistema deberá responder las consultas de búsqueda y recomendación de rutas sin demoras que impidan la interacción normal del usuario.}
\CampoRNF{Métrica y objetivo}{El percentil 95 del tiempo de respuesta deberá ser menor o igual a 3 segundos con 100 usuarios virtuales concurrentes y un conjunto de prueba de al menos 5\,000 rutas sintéticas.}
\CampoRNF{Prioridad}{Must have.}
\CampoRNF{Método de verificación}{Prueba de carga reproducible con registro de percentiles, tasa de errores y configuración del entorno.}
\CampoRNF{Evidencia}{\texttt{EV-DOC-004; EV-CUE-001; RF-03 y RF-14.}}
\end{tabularx}
\end{RNFBox}

\begin{RNFBox}{RNF-03: Oportunidad de las notificaciones}
\small
\begin{tabularx}{\linewidth}{@{}L{3.35cm}Y@{}}
\CampoRNF{Característica}{Eficiencia de desempeño.}
\CampoRNF{Requisito}{El sistema deberá entregar oportunamente las notificaciones relacionadas con reservas, cancelaciones y cambios de horario.}
\CampoRNF{Métrica y objetivo}{Al menos el 95\,\% de las notificaciones generadas deberá quedar registrada como enviada en un máximo de 5 segundos desde el evento que la origina, bajo carga normal.}
\CampoRNF{Prioridad}{Should have.}
\CampoRNF{Método de verificación}{Pruebas temporizadas con eventos controlados y comparación entre la marca de tiempo del evento y la del envío.}
\CampoRNF{Evidencia}{\texttt{EV-ENT-012; EV-DOC-004; RF-12.}}
\end{tabularx}
\end{RNFBox}

\begin{RNFBox}{RNF-04: Compatibilidad con plataformas soportadas}
\small
\begin{tabularx}{\linewidth}{@{}L{3.35cm}Y@{}}
\CampoRNF{Característica}{Compatibilidad.}
\CampoRNF{Requisito}{La interfaz web deberá funcionar de manera consistente en navegadores modernos y adaptarse a dispositivos móviles de uso común.}
\CampoRNF{Métrica y objetivo}{El 100\,\% de los flujos \emph{Must have} deberá completarse sin errores funcionales en las dos versiones estables más recientes de Chrome, Edge y Firefox, y en resoluciones desde 360\,$\times$\,640 hasta 1920\,$\times$\,1080 píxeles.}
\CampoRNF{Prioridad}{Should have.}
\CampoRNF{Método de verificación}{Matriz de pruebas cruzadas por navegador, sistema operativo y resolución.}
\CampoRNF{Evidencia}{\texttt{EV-DOC-004; DEP-01; RF-01 a RF-06.}}
\end{tabularx}
\end{RNFBox}

\begin{RNFBox}{RNF-05: Facilidad de uso de las tareas principales}
\small
\begin{tabularx}{\linewidth}{@{}L{3.35cm}Y@{}}
\CampoRNF{Característica}{Capacidad de interacción.}
\CampoRNF{Requisito}{La interfaz deberá permitir que usuarios sin capacitación previa completen las tareas principales mediante instrucciones y controles comprensibles.}
\CampoRNF{Métrica y objetivo}{Al menos el 80\,\% de los participantes deberá completar autenticación, búsqueda, reserva y publicación de ruta sin ayuda directa; la tasa de errores críticos por tarea deberá ser menor o igual al 10\,\%.}
\CampoRNF{Prioridad}{Must have.}
\CampoRNF{Método de verificación}{Prueba de usabilidad moderada con registro de éxito, errores, tiempo y solicitudes de ayuda.}
\CampoRNF{Evidencia}{\texttt{EV-CUE-001; EV-DOC-004; posteriormente EV-VAL-XXX.}}
\end{tabularx}
\end{RNFBox}

\begin{RNFBox}{RNF-06: Accesibilidad de la interfaz}
\small
\begin{tabularx}{\linewidth}{@{}L{3.35cm}Y@{}}
\CampoRNF{Característica}{Capacidad de interacción.}
\CampoRNF{Requisito}{Las pantallas prioritarias deberán ser perceptibles y operables mediante teclado, etiquetas textuales y contrastes suficientes.}
\CampoRNF{Métrica y objetivo}{Las pantallas \emph{Must have} deberán superar una revisión automática sin errores críticos de accesibilidad y permitir completar sus flujos principales utilizando únicamente teclado.}
\CampoRNF{Prioridad}{Should have.}
\CampoRNF{Método de verificación}{Inspección automática de accesibilidad, revisión manual del orden de foco y prueba de navegación por teclado.}
\CampoRNF{Evidencia}{\texttt{EV-DOC-001; EV-DOC-004; posteriormente EV-VAL-XXX.}}
\end{tabularx}
\end{RNFBox}

\begin{RNFBox}{RNF-07: Disponibilidad del servicio}
\small
\begin{tabularx}{\linewidth}{@{}L{3.35cm}Y@{}}
\CampoRNF{Característica}{Fiabilidad.}
\CampoRNF{Requisito}{El sistema deberá permanecer disponible durante los periodos de uso académico, excepto en ventanas de mantenimiento previamente comunicadas.}
\CampoRNF{Métrica y objetivo}{La disponibilidad mensual deberá ser igual o superior al 99\,\%, excluyendo mantenimientos programados anunciados con al menos 24 horas de anticipación.}
\CampoRNF{Prioridad}{Should have.}
\CampoRNF{Método de verificación}{Monitoreo de disponibilidad y cálculo mensual a partir de registros de inicio y fin de interrupciones.}
\CampoRNF{Evidencia}{\texttt{EV-DOC-004; DEP-06.}}
\end{tabularx}
\end{RNFBox}

\begin{RNFBox}{RNF-08: Integridad de reservas y cupos}
\small
\begin{tabularx}{\linewidth}{@{}L{3.35cm}Y@{}}
\CampoRNF{Característica}{Fiabilidad.}
\CampoRNF{Requisito}{Las operaciones de reserva, cancelación y actualización de cupos deberán conservar la consistencia de los datos ante errores, concurrencia o interrupciones.}
\CampoRNF{Métrica y objetivo}{En 100 ejecuciones concurrentes de reserva sobre el último cupo disponible no deberá registrarse sobreventa; toda operación fallida deberá revertirse sin dejar registros parciales.}
\CampoRNF{Prioridad}{Must have.}
\CampoRNF{Método de verificación}{Pruebas de concurrencia, transacciones interrumpidas y recuperación, comparando cupos y reservas antes y después de cada ejecución.}
\CampoRNF{Evidencia}{\texttt{EV-DOC-004; RF-05, RF-06 y RF-21.}}
\end{tabularx}
\end{RNFBox}

\begin{RNFBox}{RNF-09: Control de acceso y bloqueo de intentos}
\small
\begin{tabularx}{\linewidth}{@{}L{3.35cm}Y@{}}
\CampoRNF{Característica}{Seguridad.}
\CampoRNF{Requisito}{El sistema deberá impedir el acceso a usuarios no autenticados y limitar los intentos consecutivos de inicio de sesión.}
\CampoRNF{Métrica y objetivo}{Después de cinco intentos fallidos consecutivos, la cuenta o el origen de la solicitud deberá bloquearse durante 15 minutos; ningún rol podrá ejecutar funciones no autorizadas en la matriz de permisos.}
\CampoRNF{Prioridad}{Must have.}
\CampoRNF{Método de verificación}{Pruebas de autenticación, fuerza bruta controlada y autorización por rol.}
\CampoRNF{Evidencia}{\texttt{EV-CUE-001; EV-ENT-001; EV-ENT-006; EV-DOC-004.}}
\end{tabularx}
\end{RNFBox}

\begin{RNFBox}{RNF-10: Protección de credenciales y datos personales}
\small
\begin{tabularx}{\linewidth}{@{}L{3.35cm}Y@{}}
\CampoRNF{Característica}{Seguridad.}
\CampoRNF{Requisito}{El sistema deberá proteger las credenciales y los datos personales durante su transmisión y almacenamiento, aplicando minimización y separación de información identificable.}
\CampoRNF{Métrica y objetivo}{Ninguna contraseña deberá almacenarse en texto plano; todas las comunicaciones autenticadas deberán utilizar cifrado TLS 1.2 o superior; las pruebas de inspección no deberán revelar datos personales en registros de aplicación.}
\CampoRNF{Prioridad}{Must have.}
\CampoRNF{Método de verificación}{Inspección de configuración, revisión del esquema de almacenamiento, análisis de registros y prueba de tráfico cifrado.}
\CampoRNF{Evidencia}{\texttt{EV-DOC-003; EV-DOC-004;} Ley Orgánica de Protección de Datos Personales \cite{lopdp2021}.}
\end{tabularx}
\end{RNFBox}

\begin{RNFBox}{RNF-11: Privacidad y revocación de geolocalización}
\small
\begin{tabularx}{\linewidth}{@{}L{3.35cm}Y@{}}
\CampoRNF{Característica}{Seguridad.}
\CampoRNF{Requisito}{La ubicación de una persona solo podrá compartirse durante un viaje activo, con autorización explícita y revocable, utilizando una representación general o aproximada.}
\CampoRNF{Métrica y objetivo}{El 100\,\% de las solicitudes sin consentimiento deberá ser rechazado; la visualización deberá cesar en un máximo de 5 segundos tras la revocación o cierre del viaje; no se mostrarán domicilios ni coordenadas exactas en el MVP.}
\CampoRNF{Prioridad}{Must have.}
\CampoRNF{Método de verificación}{Pruebas de consentimiento, revocación, expiración y revisión de precisión de los datos mostrados.}
\CampoRNF{Evidencia}{\texttt{EV-ENT-003; EV-ENT-010; EV-ENT-014; EV-DOC-002; RF-17 y RF-19.}}
\end{tabularx}
\end{RNFBox}

\begin{RNFBox}{RNF-12: Mantenibilidad y cobertura de pruebas}
\small
\begin{tabularx}{\linewidth}{@{}L{3.35cm}Y@{}}
\CampoRNF{Característica}{Mantenibilidad.}
\CampoRNF{Requisito}{El código deberá organizarse en módulos con responsabilidades separadas y mantenerse sincronizado con la documentación técnica.}
\CampoRNF{Métrica y objetivo}{La capa de lógica de negocio deberá alcanzar una cobertura de pruebas automatizadas igual o superior al 70\,\%; ninguna función deberá superar una complejidad ciclomática de 15 sin justificación registrada.}
\CampoRNF{Prioridad}{Should have.}
\CampoRNF{Método de verificación}{Informe de cobertura, análisis estático e inspección de correspondencia entre módulos, README y diagramas.}
\CampoRNF{Evidencia}{\texttt{EV-DOC-004; DEP-06.}}
\end{tabularx}
\end{RNFBox}

\begin{RNFBox}{RNF-13: Despliegue reproducible}
\small
\begin{tabularx}{\linewidth}{@{}L{3.35cm}Y@{}}
\CampoRNF{Característica}{Portabilidad.}
\CampoRNF{Requisito}{El MVP deberá poder instalarse en un equipo limpio mediante instrucciones versionadas y dependencias declaradas.}
\CampoRNF{Métrica y objetivo}{Una persona con conocimientos técnicos deberá completar el despliegue local en Windows 11 o una distribución Linux soportada en un máximo de 15 minutos, utilizando como máximo cinco comandos después de instalar los prerrequisitos.}
\CampoRNF{Prioridad}{Should have.}
\CampoRNF{Método de verificación}{Prueba de instalación desde un clon nuevo del repositorio y registro del tiempo, comandos e incidencias.}
\CampoRNF{Evidencia}{\texttt{EV-DOC-001; EV-DOC-004; DEP-06.}}
\end{tabularx}
\end{RNFBox}

\begin{RNFBox}{RNF-14: Extensibilidad del componente de recomendación}
\small
\begin{tabularx}{\linewidth}{@{}L{3.35cm}Y@{}}
\CampoRNF{Característica}{Flexibilidad.}
\CampoRNF{Requisito}{La arquitectura deberá permitir incorporar o retirar factores de recomendación sin modificar los módulos de autenticación, reservas y notificaciones.}
\CampoRNF{Métrica y objetivo}{La incorporación de un nuevo factor de ponderación deberá requerir cambios únicamente en el módulo de recomendación y su configuración, sin producir regresiones en el 100\,\% de las pruebas de los módulos no relacionados.}
\CampoRNF{Prioridad}{Could have.}
\CampoRNF{Método de verificación}{Escenario de cambio controlado, revisión de archivos modificados y ejecución de pruebas de regresión.}
\CampoRNF{Evidencia}{\texttt{EV-DOC-004; RF-03, RF-22 y RF-23.}}
\end{tabularx}
\end{RNFBox}

\begin{RNFBox}{RNF-15: Explicabilidad de las recomendaciones}
\small
\begin{tabularx}{\linewidth}{@{}L{3.35cm}Y@{}}
\CampoRNF{Característica}{Explicabilidad.}
\CampoRNF{Requisito}{Cada recomendación automática deberá presentar una explicación breve de los factores principales que influyeron en la coincidencia y aclarar que la persona usuaria puede rechazarla.}
\CampoRNF{Métrica y objetivo}{La explicación deberá mostrarse en un máximo de 2 segundos, contener como máximo 60 palabras y alcanzar una puntuación media de comprensión igual o superior a 4 sobre 5 en una escala Likert aplicada a participantes técnicos y no técnicos.}
\CampoRNF{Prioridad}{Must have.}
\CampoRNF{Método de verificación}{Dos rondas de walkthrough con al menos tres usuarios técnicos y tres no técnicos por ronda, registrando tiempo, comprensión, utilidad y observaciones.}
\CampoRNF{Evidencia}{\texttt{EV-DOC-001; EV-DOC-002; RF-22; posteriormente EV-VAL-XXX.}}
\end{tabularx}
\end{RNFBox}

\subsubsection{Control de cuantificación y trazabilidad de los RNF}

Los 15 requisitos no funcionales poseen una métrica, un valor objetivo y un método de verificación. La matriz de trazabilidad deberá relacionarlos con las evidencias, requisitos funcionales afectados, casos de uso, componentes y pruebas. Los códigos de validación aún no ejecutados se muestran como \texttt{EV-VAL-XXX}; deberán sustituirse por los identificadores reales de las actas y grabaciones antes del cierre de la Entrega 3. La existencia de una meta técnica no implica que el resultado ya haya sido alcanzado: el cumplimiento se declarará únicamente después de ejecutar las pruebas y conservar sus salidas en el repositorio.



% =====================================================
% 3.6 REQUISITOS LEGALES
% =====================================================
\subsection{Requisitos legales y de protección de datos}

RoutTech tratará datos personales vinculados con identidad institucional, información de contacto, preferencias de viaje, reputación y referencias generales de ubicación. Por esta razón, el diseño deberá observar la Ley Orgánica de Protección de Datos Personales del Ecuador (LOPDP) \cite{lopdp2021}. La Tabla~\ref{tab:requisitos-legales} convierte las obligaciones aplicables en requisitos verificables y las relaciona con los RF y RNF que las implementan. Esta relación se incorporará posteriormente en la matriz extendida \emph{Ley--Artículo--RF/RNF/RD}.

\begin{longtable}{C{1.35cm} L{2.25cm} L{5.75cm} L{4.25cm}}
\caption{Requisitos legales y trazabilidad con la LOPDP}\label{tab:requisitos-legales}\\
\rowcolor{DocNavy}
\textcolor{white}{\textbf{ID}} &
\textcolor{white}{\textbf{Base legal}} &
\textcolor{white}{\textbf{Requisito legal verificable}} &
\textcolor{white}{\textbf{Trazabilidad inicial}}\\
\endfirsthead
\rowcolor{DocNavy}
\textcolor{white}{\textbf{ID}} &
\textcolor{white}{\textbf{Base legal}} &
\textcolor{white}{\textbf{Requisito legal verificable}} &
\textcolor{white}{\textbf{Trazabilidad inicial}}\\
\endhead
RL-01 & Arts. 7 y 8: legitimidad y consentimiento. & Antes de tratar datos cuya base legitimadora sea el consentimiento, el sistema deberá obtener una manifestación libre, específica, informada e inequívoca. La aceptación de la cuenta no deberá implicar automáticamente la autorización de funciones opcionales, como compartir ubicación. & RF-01, RF-17, RF-19; RNF-10 y RNF-11; EV-DOC-002 y EV-DOC-003.\\
RL-02 & Art. 10: finalidad, minimización y proporcionalidad. & El sistema deberá recopilar únicamente los datos necesarios para autenticar, coordinar y evaluar viajes. No almacenará domicilios exactos ni coordenadas permanentes; las ubicaciones se representarán mediante campus, sector, barrio, parroquia o punto general autorizado. & RF-03, RF-04, RF-17, RF-19 y RF-23; RNF-10 y RNF-11; EV-DOC-002 y EV-DOC-003.\\
RL-03 & Art. 10: calidad, exactitud y limitación de conservación. & Los datos editables deberán poder actualizarse y los registros personales se conservarán únicamente durante el periodo definido en el Plan de Gestión de Datos. Al finalizar dicho periodo se aplicará eliminación segura o anonimización irreversible, según corresponda. & RF-07, RF-08 y RF-23; RNF-10; EV-DOC-002 y EV-DOC-003.\\
RL-04 & Art. 12: derecho a la información. & Antes de la recolección, RoutTech deberá presentar un aviso comprensible que indique finalidades, base legal, datos tratados, tiempo de conservación, destinatarios, mecanismos de revocación, ejercicio de derechos y existencia de recomendaciones automatizadas. & RF-01, RF-07, RF-22 y RF-23; RNF-10 y RNF-15; EV-DOC-002 y EV-DOC-003.\\
RL-05 & Art. 8 y art. 12, nums. 13--16: revocación y ejercicio de derechos. & El sistema deberá ofrecer mecanismos accesibles para revocar consentimientos opcionales y para solicitar acceso, rectificación, eliminación, oposición o limitación del tratamiento. La revocación de ubicación deberá ejecutarse sin afectar el tratamiento previo realizado lícitamente. & RF-07, RF-17 y RF-19; RNF-11; EV-DOC-002 y EV-DOC-003.\\
RL-06 & Art. 10, lit. j: seguridad de datos personales. & RoutTech deberá aplicar medidas técnicas y organizativas para proteger confidencialidad, integridad y disponibilidad, incluyendo autenticación, autorización por roles, cifrado, registros de auditoría y gestión documentada de incidentes. & RF-01, RF-02, RF-13 y RF-16; RNF-08, RNF-09, RNF-10, RNF-11 y RNF-12; EV-DOC-003 y EV-DOC-004.\\
RL-07 & Art. 12, nums. 14 y 17: decisiones automatizadas. & Las recomendaciones deberán identificarse como sugerencias no obligatorias, mostrar los factores principales que las originaron y permitir que la persona usuaria las rechace sin que se cree una reserva automática. & RF-03, RF-22 y RF-23; RNF-14 y RNF-15; EV-DOC-001, EV-DOC-002 y posteriormente EV-VAL-XXX.\\
\end{longtable}

\subsubsection{Criterios de verificación legal}

La verificación de estos requisitos se realizará mediante inspección de interfaces, revisión de políticas, pruebas de consentimiento y revocación, análisis de permisos, revisión de registros y walkthrough con usuarios. El cumplimiento legal no se declarará únicamente por la existencia de esta tabla: cada fila deberá quedar vinculada con pruebas y evidencias reales en la matriz de trazabilidad.

% =====================================================
% 3.7 RESTRICCIONES DE DISEÑO
% =====================================================
\subsection{Restricciones de diseño}

Las restricciones de diseño delimitan decisiones que el equipo deberá respetar durante el desarrollo del MVP. No describen servicios ofrecidos al usuario, sino condiciones obligatorias de construcción, despliegue o tratamiento de información. La Tabla~\ref{tab:restricciones-diseno} contiene ocho restricciones iniciales, superando el mínimo de cinco solicitado para la Entrega 3.

\begin{longtable}{C{1.35cm} L{3.55cm} L{7.2cm} L{2.15cm}}
\caption{Restricciones de diseño de RoutTech}\label{tab:restricciones-diseno}\\
\rowcolor{DocNavy}
\textcolor{white}{\textbf{ID}} &
\textcolor{white}{\textbf{Nombre}} &
\textcolor{white}{\textbf{Restricción}} &
\textcolor{white}{\textbf{Evidencia}}\\
\endfirsthead
\rowcolor{DocNavy}
\textcolor{white}{\textbf{ID}} &
\textcolor{white}{\textbf{Nombre}} &
\textcolor{white}{\textbf{Restricción}} &
\textcolor{white}{\textbf{Evidencia}}\\
\endhead
RD-01 & Plataforma web adaptable & El MVP se implementará como aplicación web adaptable a escritorio y dispositivos móviles. No se exigirá una aplicación nativa independiente durante esta entrega. & EV-DOC-001.\\
RD-02 & Acceso institucional & El acceso operativo estará limitado a cuentas verificadas de la comunidad universitaria y a roles administrativos autorizados. Las demostraciones externas utilizarán cuentas ficticias. & EV-ENT-001; EV-ENT-006.\\
RD-03 & Ausencia de pagos & La versión académica no procesará tarifas, transferencias, comisiones ni pagos entre pasajeros y conductores. & EV-DOC-001.\\
RD-04 & Geolocalización generalizada & El MVP no almacenará domicilios exactos ni realizará seguimiento permanente. Solo manejará ubicaciones generales y puntos de encuentro autorizados; cualquier ubicación temporal deberá ser revocable. & EV-DOC-002; RF-17 y RF-19.\\
RD-05 & Datos sintéticos en demostraciones & La base de datos del MVP, los mockups, las capturas y el video de demostración utilizarán datos ficticios o sintéticos. Los datos primarios de investigación permanecerán separados del entorno de ejecución. & EV-DOC-002 y EV-DOC-003.\\
RD-06 & Participantes adultos & La fase académica y la validación descritas en este ERS se limitarán a personas mayores de edad. No se diseñarán flujos específicos para niñas, niños o adolescentes en esta versión. & EV-DOC-002.\\
RD-07 & Recomendación no vinculante & El módulo de recomendación no podrá confirmar reservas, cancelar viajes ni modificar rutas por decisión automática. La persona usuaria conservará la decisión final. & RF-03, RF-21 y RF-22; RNF-15.\\
RD-08 & Separación modular & Autenticación, usuarios, rutas, reservas, notificaciones y recomendación deberán mantenerse como módulos con interfaces definidas, de modo que los cambios del algoritmo de recomendación no alteren los módulos transaccionales. & RNF-12 y RNF-14; EV-DOC-004.\\
\end{longtable}

\subsubsection{Control de trazabilidad de requisitos legales y restricciones}

En la Sección 5, cada requisito legal y cada restricción de diseño se relacionará con el objetivo del proyecto, stakeholder, evidencia, RF o RNF afectado, caso de uso, historia de usuario, criterio de aceptación, componente y mockup correspondiente. Los identificadores \texttt{RL-XX} y \texttt{RD-XX} se conservarán sin cambios en el archivo \texttt{04\_Trazabilidad/matriz\_trazabilidad.csv}.



% =====================================================
% 3.8 HISTORIAS DE USUARIO
% =====================================================
\subsection{Historias de usuario de prioridad Must have}

Las historias de usuario se redactan con la estructura Connextra y se derivan exclusivamente de los requisitos funcionales clasificados como \emph{Must have}. Cada historia conserva la trazabilidad hacia su RF y hacia los identificadores de evidencia ya asociados en el catálogo funcional. La Tabla~\ref{tab:resumen-hu-must} presenta la correspondencia general.

\begin{longtable}{C{1.25cm} C{1.35cm} L{9.15cm} C{1.80cm}}
\caption{Historias de usuario asociadas con requisitos Must have}\label{tab:resumen-hu-must}\\
\rowcolor{DocNavy}
\textcolor{white}{\textbf{HU}} &
\textcolor{white}{\textbf{RF}} &
\textcolor{white}{\textbf{Historia Connextra}} &
\textcolor{white}{\textbf{Prioridad}}\\
\endfirsthead
\rowcolor{DocNavy}
\textcolor{white}{\textbf{HU}} &
\textcolor{white}{\textbf{RF}} &
\textcolor{white}{\textbf{Historia Connextra}} &
\textcolor{white}{\textbf{Prioridad}}\\
\endhead
HU-01 & RF-01 & Como miembro registrado de la comunidad UTEQ, quiero autenticarme con mis credenciales institucionales, para acceder de forma segura a las funciones autorizadas. & Must\\
HU-02 & RF-02 & Como administrador del sistema, quiero asignar roles y permisos a las cuentas, para controlar qué funciones puede utilizar cada tipo de usuario. & Must\\
HU-03 & RF-03 & Como pasajero universitario, quiero recibir rutas compatibles con mi origen, destino, horario y preferencias, para elegir una alternativa adecuada sin que se reserve automáticamente. & Must\\
HU-04 & RF-04 & Como conductor universitario verificado, quiero publicar una ruta con horario, cupos y punto de encuentro, para ofrecer lugares disponibles a otros miembros de la comunidad. & Must\\
HU-05 & RF-05 & Como pasajero universitario, quiero reservar un cupo disponible, para asegurar mi participación en el viaje seleccionado. & Must\\
HU-06 & RF-12 & Como participante de un viaje, quiero recibir notificaciones sobre solicitudes, confirmaciones, cambios y cancelaciones, para mantenerme informado oportunamente. & Must\\
HU-07 & RF-15 & Como conductor universitario verificado, quiero registrar el vehículo que utilizaré, para asociarlo con las rutas que publique. & Must\\
HU-08 & RF-16 & Como nuevo usuario, quiero verificar mi pertenencia institucional, para habilitar las funciones de movilidad colaborativa de RoutTech. & Must\\
HU-09 & RF-24 & Como pasajero universitario, quiero consultar la reputación del conductor antes de reservar, para tomar una decisión con mayor información y confianza. & Must\\
\end{longtable}

\begin{HUBox}{HU-01: Autenticarse con credenciales institucionales}
\small
\begin{tabularx}{\linewidth}{@{}L{3.0cm}Y@{}}
\CampoHU{Historia}{Como miembro registrado de la comunidad UTEQ, quiero autenticarme con mis credenciales institucionales, para acceder de forma segura a las funciones autorizadas.}
\CampoHU{Trazabilidad}{RF-01; RNF-09 y RNF-10; \texttt{EV-CUE-001, EV-ENT-001 y EV-ENT-005}.}
\CampoHU{Valor esperado}{Evitar accesos no autorizados y presentar al usuario únicamente las funciones de su rol.}
\end{tabularx}
\end{HUBox}

\begin{HUBox}{HU-02: Administrar roles y permisos}
\small
\begin{tabularx}{\linewidth}{@{}L{3.0cm}Y@{}}
\CampoHU{Historia}{Como administrador del sistema, quiero asignar roles y permisos a las cuentas, para controlar qué funciones puede utilizar cada tipo de usuario.}
\CampoHU{Trazabilidad}{RF-02; RNF-09; \texttt{EV-ENT-001 y EV-ENT-006}.}
\CampoHU{Valor esperado}{Aplicar separación de responsabilidades y reducir el uso indebido de funciones administrativas o de conducción.}
\end{tabularx}
\end{HUBox}

\begin{HUBox}{HU-03: Obtener recomendaciones de rutas compatibles}
\small
\begin{tabularx}{\linewidth}{@{}L{3.0cm}Y@{}}
\CampoHU{Historia}{Como pasajero universitario, quiero recibir rutas compatibles con mi origen, destino, horario y preferencias, para elegir una alternativa adecuada sin que se reserve automáticamente.}
\CampoHU{Trazabilidad}{RF-03; RNF-02 y RNF-15; \texttt{EV-CUE-001, EV-ENT-003, EV-ENT-007 y EV-ENT-010}.}
\CampoHU{Valor esperado}{Reducir el esfuerzo de búsqueda y mantener la decisión final bajo control de la persona usuaria.}
\end{tabularx}
\end{HUBox}

\begin{HUBox}{HU-04: Publicar una ruta}
\small
\begin{tabularx}{\linewidth}{@{}L{3.0cm}Y@{}}
\CampoHU{Historia}{Como conductor universitario verificado, quiero publicar una ruta con horario, cupos y punto de encuentro, para ofrecer lugares disponibles a otros miembros de la comunidad.}
\CampoHU{Trazabilidad}{RF-04; RF-15 y RF-16 como precondiciones; \texttt{EV-ENT-002, EV-ENT-012 y EV-ENT-015}.}
\CampoHU{Valor esperado}{Crear una oferta de viaje verificable, vigente y disponible para búsqueda y reserva.}
\end{tabularx}
\end{HUBox}

\begin{HUBox}{HU-05: Reservar un cupo}
\small
\begin{tabularx}{\linewidth}{@{}L{3.0cm}Y@{}}
\CampoHU{Historia}{Como pasajero universitario, quiero reservar un cupo disponible, para asegurar mi participación en el viaje seleccionado.}
\CampoHU{Trazabilidad}{RF-05; RNF-08; \texttt{EV-CUE-001, EV-ENT-004 y EV-ENT-011}.}
\CampoHU{Valor esperado}{Confirmar la participación sin superar la capacidad declarada ni duplicar reservas.}
\end{tabularx}
\end{HUBox}

\begin{HUBox}{HU-06: Recibir notificaciones de viaje}
\small
\begin{tabularx}{\linewidth}{@{}L{3.0cm}Y@{}}
\CampoHU{Historia}{Como participante de un viaje, quiero recibir notificaciones sobre solicitudes, confirmaciones, cambios y cancelaciones, para mantenerme informado oportunamente.}
\CampoHU{Trazabilidad}{RF-12; RNF-03; \texttt{EV-ENT-006, EV-ENT-011 y EV-ENT-012}.}
\CampoHU{Valor esperado}{Reducir descoordinaciones y permitir una reacción oportuna ante cambios del viaje.}
\end{tabularx}
\end{HUBox}

\begin{HUBox}{HU-07: Registrar un vehículo}
\small
\begin{tabularx}{\linewidth}{@{}L{3.0cm}Y@{}}
\CampoHU{Historia}{Como conductor universitario verificado, quiero registrar el vehículo que utilizaré, para asociarlo con las rutas que publique.}
\CampoHU{Trazabilidad}{RF-15; RF-16 como precondición; \texttt{EV-ENT-002 y EV-ENT-015}.}
\CampoHU{Valor esperado}{Mantener información mínima y consistente del vehículo disponible para la publicación de rutas.}
\end{tabularx}
\end{HUBox}

\begin{HUBox}{HU-08: Verificar la pertenencia institucional}
\small
\begin{tabularx}{\linewidth}{@{}L{3.0cm}Y@{}}
\CampoHU{Historia}{Como nuevo usuario, quiero verificar mi pertenencia institucional, para habilitar las funciones de movilidad colaborativa de RoutTech.}
\CampoHU{Trazabilidad}{RF-16; RF-01 y RF-02; \texttt{EV-CUE-001, EV-ENT-001, EV-ENT-006 y EV-ENT-009}.}
\CampoHU{Valor esperado}{Restringir la plataforma a miembros autorizados de la comunidad universitaria.}
\end{tabularx}
\end{HUBox}

\begin{HUBox}{HU-09: Consultar la reputación del conductor}
\small
\begin{tabularx}{\linewidth}{@{}L{3.0cm}Y@{}}
\CampoHU{Historia}{Como pasajero universitario, quiero consultar la reputación del conductor antes de reservar, para tomar una decisión con mayor información y confianza.}
\CampoHU{Trazabilidad}{RF-24; RF-09; \texttt{EV-CUE-001, EV-ENT-003, EV-ENT-009 y EV-ENT-014}.}
\CampoHU{Valor esperado}{Presentar un resumen útil sin revelar la identidad de quienes emitieron las evaluaciones.}
\end{tabularx}
\end{HUBox}

% =====================================================
% 3.9 VALIDACION INVEST
% =====================================================
\subsection{Verificación INVEST de las historias de usuario}

La revisión INVEST comprueba que cada historia sea independiente en su alcance, negociable en su solución, valiosa, estimable, suficientemente pequeña y verificable. Las dependencias funcionales se documentan como precondiciones y no impiden estimar ni probar la historia por separado.

\begin{longtable}{C{1.35cm} C{0.70cm} C{0.70cm} C{0.70cm} C{0.70cm} C{0.70cm} C{0.70cm} L{7.45cm}}
\caption{Validación INVEST de las historias Must have}\label{tab:invest-historias}\\
\rowcolor{DocNavy}
\textcolor{white}{\textbf{HU}} &
\textcolor{white}{\textbf{I}} &
\textcolor{white}{\textbf{N}} &
\textcolor{white}{\textbf{V}} &
\textcolor{white}{\textbf{E}} &
\textcolor{white}{\textbf{S}} &
\textcolor{white}{\textbf{T}} &
\textcolor{white}{\textbf{Observación}}\\
\endfirsthead
\rowcolor{DocNavy}
\textcolor{white}{\textbf{HU}} &
\textcolor{white}{\textbf{I}} &
\textcolor{white}{\textbf{N}} &
\textcolor{white}{\textbf{V}} &
\textcolor{white}{\textbf{E}} &
\textcolor{white}{\textbf{S}} &
\textcolor{white}{\textbf{T}} &
\textcolor{white}{\textbf{Observación}}\\
\endhead
HU-01 & Sí & Sí & Sí & Sí & Sí & Sí & Alcance limitado al inicio de sesión; verificable con credenciales válidas, inválidas, inactivas y bloqueadas.\\
HU-02 & Sí & Sí & Sí & Sí & Sí & Sí & Se centra en asignación y actualización de permisos; la autorización puede probarse por rol.\\
HU-03 & Sí & Sí & Sí & Sí & Sí & Sí & Produce una lista sin crear reservas; criterios de compatibilidad y orden pueden validarse con datos controlados.\\
HU-04 & Sí & Sí & Sí & Sí & Sí & Sí & La verificación y el vehículo son precondiciones documentadas; la publicación constituye un incremento funcional separado.\\
HU-05 & Sí & Sí & Sí & Sí & Sí & Sí & Se limita a reservar y actualizar cupos de forma atómica; incluye casos de concurrencia y ruta llena.\\
HU-06 & Sí & Sí & Sí & Sí & Sí & Sí & Cada tipo de evento puede ejecutarse y verificarse de forma independiente, incluyendo control de duplicados.\\
HU-07 & Sí & Sí & Sí & Sí & Sí & Sí & Comprende alta y mantenimiento mínimo del vehículo; es estimable y verificable sin publicar una ruta.\\
HU-08 & Sí & Sí & Sí & Sí & Sí & Sí & La respuesta del servicio institucional puede simularse para probar estados válido, rechazado, pendiente y no disponible.\\
HU-09 & Sí & Sí & Sí & Sí & Sí & Sí & Presenta un resumen previo a la reserva y puede verificarse con historiales vacíos, múltiples y moderados.\\
\end{longtable}

\subsubsection{Resultado de la revisión INVEST}

Las nueve historias cumplen los seis criterios en su formulación actual. Esta evaluación documental deberá contrastarse durante las sesiones de walkthrough; cualquier división, combinación o cambio de alcance aprobado se registrará en el control de cambios y se reflejará en la matriz de trazabilidad.

% =====================================================
% 3.10 CRITERIOS DE ACEPTACION GHERKIN
% =====================================================
\subsection{Criterios de aceptación en formato Gherkin}
\label{sec:criterios-gherkin}

Los criterios siguientes se redactan con la estructura \emph{Dado--Cuando--Entonces}. Cada escenario posee un identificador CA y se vincula con una historia y un requisito funcional. Los escenarios constituyen la base para los casos de prueba de aceptación y no representan resultados ya ejecutados.

\subsubsection{HU-01 / RF-01: Autenticación institucional}

\begin{GherkinBox}{CA-01.1: Inicio de sesión aceptado}
\textbf{Dado} que la cuenta institucional existe, está activa y no se encuentra bloqueada,\\
\textbf{cuando} el usuario ingresa un correo institucional y una contraseña válidos,\\
\textbf{entonces} el sistema inicia la sesión y muestra únicamente las funciones autorizadas para su rol.
\end{GherkinBox}

\begin{GherkinBox}{CA-01.2: Credenciales rechazadas}
\textbf{Dado} que una persona se encuentra en la pantalla de autenticación,\\
\textbf{cuando} ingresa credenciales inválidas,\\
\textbf{entonces} el sistema rechaza el acceso, no crea una sesión y registra el intento fallido.
\end{GherkinBox}

\subsubsection{HU-02 / RF-02: Roles y permisos}

\begin{GherkinBox}{CA-02.1: Asignación de rol autorizada}
\textbf{Dado} que el administrador inició sesión y seleccionó una cuenta existente,\\
\textbf{cuando} asigna un rol válido y confirma el cambio,\\
\textbf{entonces} el sistema actualiza los permisos y registra la modificación para auditoría.
\end{GherkinBox}

\begin{GherkinBox}{CA-02.2: Acceso a función restringida}
\textbf{Dado} que un usuario posee un rol sin permiso administrativo,\\
\textbf{cuando} intenta acceder a una función restringida,\\
\textbf{entonces} el sistema deniega la operación y conserva sin cambios la información protegida.
\end{GherkinBox}

\subsubsection{HU-03 / RF-03: Recomendación de rutas}

\begin{GherkinBox}{CA-03.1: Presentación de rutas compatibles}
\textbf{Dado} que el pasajero está autenticado y existen rutas activas con cupos,\\
\textbf{cuando} registra origen y destino generales, fecha, horario y preferencias,\\
\textbf{entonces} el sistema presenta una lista ordenada de rutas compatibles sin crear una reserva automática.
\end{GherkinBox}

\begin{GherkinBox}{CA-03.2: Ausencia de coincidencias}
\textbf{Dado} que no existe una ruta que cumpla los criterios indicados,\\
\textbf{cuando} el pasajero solicita recomendaciones,\\
\textbf{entonces} el sistema informa que no encontró coincidencias y permite modificar los criterios de búsqueda.
\end{GherkinBox}

\subsubsection{HU-04 / RF-04: Publicación de rutas}

\begin{GherkinBox}{CA-04.1: Ruta publicada correctamente}
\textbf{Dado} que el conductor está autenticado, verificado y posee un vehículo habilitado,\\
\textbf{cuando} completa los datos obligatorios con una fecha futura y una cantidad válida de cupos,\\
\textbf{entonces} el sistema crea la ruta con un identificador y estado activo, disponible para búsqueda.
\end{GherkinBox}

\begin{GherkinBox}{CA-04.2: Datos de publicación inválidos}
\textbf{Dado} que el conductor intenta publicar una ruta,\\
\textbf{cuando} omite un campo obligatorio, usa una fecha pasada o registra una capacidad inválida,\\
\textbf{entonces} el sistema no publica la ruta e indica los campos que deben corregirse.
\end{GherkinBox}

\subsubsection{HU-05 / RF-05: Reserva de cupos}

\begin{GherkinBox}{CA-05.1: Reserva confirmada}
\textbf{Dado} que el pasajero está autenticado y la ruta permanece activa con cupos suficientes,\\
\textbf{cuando} solicita un cupo y confirma la operación,\\
\textbf{entonces} el sistema registra la reserva y disminuye la disponibilidad de la ruta de forma atómica.
\end{GherkinBox}

\begin{GherkinBox}{CA-05.2: Último cupo solicitado simultáneamente}
\textbf{Dado} que una ruta posee un solo cupo disponible,\\
\textbf{cuando} dos solicitudes concurrentes intentan reservarlo,\\
\textbf{entonces} el sistema confirma únicamente una reserva y rechaza la otra sin generar sobreventa.
\end{GherkinBox}

\subsubsection{HU-06 / RF-12: Notificaciones de viaje}

\begin{GherkinBox}{CA-06.1: Notificación generada por un evento}
\textbf{Dado} que existe una reserva o ruta relacionada con usuarios identificados,\\
\textbf{cuando} ocurre una aceptación, rechazo, cambio, cancelación o recordatorio,\\
\textbf{entonces} el sistema genera la notificación correspondiente y registra su estado y marca temporal.
\end{GherkinBox}

\begin{GherkinBox}{CA-06.2: Control de notificaciones duplicadas}
\textbf{Dado} que un evento ya produjo una notificación para un destinatario,\\
\textbf{cuando} el mismo evento se procesa nuevamente con el mismo identificador,\\
\textbf{entonces} el sistema evita enviar un aviso duplicado y conserva la trazabilidad del procesamiento.
\end{GherkinBox}

\subsubsection{HU-07 / RF-15: Registro de vehículo}

\begin{GherkinBox}{CA-07.1: Vehículo registrado}
\textbf{Dado} que el conductor está autenticado y verificado,\\
\textbf{cuando} registra tipo, marca, modelo, color y capacidad mediante datos válidos,\\
\textbf{entonces} el sistema asocia el vehículo con su perfil y lo habilita para futuras publicaciones.
\end{GherkinBox}

\begin{GherkinBox}{CA-07.2: Vehículo con datos inválidos o duplicados}
\textbf{Dado} que el conductor se encuentra en el formulario de vehículo,\\
\textbf{cuando} indica una capacidad inválida o intenta registrar un vehículo duplicado,\\
\textbf{entonces} el sistema rechaza el registro y muestra la causa de validación.
\end{GherkinBox}

\subsubsection{HU-08 / RF-16: Verificación institucional}

\begin{GherkinBox}{CA-08.1: Identidad institucional validada}
\textbf{Dado} que el usuario completó el registro y aceptó las condiciones aplicables,\\
\textbf{cuando} el servicio institucional confirma que la identidad es válida,\\
\textbf{entonces} el sistema marca la cuenta como verificada y habilita las funciones autorizadas.
\end{GherkinBox}

\begin{GherkinBox}{CA-08.2: Verificación rechazada o no disponible}
\textbf{Dado} que una cuenta se encuentra pendiente de verificación,\\
\textbf{cuando} el servicio rechaza la identidad o no se encuentra disponible,\\
\textbf{entonces} el sistema mantiene restringidas las funciones de movilidad e informa el estado sin crear una verificación falsa.
\end{GherkinBox}

\subsubsection{HU-09 / RF-24: Historial de reputación}

\begin{GherkinBox}{CA-09.1: Reputación disponible antes de reservar}
\textbf{Dado} que el pasajero consulta una ruta disponible cuyo conductor posee evaluaciones,\\
\textbf{cuando} abre el resumen del perfil del conductor,\\
\textbf{entonces} el sistema muestra el promedio, la cantidad de viajes evaluados y las valoraciones permitidas sin revelar autores.
\end{GherkinBox}

\begin{GherkinBox}{CA-09.2: Conductor sin evaluaciones}
\textbf{Dado} que el conductor todavía no posee calificaciones,\\
\textbf{cuando} el pasajero consulta su reputación,\\
\textbf{entonces} el sistema informa que aún no existen evaluaciones y no presenta un promedio engañoso.
\end{GherkinBox}

\subsubsection{Control de trazabilidad de historias y criterios}

Los identificadores \texttt{HU-01} a \texttt{HU-09} y \texttt{CA-01.1} a \texttt{CA-09.2} se conservarán en la matriz extendida. Durante la validación, cada criterio deberá relacionarse con el caso de uso, componente, mockup y resultado de prueba correspondiente. Las modificaciones aprobadas deberán actualizar simultáneamente el RF, la historia, el criterio y la matriz de trazabilidad.




% =====================================================
% 3.11 MOCKUPS DE PRIORIDAD MUST HAVE
% =====================================================
\subsection{Prototipos de interfaz asociados a requisitos Must have}

Los prototipos de interfaz de alta fidelidad se identificarán con el prefijo \texttt{MU} y se mantendrán trazados a los requisitos funcionales, historias de usuario y criterios de aceptación correspondientes. Cada mockup deberá conservar una versión PNG o SVG dentro de \texttt{03\_Modelado/Mockups/}; el enlace de una herramienta externa no sustituye al archivo versionado en el repositorio.

La Tabla~\ref{tab:catalogo-mockups-must} establece el catálogo mínimo de pantallas que deberá incorporarse al documento. La imagen de cada interfaz se insertará posteriormente sin modificar sus identificadores ni relaciones de trazabilidad.

\begin{longtable}{C{1.20cm} L{3.35cm} L{2.15cm} L{2.25cm} L{4.95cm}}
\caption{Catálogo de mockups asociados a requisitos Must have}\label{tab:catalogo-mockups-must}\\
\rowcolor{DocNavy}
\textcolor{white}{\textbf{ID}} &
\textcolor{white}{\textbf{Interfaz}} &
\textcolor{white}{\textbf{RF / HU}} &
\textcolor{white}{\textbf{CA}} &
\textcolor{white}{\textbf{Propósito funcional}}\\
\endfirsthead
\rowcolor{DocNavy}
\textcolor{white}{\textbf{ID}} &
\textcolor{white}{\textbf{Interfaz}} &
\textcolor{white}{\textbf{RF / HU}} &
\textcolor{white}{\textbf{CA}} &
\textcolor{white}{\textbf{Propósito funcional}}\\
\endhead
MU-01 & Acceso institucional & RF-01 / HU-01 & CA-01.1, CA-01.2 & Permitir el ingreso con credenciales institucionales, presentar mensajes de validación y restringir el acceso cuando la cuenta no sea válida.\\
MU-02 & Administración de roles y permisos & RF-02 / HU-02 & CA-02.1, CA-02.2 & Permitir al administrador consultar cuentas, asignar roles y verificar las funciones habilitadas para cada perfil.\\
MU-03 & Búsqueda y recomendación de rutas & RF-03 / HU-03 & CA-03.1, CA-03.2 & Recopilar criterios generales de origen, destino, fecha, horario y preferencias; mostrar coincidencias sin crear reservas automáticas.\\
MU-04 & Publicación de ruta & RF-04 / HU-04 & CA-04.1, CA-04.2 & Registrar horario, cupos, puntos generales de encuentro y datos del recorrido, con validaciones antes de publicar.\\
MU-05 & Reserva de cupo & RF-05 / HU-05 & CA-05.1, CA-05.2 & Presentar la información esencial de la ruta y confirmar una reserva sin superar la capacidad disponible.\\
MU-06 & Centro de notificaciones & RF-12 / HU-06 & CA-06.1, CA-06.2 & Mostrar avisos de solicitudes, confirmaciones, cambios, cancelaciones y recordatorios, evitando duplicados.\\
MU-07 & Registro de vehículo & RF-15 / HU-07 & CA-07.1, CA-07.2 & Registrar y validar los datos mínimos del vehículo que será asociado con las rutas del conductor.\\
MU-08 & Verificación institucional & RF-16 / HU-08 & CA-08.1, CA-08.2 & Comunicar el estado de la verificación institucional y las funciones habilitadas, pendientes o restringidas.\\
MU-09 & Reputación del conductor & RF-24 / HU-09 & CA-09.1, CA-09.2 & Mostrar promedio, cantidad de viajes evaluados y valoraciones permitidas antes de confirmar una reserva.\\
\end{longtable}

\subsubsection{Estados mínimos representados en los mockups}

Cada interfaz deberá representar, cuando corresponda, los siguientes estados para facilitar su validación durante el walkthrough:

\begin{itemize}
  \item \textbf{Estado inicial:} pantalla disponible antes de ingresar información o ejecutar una acción.
  \item \textbf{Estado con datos válidos:} contenido completo y preparado para confirmar la operación.
  \item \textbf{Estado vacío:} ausencia de rutas, notificaciones, evaluaciones o resultados, sin mostrar información engañosa.
  \item \textbf{Estado de validación:} campos incompletos, datos incompatibles, permisos insuficientes o indisponibilidad temporal.
  \item \textbf{Estado de confirmación:} resultado exitoso de una autenticación, publicación, reserva, verificación o registro.
\end{itemize}

\subsubsection{Criterios generales de diseño de interfaz}

Los mockups deberán aplicar una jerarquía visual consistente, lenguaje comprensible y retroalimentación inmediata. Las ubicaciones se expresarán mediante sectores o puntos generales de encuentro; no se mostrarán domicilios particulares ni coordenadas exactas. Las recomendaciones deberán distinguirse de las decisiones confirmadas y, cuando se presente una coincidencia automática, la interfaz deberá comunicar los factores principales que la originaron, de acuerdo con el RNF de explicabilidad.

La revisión de los mockups se documentará mediante actas de walkthrough. Los cambios aprobados deberán actualizar el identificador del mockup, el requisito relacionado, la historia de usuario, los criterios de aceptación y la matriz de trazabilidad extendida.

\subsubsection{Mockups actuales disponibles para el MVP}

En la iteración actual del proyecto ya se dispone de un conjunto de pantallas funcionales de alta fidelidad que servirán como insumo directo para el walkthrough y para la demostración del MVP. Estas interfaces no sustituyen el catálogo de trazabilidad presentado en la Tabla~\ref{tab:catalogo-mockups-must}; más bien, constituyen la evidencia visual disponible a la fecha y permiten cubrir el flujo principal de autenticación, navegación, búsqueda, reservas, historial, notificaciones, administración, recomendaciones e incidencias.

Las pantallas actualmente disponibles cubren de forma directa los flujos de acceso institucional, exploración de rutas, reserva de cupos, revisión del historial, consulta de notificaciones, administración del sistema y explicación básica de recomendaciones. Los mockups específicos para publicación detallada de rutas, registro de vehículo y verificación institucional podrán añadirse en una siguiente iteración manteniendo los identificadores MU ya definidos.

\FiguraDocumento{imagenes/mockups/MU_Login.jpg}{0.26}{Mockup MU-01. Pantalla de acceso institucional de RoutTech.}{fig:mu01-login}

\FiguraDocumento{imagenes/mockups/MU_Registro.jpg}{0.26}{Pantalla complementaria de registro de cuenta asociada con el flujo de incorporación de nuevos usuarios.}{fig:mu-registro}

\FiguraDocumento{imagenes/mockups/MU_Recuperar.jpg}{0.26}{Pantalla de recuperación de cuenta vinculada con el soporte del acceso institucional.}{fig:mu-recuperar}

\FiguraDocumento{imagenes/mockups/MU_Inicio.jpg}{0.26}{Vista principal del usuario autenticado, desde la cual se accede a búsqueda de rutas, reservas, historial, perfil y funciones administrativas.}{fig:mu-inicio}

\FiguraDocumento{imagenes/mockups/MU_BuscarRutas.jpg}{0.26}{Mockup MU-03. Pantalla de búsqueda de rutas con filtros de origen, destino, fecha, hora y listado de coincidencias.}{fig:mu03-busqueda}

\FiguraDocumento{imagenes/mockups/MU_IA.jpg}{0.26}{Pantalla de recomendaciones IA asociada con el requisito de recomendación explicable de rutas.}{fig:mu-ia}

\FiguraDocumento{imagenes/mockups/MU_Reservas.jpg}{0.26}{Mockup MU-05. Pantalla de gestión de reservas activas, pendientes y canceladas.}{fig:mu05-reservas}

\FiguraDocumento{imagenes/mockups/MU_Historial.jpg}{0.26}{Pantalla de historial de viajes con opción de calificación y consulta de detalles.}{fig:mu-historial}

\FiguraDocumento{imagenes/mockups/MU_Notificaciones.jpg}{0.26}{Mockup MU-06. Centro de notificaciones con avisos de confirmación, cambios y cancelaciones.}{fig:mu06-notificaciones}

\FiguraDocumento{imagenes/mockups/MU_Perfil.jpg}{0.26}{Pantalla de perfil del usuario con datos personales, edición básica y cambio de contraseña.}{fig:mu-perfil}

\FiguraDocumento{imagenes/mockups/MU_Admin.jpg}{0.26}{Mockup MU-02. Panel de administración con indicadores operativos, incidencias y estadísticas generales.}{fig:mu02-admin}

\FiguraDocumento{imagenes/mockups/MU_Incidencia.jpg}{0.26}{Pantalla de reporte de incidencias asociada con el seguimiento de eventos durante el viaje.}{fig:mu-incidencia}

\FiguraDocumento{imagenes/mockups/MU_Reportes_01.jpg}{0.26}{Pantalla preliminar de reportes de movilidad con filtros temporales, indicadores y vista previa del reporte.}{fig:mu-reportes-01}

\FiguraDocumento{imagenes/mockups/MU_Reportes_02.jpg}{0.26}{Pantalla complementaria de reportes de movilidad con opciones de exportación en PDF y Excel.}{fig:mu-reportes-02}

\begin{longtable}{C{1.35cm} L{4.25cm} L{2.45cm} L{5.25cm}}
\caption{Cobertura actual de mockups disponibles para el MVP}\label{tab:cobertura-mockups-actual}\\
\rowcolor{DocNavy}
\textcolor{white}{\textbf{ID}} &
\textcolor{white}{\textbf{Pantalla actual}} &
\textcolor{white}{\textbf{Trazabilidad principal}} &
\textcolor{white}{\textbf{Observación de cobertura}}\\
\endfirsthead
\rowcolor{DocNavy}
\textcolor{white}{\textbf{ID}} &
\textcolor{white}{\textbf{Pantalla actual}} &
\textcolor{white}{\textbf{Trazabilidad principal}} &
\textcolor{white}{\textbf{Observación de cobertura}}\\
\endhead
MU-01 & Inicio de sesión / recuperación de cuenta & RF-01; HU-01; CA-01.1--CA-01.2 & Cobertura directa del acceso institucional y recuperación de credenciales.\\
MU-02 & Panel de administración & RF-02; HU-02 & Cobertura parcial; evidencia acceso administrativo e indicadores, pero la asignación detallada de roles podrá ampliarse.\\
MU-03 & Buscar rutas / recomendaciones IA & RF-03; HU-03; RNF-15 & Cobertura directa de búsqueda y soporte visual del componente de recomendación explicable.\\
MU-04 & Inicio del usuario autenticado & RF-04; HU-04 & Cobertura indirecta; existe acceso desde el panel principal, aunque el formulario específico de publicación deberá añadirse luego.\\
MU-05 & Mis reservas & RF-05; HU-05 & Cobertura directa del seguimiento y cancelación de reservas.\\
MU-06 & Notificaciones & RF-12; HU-06 & Cobertura directa del centro de notificaciones internas.\\
MU-07 & Perfil del usuario & RF-15; HU-07 & Cobertura parcial; sirve como base de gestión de datos, aunque el registro específico del vehículo todavía está pendiente.\\
MU-08 & Registro de cuenta & RF-16; HU-08 & Cobertura parcial; cubre el alta inicial, pero la verificación institucional explícita se documentará en una iteración posterior.\\
MU-09 & Historial de viajes / calificación & RF-24; HU-09 & Cobertura directa del flujo de consulta y calificación posterior al viaje.\\
EX-01 & Reportar incidencia & RF-11; RF-13 & Pantalla complementaria útil para walkthrough y cobertura del seguimiento de incidencias.\\
EX-02 & Reportes de movilidad & RF-20; RF-21 & Pantalla administrativa complementaria para visualizar indicadores y exportaciones.\\
\end{longtable}

\subsubsection{Observación sobre el video del MVP}

Las pantallas anteriores son suficientes para preparar un video corto del MVP. La grabación podrá estructurarse en una secuencia de acceso, búsqueda, recomendación, reserva, notificación, historial y administración. Cuando se prepare ese video, deberá enlazarse con los módulos M01 a M06 descritos en la Sección~6 y conservarse como evidencia audiovisual dentro del repositorio del proyecto.


% =====================================================
% 4. MODELADO UML DEL SISTEMA
% =====================================================
\section{Modelado UML del sistema}

El modelado de RoutTech representa de forma complementaria el comportamiento, la estructura y la distribución lógica del sistema. Los modelos se derivan del catálogo de requisitos, las historias de usuario y los criterios de aceptación definidos en la Sección 3. En esta etapa se consolida la especificación textual y se integran los diagramas principales de casos de uso, secuencia, actividad, componentes, paquetes y datos, manteniendo los identificadores definidos para su trazabilidad con requisitos, historias de usuario y criterios de aceptación.

\subsection{Modelo general de casos de uso}

El modelo general considera como actores primarios al pasajero universitario, al conductor universitario y al administrador del sistema. Como actores externos se incluyen el servicio de autenticación institucional, el servicio de mapas, el servicio de notificaciones y el módulo de recomendación. La Tabla~\ref{tab:catalogo-cu} resume los diez casos de uso detallados que cubren todos los requisitos funcionales de prioridad \emph{Must have} y un flujo adicional de cancelación.

\begin{longtable}{C{1.20cm} L{4.20cm} L{3.15cm} L{2.10cm} L{2.50cm}}
\caption{Catálogo de casos de uso detallados de RoutTech}\label{tab:catalogo-cu}\\
\rowcolor{DocNavy}
\textcolor{white}{\textbf{ID}} &
\textcolor{white}{\textbf{Caso de uso}} &
\textcolor{white}{\textbf{Actor principal}} &
\textcolor{white}{\textbf{RF}} &
\textcolor{white}{\textbf{HU / CA}}\\
\endfirsthead
\rowcolor{DocNavy}
\textcolor{white}{\textbf{ID}} &
\textcolor{white}{\textbf{Caso de uso}} &
\textcolor{white}{\textbf{Actor principal}} &
\textcolor{white}{\textbf{RF}} &
\textcolor{white}{\textbf{HU / CA}}\\
\endhead
CU-01 & Autenticar usuario institucional & Usuario registrado & RF-01 & HU-01 / CA-01.1, CA-01.2\\
CU-02 & Administrar roles y permisos & Administrador & RF-02 & HU-02 / CA-02.1, CA-02.2\\
CU-03 & Buscar y recomendar rutas & Pasajero universitario & RF-03, RF-14, RF-22, RF-23 & HU-03 / CA-03.1, CA-03.2\\
CU-04 & Publicar una ruta & Conductor universitario & RF-04, RF-15, RF-19 & HU-04 / CA-04.1, CA-04.2\\
CU-05 & Reservar un cupo & Pasajero universitario & RF-05, RF-24 & HU-05 / CA-05.1, CA-05.2\\
CU-06 & Cancelar una reserva o ruta & Pasajero o conductor & RF-06, RF-12, RF-21 & --\\
CU-07 & Gestionar notificaciones de viaje & Sistema & RF-12 & HU-06 / CA-06.1, CA-06.2\\
CU-08 & Registrar un vehículo & Conductor universitario & RF-15 & HU-07 / CA-07.1, CA-07.2\\
CU-09 & Verificar identidad institucional & Usuario y servicio UTEQ & RF-16 & HU-08 / CA-08.1, CA-08.2\\
CU-10 & Consultar reputación del conductor & Pasajero universitario & RF-24, RF-09 & HU-09 / CA-09.1, CA-09.2\\
\end{longtable}


El modelo gráfico de casos de uso se presenta en las Figuras~\ref{fig:cu-general-routtech} a~\ref{fig:cu-admin-routtech}. La vista general sintetiza el alcance funcional del sistema, mientras que las vistas por actor permiten distinguir con mayor claridad las responsabilidades del pasajero, del conductor y del administrador dentro del ecosistema de movilidad colaborativa.

\FiguraDocumento{03_Modelado/Diagramas_UML/Diagrama de caso de uso general.jpg}{0.95}{Diagrama general de casos de uso de RoutTech.}{fig:cu-general-routtech}

\FiguraDocumento{03_Modelado/Diagramas_UML/CU Pasajero.jpg}{0.92}{Casos de uso asociados con el actor pasajero universitario.}{fig:cu-pasajero-routtech}

\FiguraDocumento{03_Modelado/Diagramas_UML/CU Conductor.jpg}{0.92}{Casos de uso asociados con el actor conductor universitario.}{fig:cu-conductor-routtech}

\FiguraDocumento{03_Modelado/Diagramas_UML/CU Administración.jpg}{0.92}{Casos de uso asociados con el actor administrador del sistema.}{fig:cu-admin-routtech}

\subsection{Especificación textual de casos de uso}

\begin{CUBox}{CU-01: Autenticar usuario institucional}
\small
\begin{tabularx}{\linewidth}{@{}L{3.20cm}Y@{}}
\CampoCU{Objetivo}{Permitir que una cuenta institucional válida acceda a las funciones autorizadas según su rol.}
\CampoCU{Actor principal}{Usuario registrado.}
\CampoCU{Actores secundarios}{Servicio de autenticación institucional y módulo de autorización.}
\CampoCU{Precondiciones}{La cuenta existe, está activa y no se encuentra bloqueada.}
\CampoCU{Disparador}{El usuario selecciona la opción de iniciar sesión.}
\CampoCU{Flujo principal}{1) El sistema solicita correo institucional y contraseña. 2) El usuario ingresa los datos. 3) El sistema valida el formato y consulta la autenticación. 4) El servicio confirma las credenciales. 5) El sistema obtiene el rol vigente. 6) Se crea la sesión y se muestra el menú autorizado.}
\CampoCU{Flujos alternativos}{A1: Si las credenciales son inválidas, se rechaza el acceso. A2: Si la cuenta está inactiva o bloqueada, se informa el estado. A3: Si el servicio externo no responde, no se crea una sesión y se muestra indisponibilidad temporal.}
\CampoCU{Postcondiciones}{La sesión queda iniciada con el rol correspondiente o el intento fallido queda registrado.}
\CampoCU{Trazabilidad}{RF-01; HU-01; CA-01.1 y CA-01.2; \texttt{EV-CUE-001, EV-ENT-001, EV-ENT-005}.}
\end{tabularx}
\end{CUBox}

\begin{CUBox}{CU-02: Administrar roles y permisos}
\small
\begin{tabularx}{\linewidth}{@{}L{3.20cm}Y@{}}
\CampoCU{Objetivo}{Asignar, actualizar o retirar roles y permisos sin permitir que usuarios no autorizados modifiquen la matriz de acceso.}
\CampoCU{Actor principal}{Administrador del sistema.}
\CampoCU{Actores secundarios}{Módulo de autorización y registro de auditoría.}
\CampoCU{Precondiciones}{El administrador inició sesión y la cuenta objetivo existe.}
\CampoCU{Disparador}{El administrador abre la gestión de usuarios y selecciona una cuenta.}
\CampoCU{Flujo principal}{1) El sistema muestra los datos permitidos de la cuenta. 2) El administrador selecciona un rol válido. 3) El sistema presenta los permisos asociados. 4) El administrador confirma. 5) El sistema actualiza la autorización y registra el cambio.}
\CampoCU{Flujos alternativos}{A1: Si el rol no existe, se rechaza la operación. A2: Si el administrador intenta retirar su propio acceso crítico sin reemplazo autorizado, el sistema solicita una validación adicional. A3: Si la cuenta está eliminada o suspendida, no se aplican cambios.}
\CampoCU{Postcondiciones}{La matriz de permisos queda actualizada y el evento queda disponible para auditoría.}
\CampoCU{Trazabilidad}{RF-02; HU-02; CA-02.1 y CA-02.2; \texttt{EV-ENT-001, EV-ENT-006}.}
\end{tabularx}
\end{CUBox}

\begin{CUBox}{CU-03: Buscar y recomendar rutas}
\small
\begin{tabularx}{\linewidth}{@{}L{3.20cm}Y@{}}
\CampoCU{Objetivo}{Presentar rutas activas compatibles con los criterios del pasajero y explicar los factores principales de la recomendación.}
\CampoCU{Actor principal}{Pasajero universitario.}
\CampoCU{Actores secundarios}{Módulo de recomendación, servicio de mapas y catálogo de rutas.}
\CampoCU{Precondiciones}{El usuario está autenticado y existen rutas activas con cupos disponibles.}
\CampoCU{Disparador}{El pasajero registra o modifica los criterios de búsqueda.}
\CampoCU{Flujo principal}{1) El sistema solicita origen y destino generales, fecha e intervalo horario. 2) El usuario puede aplicar preferencias. 3) El sistema consulta rutas activas. 4) Se descartan rutas incompatibles o sin cupos. 5) El módulo calcula coincidencias. 6) El sistema ordena y muestra los resultados. 7) El usuario puede consultar la explicación de una recomendación.}
\CampoCU{Flujos alternativos}{A1: Si no existen coincidencias, se informa el resultado vacío y se permite cambiar filtros. A2: Si el módulo de recomendación no está disponible, se muestran resultados por criterios básicos. A3: Si faltan datos obligatorios, se solicita completarlos.}
\CampoCU{Postcondiciones}{Se presenta una lista de rutas sin crear reservas ni modificar publicaciones.}
\CampoCU{Trazabilidad}{RF-03, RF-14, RF-22 y RF-23; HU-03; CA-03.1 y CA-03.2; \texttt{EV-CUE-001, EV-ENT-003, EV-ENT-007, EV-ENT-010}.}
\end{tabularx}
\end{CUBox}

\begin{CUBox}{CU-04: Publicar una ruta}
\small
\begin{tabularx}{\linewidth}{@{}L{3.20cm}Y@{}}
\CampoCU{Objetivo}{Registrar una ruta válida y habilitarla para búsqueda y reserva.}
\CampoCU{Actor principal}{Conductor universitario.}
\CampoCU{Actores secundarios}{Servicio de mapas, catálogo de vehículos y módulo de notificaciones.}
\CampoCU{Precondiciones}{El conductor está autenticado, verificado y posee un vehículo habilitado.}
\CampoCU{Disparador}{El conductor selecciona la opción de publicar ruta.}
\CampoCU{Flujo principal}{1) El sistema solicita vehículo, origen y destino generales, fecha, hora, cupos y punto de encuentro. 2) El conductor completa los datos. 3) El sistema valida fecha, capacidad y campos obligatorios. 4) Se genera un resumen. 5) El conductor confirma. 6) El sistema crea la ruta en estado activo y asigna un identificador.}
\CampoCU{Flujos alternativos}{A1: Si la fecha es pasada o la capacidad es inválida, no se publica. A2: Si el vehículo no está habilitado, se dirige al registro de vehículo. A3: Si el servicio de mapas no responde, el sistema conserva los datos y solicita reintentar la validación geográfica.}
\CampoCU{Postcondiciones}{La ruta queda disponible mientras mantenga vigencia y cupos.}
\CampoCU{Trazabilidad}{RF-04, RF-15 y RF-19; HU-04; CA-04.1 y CA-04.2; \texttt{EV-ENT-002, EV-ENT-012, EV-ENT-015}.}
\end{tabularx}
\end{CUBox}

\begin{CUBox}{CU-05: Reservar un cupo}
\small
\begin{tabularx}{\linewidth}{@{}L{3.20cm}Y@{}}
\CampoCU{Objetivo}{Registrar una reserva sin superar la capacidad de la ruta y permitir una decisión informada.}
\CampoCU{Actor principal}{Pasajero universitario.}
\CampoCU{Actores secundarios}{Módulo de reservas, módulo de reputación y notificaciones.}
\CampoCU{Precondiciones}{La ruta está activa, posee cupos y el pasajero está autenticado.}
\CampoCU{Disparador}{El pasajero selecciona una ruta y solicita reservar.}
\CampoCU{Flujo principal}{1) El sistema muestra datos esenciales de la ruta y reputación del conductor. 2) El pasajero indica la cantidad permitida de cupos. 3) El sistema verifica disponibilidad y duplicados. 4) Se muestra un resumen. 5) El pasajero confirma. 6) La reserva se registra y los cupos se actualizan de forma atómica. 7) Se generan las notificaciones.}
\CampoCU{Flujos alternativos}{A1: Si no hay cupos, se rechaza la reserva. A2: Si existe una solicitud duplicada, se informa el registro vigente. A3: Si dos solicitudes compiten por el último cupo, solo una se confirma. A4: Si la ruta fue cancelada durante el proceso, no se registra la reserva.}
\CampoCU{Postcondiciones}{La reserva queda confirmada y la disponibilidad refleja la operación realizada.}
\CampoCU{Trazabilidad}{RF-05 y RF-24; HU-05; CA-05.1 y CA-05.2; \texttt{EV-CUE-001, EV-ENT-004, EV-ENT-011}.}
\end{tabularx}
\end{CUBox}

\begin{CUBox}{CU-06: Cancelar una reserva o ruta}
\small
\begin{tabularx}{\linewidth}{@{}L{3.20cm}Y@{}}
\CampoCU{Objetivo}{Cancelar una reserva o ruta vigente, liberar los cupos correspondientes y avisar a las personas afectadas.}
\CampoCU{Actor principal}{Pasajero universitario o conductor universitario.}
\CampoCU{Actores secundarios}{Módulo de reservas, módulo de rutas y notificaciones.}
\CampoCU{Precondiciones}{El elemento existe, está activo y pertenece al usuario que solicita la cancelación.}
\CampoCU{Disparador}{El usuario selecciona cancelar y confirma el motivo.}
\CampoCU{Flujo principal}{1) El sistema muestra el elemento activo. 2) El usuario solicita cancelar. 3) El sistema solicita motivo y confirmación. 4) Se valida la propiedad y el estado. 5) El sistema cambia el estado a cancelado. 6) Se liberan cupos cuando corresponde. 7) Se notifican los usuarios afectados.}
\CampoCU{Flujos alternativos}{A1: Si el elemento ya está cancelado o finalizado, no se modifica. A2: Si el usuario no es propietario, se rechaza la operación. A3: Si la notificación falla, la cancelación se mantiene y el intento queda pendiente de reenvío.}
\CampoCU{Postcondiciones}{La cancelación queda registrada, los cupos se actualizan y existe trazabilidad del motivo.}
\CampoCU{Trazabilidad}{RF-06, RF-12 y RF-21; \texttt{EV-ENT-008, EV-ENT-016}.}
\end{tabularx}
\end{CUBox}

\begin{CUBox}{CU-07: Gestionar notificaciones de viaje}
\small
\begin{tabularx}{\linewidth}{@{}L{3.20cm}Y@{}}
\CampoCU{Objetivo}{Generar y registrar avisos derivados de eventos relevantes sin producir mensajes duplicados.}
\CampoCU{Actor principal}{Sistema.}
\CampoCU{Actores secundarios}{Servicio de notificaciones, pasajero y conductor.}
\CampoCU{Precondiciones}{Existe un evento válido asociado con una ruta, reserva, viaje o incidencia.}
\CampoCU{Disparador}{Se registra una aceptación, rechazo, cambio, cancelación o recordatorio.}
\CampoCU{Flujo principal}{1) El sistema identifica el tipo de evento. 2) Determina destinatarios. 3) Construye un mensaje sin datos innecesarios. 4) Verifica que el evento no haya sido notificado. 5) Envía el aviso. 6) Registra marca temporal y estado de entrega.}
\CampoCU{Flujos alternativos}{A1: Si el evento ya fue procesado, no se duplica. A2: Si el canal no está disponible, se registra un intento pendiente. A3: Si el usuario desactivó un aviso opcional, el sistema respeta la preferencia.}
\CampoCU{Postcondiciones}{La notificación queda entregada, pendiente o fallida con trazabilidad del resultado.}
\CampoCU{Trazabilidad}{RF-12; HU-06; CA-06.1 y CA-06.2; \texttt{EV-ENT-006, EV-ENT-011, EV-ENT-012}.}
\end{tabularx}
\end{CUBox}

\begin{CUBox}{CU-08: Registrar un vehículo}
\small
\begin{tabularx}{\linewidth}{@{}L{3.20cm}Y@{}}
\CampoCU{Objetivo}{Asociar un vehículo válido con el perfil de un conductor verificado.}
\CampoCU{Actor principal}{Conductor universitario.}
\CampoCU{Actores secundarios}{Módulo de perfiles y catálogo de vehículos.}
\CampoCU{Precondiciones}{El conductor inició sesión y posee estado institucional verificado.}
\CampoCU{Disparador}{El conductor selecciona registrar vehículo.}
\CampoCU{Flujo principal}{1) El sistema solicita tipo, marca, modelo, color, capacidad y datos autorizados. 2) El conductor completa el formulario. 3) El sistema valida formatos, capacidad y duplicados. 4) Se presenta un resumen. 5) El conductor confirma. 6) El vehículo queda asociado al perfil.}
\CampoCU{Flujos alternativos}{A1: Si la capacidad es inválida, no se registra. A2: Si existe un registro duplicado, se informa al usuario. A3: Si el conductor no está verificado, se mantiene restringida la operación.}
\CampoCU{Postcondiciones}{El vehículo queda disponible para futuras publicaciones.}
\CampoCU{Trazabilidad}{RF-15; HU-07; CA-07.1 y CA-07.2; \texttt{EV-ENT-002, EV-ENT-015}.}
\end{tabularx}
\end{CUBox}

\begin{CUBox}{CU-09: Verificar identidad institucional}
\small
\begin{tabularx}{\linewidth}{@{}L{3.20cm}Y@{}}
\CampoCU{Objetivo}{Confirmar la pertenencia institucional antes de habilitar las funciones de movilidad colaborativa.}
\CampoCU{Actor principal}{Usuario en proceso de registro.}
\CampoCU{Actores secundarios}{Servicio de autenticación institucional UTEQ.}
\CampoCU{Precondiciones}{El usuario completó el registro y aceptó las condiciones aplicables.}
\CampoCU{Disparador}{El sistema envía una solicitud de verificación institucional.}
\CampoCU{Flujo principal}{1) El sistema prepara la solicitud mínima. 2) El servicio institucional valida la identidad. 3) Devuelve un estado. 4) El sistema registra la respuesta. 5) Si es válida, marca la cuenta como verificada y habilita las funciones correspondientes.}
\CampoCU{Flujos alternativos}{A1: Si la identidad es rechazada, la cuenta permanece restringida. A2: Si la respuesta queda pendiente, se informa el estado. A3: Si el servicio está indisponible, no se crea una verificación falsa y se permite reintentar.}
\CampoCU{Postcondiciones}{La cuenta queda verificada, rechazada o pendiente con una marca temporal.}
\CampoCU{Trazabilidad}{RF-16; HU-08; CA-08.1 y CA-08.2; \texttt{EV-CUE-001, EV-ENT-001, EV-ENT-006, EV-ENT-009}.}
\end{tabularx}
\end{CUBox}

\begin{CUBox}{CU-10: Consultar reputación del conductor}
\small
\begin{tabularx}{\linewidth}{@{}L{3.20cm}Y@{}}
\CampoCU{Objetivo}{Mostrar al pasajero un resumen de reputación antes de confirmar una reserva.}
\CampoCU{Actor principal}{Pasajero universitario.}
\CampoCU{Actores secundarios}{Módulo de reputación y catálogo de calificaciones.}
\CampoCU{Precondiciones}{La ruta está visible y el conductor posee un perfil consultable.}
\CampoCU{Disparador}{El pasajero abre el detalle de una ruta o el perfil resumido del conductor.}
\CampoCU{Flujo principal}{1) El sistema obtiene las calificaciones permitidas. 2) Calcula el promedio vigente y el número de viajes evaluados. 3) Aplica moderación y anonimización a los comentarios. 4) Muestra el resumen antes de la reserva.}
\CampoCU{Flujos alternativos}{A1: Si no existen evaluaciones, se muestra el estado ``sin calificaciones''. A2: Si una valoración fue retirada por moderación, no se incluye. A3: Si el perfil no está disponible, se informa sin bloquear la consulta de la ruta.}
\CampoCU{Postcondiciones}{La información se presenta sin revelar la identidad de quienes calificaron.}
\CampoCU{Trazabilidad}{RF-24 y RF-09; HU-09; CA-09.1 y CA-09.2; \texttt{EV-CUE-001, EV-ENT-003, EV-ENT-009, EV-ENT-014}.}
\end{tabularx}
\end{CUBox}

\subsection{Modelo conceptual de clases}

El modelo conceptual organiza las entidades del dominio, sus responsabilidades y relaciones principales. La Tabla~\ref{tab:clases-conceptuales} define las clases previstas para el diagrama de clases refinado.

\begin{longtable}{L{3.10cm} L{5.20cm} L{5.30cm}}
\caption{Clases conceptuales de RoutTech}\label{tab:clases-conceptuales}\\
\rowcolor{DocNavy}
\textcolor{white}{\textbf{Clase}} &
\textcolor{white}{\textbf{Atributos principales}} &
\textcolor{white}{\textbf{Operaciones principales}}\\
\endfirsthead
\rowcolor{DocNavy}
\textcolor{white}{\textbf{Clase}} &
\textcolor{white}{\textbf{Atributos principales}} &
\textcolor{white}{\textbf{Operaciones principales}}\\
\endhead
Usuario & idUsuario, correoInstitucional, estadoCuenta, fechaRegistro & autenticar(), actualizarPerfil(), consultarHistorial()\\
Rol & idRol, nombre, descripción & asignarPermiso(), retirarPermiso()\\
Permiso & idPermiso, código, descripción & validarAcceso()\\
PerfilUsuario & nombres, contactoPermitido, preferencias, reputaciónPromedio & actualizarDatos(), configurarPreferencias()\\
\texttt{Verificacion\allowbreak Institucional} & idVerificacion, estado, fechaSolicitud, fechaRespuesta & solicitar(), aprobar(), rechazar()\\
Vehiculo & idVehiculo, tipo, marca, modelo, color, capacidad, estado & registrar(), actualizar(), habilitar()\\
Ruta & idRuta, origenGeneral, destinoGeneral, fechaHora, cupos, estado & publicar(), actualizarCupos(), cancelar(), finalizar()\\
PuntoEncuentro & idPunto, nombreGeneral, referencia, estado & registrar(), modificar(), desactivar()\\
Reserva & idReserva, cantidadCupos, fechaSolicitud, estado & confirmar(), rechazar(), cancelar()\\
Notificacion & idNotificacion, tipo, mensaje, fechaHora, estadoEntrega & generar(), enviar(), marcarLeída()\\
Calificacion & idCalificacion, puntuación, comentarioModerado, fecha & registrar(), moderar()\\
Recomendacion & idRecomendacion, porcentajeCoincidencia, factores, fecha & calcular(), explicar(), descartar()\\
Incidencia & idIncidencia, categoría, descripción, estado, fecha & registrar(), asignar(), actualizarEstado(), cerrar()\\
\texttt{Reporte\allowbreak Administrativo} & idReporte, tipo, periodo, fechaGeneración & generar(), exportar()\\
RegistroAuditoria & idRegistro, actor, acción, fechaHora, resultado & registrarEvento(), consultar()\\
\end{longtable}

\subsubsection{Relaciones principales del modelo de clases}

\begin{itemize}
  \item Un \textbf{Usuario} posee un \textbf{PerfilUsuario} y puede tener uno o más \textbf{Rol}; cada rol agrupa permisos.
  \item Un usuario con rol de conductor puede registrar uno o más \textbf{Vehiculo}; cada \textbf{Ruta} se publica con un vehículo habilitado.
  \item Una \textbf{Ruta} define uno o más \textbf{PuntoEncuentro} y puede recibir varias \textbf{Reserva}; cada reserva pertenece a un pasajero y a una sola ruta.
  \item Las operaciones relevantes generan \textbf{Notificacion} y \textbf{RegistroAuditoria}.
  \item Una \textbf{Calificacion} se asocia con un viaje finalizado y contribuye al resumen de reputación del perfil evaluado.
  \item Una \textbf{Recomendacion} relaciona preferencias del usuario con rutas activas y conserva los factores usados para su explicación.
  \item Las \textbf{Incidencia} pueden relacionarse con una ruta, reserva o viaje y son gestionadas por usuarios con permisos administrativos.
  \item Los \textbf{ReporteAdministrativo} se construyen con información agregada y no exponen recorridos individuales identificables.
\end{itemize}


\subsection{Diagrama entidad--relación}

A partir del modelo conceptual se estructuró también una vista de datos orientada a persistencia. La Figura~\ref{fig:er-routtech} resume las entidades principales del dominio, sus atributos esenciales y las relaciones requeridas para soportar usuarios, rutas, reservas, recomendaciones, incidencias y notificaciones.

\FiguraDocumento{03_Modelado/Diagramas_UML/Diagrama entidad-relación.jpg}{0.96}{Diagrama entidad--relación de RoutTech.}{fig:er-routtech}

\subsection{Diagramas de secuencia}

Los diagramas de secuencia deberán representar la interacción temporal entre interfaz, servicios de aplicación, repositorios y servicios externos. La Tabla~\ref{tab:catalogo-secuencia} define el conjunto mínimo previsto.

\begin{longtable}{C{1.35cm} L{4.10cm} L{7.60cm}}
\caption{Catálogo de diagramas de secuencia}\label{tab:catalogo-secuencia}\\
\rowcolor{DocNavy}
\textcolor{white}{\textbf{ID}} &
\textcolor{white}{\textbf{Caso asociado}} &
\textcolor{white}{\textbf{Participantes principales}}\\
\endfirsthead
\rowcolor{DocNavy}
\textcolor{white}{\textbf{ID}} &
\textcolor{white}{\textbf{Caso asociado}} &
\textcolor{white}{\textbf{Participantes principales}}\\
\endhead
DS-01 & CU-01 Autenticación & Usuario, interfaz de acceso, servicio de autenticación, autorización y auditoría.\\
DS-02 & CU-02 Roles y permisos & Administrador, interfaz administrativa, servicio de usuarios, autorización y auditoría.\\
DS-03 & CU-03 Búsqueda y recomendación & Pasajero, interfaz de búsqueda, servicio de rutas, recomendación y mapas.\\
DS-04 & CU-04 Publicación de ruta & Conductor, interfaz de rutas, servicio de vehículos, rutas, mapas y persistencia.\\
DS-05 & CU-05 Reserva & Pasajero, interfaz de reserva, servicio de rutas, reservas, reputación y notificaciones.\\
DS-06 & CU-06 Cancelación & Usuario, servicio de rutas o reservas, notificaciones y auditoría.\\
DS-07 & CU-07 Notificación & Evento de negocio, servicio de notificaciones, canal externo y registro de entrega.\\
DS-08 & CU-08 Registro de vehículo & Conductor, interfaz de vehículo, servicio de perfiles, validación y repositorio.\\
DS-09 & CU-09 Verificación institucional & Usuario, servicio de verificación, servicio institucional UTEQ y autorización.\\
DS-10 & CU-10 Reputación & Pasajero, interfaz de detalle, servicio de reputación y repositorio de calificaciones.\\
\end{longtable}


Los diagramas incorporados cubren los escenarios transaccionales y de interacción más relevantes del sistema. En conjunto muestran la participación de la interfaz, la lógica de aplicación, la persistencia y los servicios externos en los flujos críticos del producto.

\FiguraDocumento{03_Modelado/Diagramas_UML/Secuencia iniciar sesión.jpg}{0.96}{Diagrama de secuencia DS-01: inicio de sesión institucional.}{fig:seq-inicio-sesion}

\FiguraDocumento{03_Modelado/Diagramas_UML/Secuencia buscar rutas.jpg}{0.96}{Diagrama de secuencia DS-03: búsqueda de rutas.}{fig:seq-buscar-rutas}

\FiguraDocumento{03_Modelado/Diagramas_UML/Secuencia publicar ruta.jpg}{0.96}{Diagrama de secuencia DS-04: publicación de una ruta.}{fig:seq-publicar-ruta}

\FiguraDocumento{03_Modelado/Diagramas_UML/Secuencia reservar cupo.jpg}{0.96}{Diagrama de secuencia DS-05: reserva de cupo.}{fig:seq-reservar-cupo}

\FiguraDocumento{03_Modelado/Diagramas_UML/Secuencia confirmar reserva.jpg}{0.96}{Diagrama de secuencia complementario: confirmación de reserva.}{fig:seq-confirmar-reserva}

\FiguraDocumento{03_Modelado/Diagramas_UML/Secuencia notificaciones.jpg}{0.96}{Diagrama de secuencia DS-07: gestión de notificaciones.}{fig:seq-notificaciones}

\FiguraDocumento{03_Modelado/Diagramas_UML/Secuencia reportar incidencia.jpg}{0.96}{Diagrama de secuencia asociado con el registro de incidencias.}{fig:seq-reportar-incidencia}

\FiguraDocumento{03_Modelado/Diagramas_UML/Secuencia recomendación explicable.jpg}{0.96}{Diagrama de secuencia asociado con la recomendación explicable de rutas.}{fig:seq-recomendacion-explicable}

\FiguraDocumento{03_Modelado/Diagramas_UML/Secuencia generar reporte.jpg}{0.96}{Diagrama de secuencia asociado con la generación de reportes administrativos.}{fig:seq-generar-reporte}

\FiguraDocumento{03_Modelado/Diagramas_UML/Secuencia calificar viaje.jpg}{0.96}{Diagrama de secuencia asociado con la calificación de viajes y reputación.}{fig:seq-calificar-viaje}

\subsection{Diagramas de actividad}

Los diagramas de actividad describirán decisiones, validaciones, concurrencia y rutas alternativas de los procesos principales. Se establecen los siguientes artefactos:

\begin{longtable}{C{1.35cm} L{4.40cm} L{7.30cm}}
\caption{Catálogo de diagramas de actividad}\label{tab:catalogo-actividad}\\
\rowcolor{DocNavy}
\textcolor{white}{\textbf{ID}} &
\textcolor{white}{\textbf{Proceso}} &
\textcolor{white}{\textbf{Decisiones representadas}}\\
\endfirsthead
\rowcolor{DocNavy}
\textcolor{white}{\textbf{ID}} &
\textcolor{white}{\textbf{Proceso}} &
\textcolor{white}{\textbf{Decisiones representadas}}\\
\endhead
DA-01 & Registro, verificación y autenticación & Cuenta existente, identidad válida, servicio disponible, rol autorizado y acceso aceptado o rechazado.\\
DA-02 & Publicación de ruta & Conductor verificado, vehículo habilitado, datos completos, fecha futura y capacidad válida.\\
DA-03 & Búsqueda y recomendación & Criterios válidos, existencia de rutas, disponibilidad de recomendación y consulta de explicación.\\
DA-04 & Reserva de cupo & Ruta activa, cupos disponibles, solicitud duplicada, concurrencia y confirmación.\\
DA-05 & Cancelación y recuperación & Propiedad del elemento, estado cancelable, liberación de cupos, aviso y búsqueda alternativa.\\
DA-06 & Calificación y reputación & Viaje finalizado, participante válido, evaluación previa, moderación y actualización del promedio.\\
\end{longtable}


Las actividades seleccionadas permiten verificar reglas de decisión, rutas alternativas y puntos de validación que deben mantenerse consistentes con los casos de uso y con los criterios de aceptación definidos en la Sección 3.

\FiguraDocumento{03_Modelado/Diagramas_UML/Actividad iniciar sesión.jpg}{0.96}{Diagrama de actividad DA-01: inicio de sesión y validación de acceso.}{fig:act-inicio-sesion}

\FiguraDocumento{03_Modelado/Diagramas_UML/Actividad publicar una ruta.jpg}{0.96}{Diagrama de actividad DA-02: publicación de una ruta.}{fig:act-publicar-ruta}

\FiguraDocumento{03_Modelado/Diagramas_UML/Actividad buscar y recomendar rutas.jpg}{0.96}{Diagrama de actividad DA-03: búsqueda y recomendación de rutas.}{fig:act-buscar-rutas}

\FiguraDocumento{03_Modelado/Diagramas_UML/Actividad reservar cupo.jpg}{0.96}{Diagrama de actividad DA-04: reserva de cupo.}{fig:act-reservar-cupo}

\FiguraDocumento{03_Modelado/Diagramas_UML/Actividad cancelar reserva.jpg}{0.96}{Diagrama de actividad DA-05: cancelación de reserva o ruta.}{fig:act-cancelar-reserva}

\FiguraDocumento{03_Modelado/Diagramas_UML/Actividad reportar incidencia.jpg}{0.96}{Diagrama de actividad asociado con el reporte de incidencias.}{fig:act-reportar-incidencia}

\subsection{Diagramas de estados}

Las entidades con ciclos de vida no triviales se modelarán mediante diagramas de estados. La Tabla~\ref{tab:catalogo-estados} define las transiciones mínimas que deben conservarse.

\begin{longtable}{C{1.35cm} L{3.15cm} L{8.60cm}}
\caption{Catálogo de diagramas de estados}\label{tab:catalogo-estados}\\
\rowcolor{DocNavy}
\textcolor{white}{\textbf{ID}} &
\textcolor{white}{\textbf{Entidad}} &
\textcolor{white}{\textbf{Estados y transiciones principales}}\\
\endfirsthead
\rowcolor{DocNavy}
\textcolor{white}{\textbf{ID}} &
\textcolor{white}{\textbf{Entidad}} &
\textcolor{white}{\textbf{Estados y transiciones principales}}\\
\endhead
DE-01 & Usuario & Registrado, pendiente de verificación, verificado, rechazado, suspendido, bloqueado y desactivado.\\
DE-02 & Ruta & Borrador, activa, completa, modificada, en curso, finalizada, cancelada y vencida.\\
DE-03 & Reserva & Solicitada, confirmada, rechazada, cancelada, completada y no presentada.\\
DE-04 & Incidencia & Registrada, clasificada, asignada, en atención, resuelta, cerrada y reabierta.\\
DE-05 & Notificación & Generada, pendiente, enviada, entregada, leída y fallida.\\
\end{longtable}

\subsection{Diagrama de componentes}

El modelo de componentes separará las responsabilidades del sistema para mantener independencia entre la interfaz, la lógica de negocio, la persistencia y las integraciones externas. Los componentes previstos son:

\begin{itemize}
  \item \textbf{Cliente web adaptable:} interfaces para pasajero, conductor y administrador.
  \item \textbf{API de aplicación:} punto de entrada para solicitudes autenticadas y coordinación de servicios.
  \item \textbf{Componente de identidad y autorización:} autenticación, verificación, roles y permisos.
  \item \textbf{Componente de usuarios y perfiles:} datos editables, preferencias, vehículos y reputación.
  \item \textbf{Componente de rutas:} publicación, búsqueda, puntos de encuentro, vigencia y cupos.
  \item \textbf{Componente de reservas:} solicitudes, confirmaciones, cancelaciones y control de concurrencia.
  \item \textbf{Componente de recomendación:} coincidencia de rutas y explicación de factores.
  \item \textbf{Componente de notificaciones:} generación, envío y seguimiento de avisos.
  \item \textbf{Componente de incidencias y reportes:} seguimiento administrativo e indicadores agregados.
  \item \textbf{Persistencia:} repositorios de datos, auditoría y respaldo.
  \item \textbf{Adaptadores externos:} autenticación UTEQ, mapas y canal de notificaciones.
\end{itemize}

Las dependencias deberán orientarse desde la interfaz hacia la API y desde los servicios de aplicación hacia contratos de persistencia o adaptadores externos. El componente de recomendación no podrá confirmar reservas ni modificar rutas sin una acción explícita de la persona usuaria.


La Figura~\ref{fig:componentes-routtech} muestra la partición funcional propuesta para el sistema. En ella se observa la separación entre cliente, lógica de aplicación, persistencia y servicios externos, así como las dependencias de comunicación necesarias para el funcionamiento coordinado del MVP y su evolución posterior.

\FiguraDocumento{03_Modelado/Diagramas_UML/Diagrama de componentes.jpg}{0.96}{Diagrama de componentes de RoutTech.}{fig:componentes-routtech}


\subsection{Diagrama de paquetes}

Como complemento estructural, la Figura~\ref{fig:paquetes-routtech} agrupa las responsabilidades del sistema en paquetes lógicos, facilitando la organización del modelo UML y la futura correspondencia con módulos de implementación.

\FiguraDocumento{03_Modelado/Diagramas_UML/Diagrama de paquetes.jpg}{0.96}{Diagrama de paquetes de RoutTech.}{fig:paquetes-routtech}

\subsection{Diagrama de despliegue}

El despliegue lógico de RoutTech contempla nodos separados para cliente, aplicación, datos y servicios externos. Esta representación no afirma que la infraestructura definitiva ya esté implementada; define la distribución esperada para el MVP y una posible evolución posterior.

\begin{longtable}{L{3.30cm} L{4.10cm} L{5.70cm}}
\caption{Nodos del despliegue lógico de RoutTech}\label{tab:nodos-despliegue}\\
\rowcolor{DocNavy}
\textcolor{white}{\textbf{Nodo}} &
\textcolor{white}{\textbf{Artefactos desplegados}} &
\textcolor{white}{\textbf{Comunicación principal}}\\
\endfirsthead
\rowcolor{DocNavy}
\textcolor{white}{\textbf{Nodo}} &
\textcolor{white}{\textbf{Artefactos desplegados}} &
\textcolor{white}{\textbf{Comunicación principal}}\\
\endhead
Dispositivo del usuario & Navegador web, interfaz adaptable y almacenamiento temporal de sesión & HTTPS con la API de aplicación.\\
Servidor web / aplicación & API, servicios de negocio, autorización, recomendación y notificaciones & HTTPS con clientes; conexión segura con base de datos y servicios externos.\\
Servidor de base de datos & Esquema relacional, auditoría, respaldos y datos sintéticos del MVP & Conexión restringida desde la aplicación.\\
Servicio institucional UTEQ & Verificación de identidad y estado institucional & Solicitudes autenticadas desde el adaptador de identidad.\\
Servicio de mapas & Geocodificación general, distancias y visualización de rutas & API externa mediante claves protegidas.\\
Servicio de notificaciones & Entrega de correo o avisos compatibles con el MVP & Solicitudes desde el componente de notificaciones y respuesta de estado.\\
Repositorio GitHub & Código, documentación, modelos y artefactos públicos anonimizados & Git mediante HTTPS o SSH; no almacena evidencia personal sin cifrar.\\
\end{longtable}

\subsection{Trazabilidad y control de los modelos}

Cada artefacto de modelado conservará un identificador estable y deberá relacionarse con los elementos de la matriz extendida. Los archivos editables se almacenarán en \nolinkurl{03_Modelado/Diagramas_UML/}, acompañados por exportaciones PNG y SVG cuando la herramienta lo permita. La trazabilidad mínima será la siguiente:

\begin{itemize}
  \item Caso de uso: \texttt{CU-XX} $\rightarrow$ RF, HU, CA y evidencia.
  \item Diagrama de secuencia: \texttt{DS-XX} $\rightarrow$ CU y componentes participantes.
  \item Diagrama de actividad: \texttt{DA-XX} $\rightarrow$ CU y criterios de aceptación.
  \item Diagrama de estados: \texttt{DE-XX} $\rightarrow$ entidad, RF y reglas de transición.
  \item Clase o componente: identificador del modelo $\rightarrow$ RF/RNF y caso de uso soportado.
  \item Nodo de despliegue: nodo $\rightarrow$ componente desplegado y RNF técnico relacionado.
\end{itemize}

La consistencia se verificará comparando nombres, actores, estados y relaciones entre los diagramas y el catálogo de requisitos. Un cambio aprobado en un RF Must have deberá revisarse en el caso de uso, los diagramas correspondientes, la historia, los criterios de aceptación y la matriz de trazabilidad.



% =====================================================
% 5. PRIORIZACION Y TRAZABILIDAD EXTENDIDA
% =====================================================
\section{Priorización y trazabilidad extendida}

La priorización de RoutTech combina tres perspectivas complementarias: MoSCoW para definir el compromiso de alcance, el modelo de Kano para interpretar el efecto de cada función sobre la satisfacción y el método WSJF para comparar valor y esfuerzo. La estimación se considera una línea base de planificación y deberá confirmarse durante los walkthrough con stakeholders; no constituye un resultado empírico definitivo. La matriz extendida mantiene la trazabilidad exigida desde la obligación legal y el objetivo del proyecto hasta el componente y el mockup correspondiente \cite{pohl2010,uteq2026}.

\subsection{Objetivos utilizados en la trazabilidad}

Para evitar descripciones extensas dentro de la matriz, se establecen los objetivos identificados en la Tabla~\ref{tab:objetivos-trazabilidad}. Estos códigos se conservarán también en el archivo CSV de trazabilidad ubicado en la carpeta \texttt{04\_Trazabilidad}.

\begin{table}[H]
\centering
\caption{Objetivos de RoutTech utilizados en la matriz de trazabilidad}
\label{tab:objetivos-trazabilidad}
\small
\rowcolors{2}{DocLight}{white}
\begin{tabularx}{\textwidth}{C{1.55cm} Y}
\rowcolor{DocNavy}
\textcolor{white}{\textbf{ID}} & \textcolor{white}{\textbf{Objetivo}}\\
OBJ-01 & Controlar el acceso mediante identidad institucional, roles y permisos verificables.\\
OBJ-02 & Coordinar la publicación, búsqueda, reserva y actualización de viajes compartidos.\\
OBJ-03 & Incrementar la confianza, seguridad percibida y capacidad de reacción de los participantes.\\
OBJ-04 & Apoyar decisiones mediante recomendaciones comprensibles e información agregada.\\
OBJ-05 & Proteger datos personales y respetar las obligaciones aplicables de la LOPDP.\\
OBJ-06 & Asegurar calidad técnica, verificabilidad, mantenibilidad y despliegue reproducible.\\
\end{tabularx}
\end{table}

\subsection{Método de priorización combinado}

En MoSCoW, \emph{Must have} identifica funciones indispensables para que el MVP sea utilizable; \emph{Should have} reúne capacidades importantes que pueden incorporarse después de estabilizar el núcleo; \emph{Could have} agrupa mejoras valiosas pero aplazables; y \emph{Won't have} registra funciones excluidas de la versión evaluada. En RoutTech se excluyen del alcance actual los pagos, el transporte abierto al público externo, el seguimiento permanente mediante coordenadas exactas y la sustitución de servicios de emergencia.

Kano clasifica cada RF como básico, de desempeño o atractivo. Las funciones básicas provocan insatisfacción cuando faltan; las de desempeño incrementan la utilidad conforme mejoran; y las atractivas añaden valor diferencial sin impedir el uso esencial cuando se aplazan.

El valor WSJF se calculó mediante la expresión:

\begin{equation}
\mathrm{WSJF}=\frac{\text{valor para el usuario}+\text{criticidad temporal}+\text{reducción de riesgo u oportunidad}}{\text{tamaño del trabajo}}.
\label{eq:wsjf}
\end{equation}

Las tres dimensiones del numerador se estimaron en una escala de 1 a 10. El tamaño del trabajo utiliza la serie 1, 2, 3, 5, 8 y 13. Un valor WSJF mayor indica una relación valor--esfuerzo más favorable, pero no reemplaza las obligaciones legales, de seguridad ni de explicabilidad.

\begin{landscape}
\subsection{Resultado de la priorización de requisitos funcionales}

La Tabla~\ref{tab:priorizacion-combinada} presenta la línea base combinada. El RF-22 conserva prioridad \emph{Could have} para una vista ampliada de explicación; sin embargo, la explicación mínima integrada en la recomendación es obligatoria por RNF-15 y se representa dentro de MU-03.
\begingroup
\scriptsize
\setlength{\tabcolsep}{3.5pt}
\renewcommand{\arraystretch}{0.98}
\begin{longtable}{C{1.25cm} C{1.65cm} L{2.55cm} C{1.10cm} C{1.10cm} C{1.10cm} C{1.10cm} C{1.30cm} C{1.15cm}}
\caption{Priorización combinada MoSCoW, Kano y WSJF}\label{tab:priorizacion-combinada}\\
\rowcolor{DocNavy}
\textcolor{white}{\textbf{RF}} &
\textcolor{white}{\textbf{MoSCoW}} &
\textcolor{white}{\textbf{Kano}} &
\textcolor{white}{\textbf{VU}} &
\textcolor{white}{\textbf{CT}} &
\textcolor{white}{\textbf{RR/O}} &
\textcolor{white}{\textbf{TW}} &
\textcolor{white}{\textbf{WSJF}} &
\textcolor{white}{\textbf{Orden}}\\
\endfirsthead
\rowcolor{DocNavy}
\textcolor{white}{\textbf{RF}} &
\textcolor{white}{\textbf{MoSCoW}} &
\textcolor{white}{\textbf{Kano}} &
\textcolor{white}{\textbf{VU}} &
\textcolor{white}{\textbf{CT}} &
\textcolor{white}{\textbf{RR/O}} &
\textcolor{white}{\textbf{TW}} &
\textcolor{white}{\textbf{WSJF}} &
\textcolor{white}{\textbf{Orden}}\\
\endhead
RF-01 & Must & Básico & 10 & 10 & 10 & 5 & 6.00 & 4\\
RF-02 & Must & Básico & 9 & 9 & 10 & 5 & 5.60 & 9\\
RF-03 & Must & Desempeño & 10 & 9 & 8 & 8 & 3.38 & 19\\
RF-04 & Must & Básico & 10 & 10 & 9 & 5 & 5.80 & 7\\
RF-05 & Must & Básico & 10 & 10 & 10 & 5 & 6.00 & 5\\
RF-06 & Should & Desempeño & 8 & 8 & 8 & 5 & 4.80 & 11\\
RF-07 & Should & Desempeño & 7 & 6 & 7 & 5 & 4.00 & 16\\
RF-08 & Should & Desempeño & 6 & 5 & 5 & 3 & 5.33 & 10\\
RF-09 & Should & Desempeño & 7 & 6 & 7 & 5 & 4.00 & 17\\
RF-10 & Could & Atractivo & 6 & 3 & 4 & 8 & 1.62 & 25\\
RF-11 & Could & Desempeño & 5 & 4 & 5 & 8 & 1.75 & 23\\
RF-12 & Must & Básico & 9 & 10 & 10 & 5 & 5.80 & 8\\
RF-13 & Should & Desempeño & 7 & 8 & 9 & 8 & 3.00 & 20\\
RF-14 & Should & Desempeño & 8 & 7 & 6 & 5 & 4.20 & 15\\
RF-15 & Must & Básico & 9 & 9 & 9 & 3 & 9.00 & 1\\
RF-16 & Must & Básico & 10 & 10 & 10 & 5 & 6.00 & 6\\
RF-17 & Should & Desempeño & 7 & 8 & 9 & 8 & 3.00 & 21\\
RF-18 & Should & Atractivo & 7 & 7 & 8 & 5 & 4.40 & 14\\
RF-19 & Should & Desempeño & 8 & 8 & 8 & 5 & 4.80 & 12\\
RF-20 & Should & Atractivo & 6 & 5 & 6 & 8 & 2.12 & 22\\
RF-21 & Should & Atractivo & 8 & 8 & 8 & 5 & 4.80 & 13\\
RF-22 & Could & Atractivo & 7 & 5 & 8 & 5 & 4.00 & 18\\
RF-23 & Should & Desempeño & 7 & 6 & 6 & 3 & 6.33 & 3\\
RF-24 & Must & Básico & 9 & 9 & 9 & 3 & 9.00 & 2\\
RF-25 & Could & Desempeño & 5 & 4 & 5 & 8 & 1.75 & 24\\
\end{longtable}
\endgroup
\end{landscape}

\textbf{Leyenda:} VU = valor para el usuario; CT = criticidad temporal; RR/O = reducción de riesgo u oportunidad; TW = tamaño del trabajo. El orden WSJF orienta la secuencia de implementación dentro de cada compromiso de alcance.

La Figura~\ref{fig:distribucion-priorizacion-rf} se reservará para visualizar la distribución de los 25 RF por categoría MoSCoW y Kano. La imagen deberá generarse desde \texttt{priorizacion\_moscow\_kano.csv}, evitando modificar manualmente los valores presentados en la tabla.

\FiguraDocumento{imagenes/FIG-05_Distribucion_Priorizacion_RF.jpg}{0.92}{Distribución de los requisitos funcionales según MoSCoW y Kano}{fig:distribucion-priorizacion-rf}

\subsection{Componentes identificados para la matriz}

Para compactar la trazabilidad se utilizan los códigos siguientes: CMP-01 identidad y autorización; CMP-02 usuarios y perfiles; CMP-03 rutas; CMP-04 reservas; CMP-05 recomendación; CMP-06 notificaciones y comunicación; CMP-07 incidencias y reputación; CMP-08 reportes e indicadores; CMP-09 persistencia y auditoría; CMP-10 cliente web adaptable; y CMP-11 adaptadores externos.

\subsection{Matriz de trazabilidad extendida}

La matriz contiene 48 filas: 25 correspondientes a RF, 15 a RNF y 8 a restricciones de diseño. De este modo supera el mínimo de 40 filas y conserva las columnas establecidas por la guía. Los códigos \texttt{EV-VAL-XXX} identifican evidencia de validación aún pendiente y deberán reemplazarse por códigos reales antes del corte. La Figura~\ref{fig:cadena-trazabilidad-extendida} mostrará posteriormente la cadena lógica de relaciones que se aplica en cada fila.

\FiguraDocumento{imagenes/FIG-06_Cadena_Trazabilidad_Extendida.jpg}{0.95}{Cadena de trazabilidad desde la base legal hasta el mockup}{fig:cadena-trazabilidad-extendida}

\begin{landscape}
\begingroup
\tiny
\setlength{\tabcolsep}{2.1pt}
\renewcommand{\arraystretch}{1.16}
\begin{longtable}{C{0.75cm} L{1.05cm} L{1.10cm} L{1.45cm} L{1.60cm} L{2.00cm} C{0.95cm} C{0.75cm} L{1.35cm} L{1.15cm} L{1.55cm} L{1.75cm} L{1.25cm}}
\caption{Matriz de trazabilidad extendida de RoutTech}\label{tab:matriz-trazabilidad-extendida}\\
\rowcolor{DocNavy}
\textcolor{white}{\textbf{ID}} &
\textcolor{white}{\textbf{Ley}} &
\textcolor{white}{\textbf{Artículo}} &
\textcolor{white}{\textbf{Objetivo}} &
\textcolor{white}{\textbf{Stakeholder}} &
\textcolor{white}{\textbf{ID-EV}} &
\textcolor{white}{\textbf{ID-RF}} &
\textcolor{white}{\textbf{Tipo}} &
\textcolor{white}{\textbf{ID-CU}} &
\textcolor{white}{\textbf{ID-HU}} &
\textcolor{white}{\textbf{ID-CA}} &
\textcolor{white}{\textbf{Componente}} &
\textcolor{white}{\textbf{Mockup}}\\
\endfirsthead
\rowcolor{DocNavy}
\textcolor{white}{\textbf{ID}} &
\textcolor{white}{\textbf{Ley}} &
\textcolor{white}{\textbf{Artículo}} &
\textcolor{white}{\textbf{Objetivo}} &
\textcolor{white}{\textbf{Stakeholder}} &
\textcolor{white}{\textbf{ID-EV}} &
\textcolor{white}{\textbf{ID-RF}} &
\textcolor{white}{\textbf{Tipo}} &
\textcolor{white}{\textbf{ID-CU}} &
\textcolor{white}{\textbf{ID-HU}} &
\textcolor{white}{\textbf{ID-CA}} &
\textcolor{white}{\textbf{Componente}} &
\textcolor{white}{\textbf{Mockup}}\\
\endhead
TR-001 & LOPDP & 7, 8 y 12 & OBJ-01; OBJ-05 & Usuario registrado & EV-CUE-001; EV-ENT-001; EV-ENT-005 & RF-01 & RF & CU-01 & HU-01 & CA-01.1; CA-01.2 & CMP-01; CMP-10 & MU-01\\
TR-002 & LOPDP & 10 literal j & OBJ-01; OBJ-05 & Administrador & EV-ENT-001; EV-ENT-006 & RF-02 & RF & CU-02 & HU-02 & CA-02.1; CA-02.2 & CMP-01; CMP-09; CMP-10 & MU-02\\
TR-003 & LOPDP & 10 y 12 & OBJ-02; OBJ-04; OBJ-05 & Pasajero universitario & EV-CUE-001; EV-ENT-003; EV-ENT-007; EV-ENT-010 & RF-03 & RF & CU-03 & HU-03 & CA-03.1; CA-03.2 & CMP-03; CMP-05; CMP-10 & MU-03\\
TR-004 & LOPDP & 10 & OBJ-02 & Conductor universitario & EV-ENT-002; EV-ENT-012; EV-ENT-015 & RF-04 & RF & CU-04 & HU-04 & CA-04.1; CA-04.2 & CMP-02; CMP-03; CMP-10 & MU-04\\
TR-005 & No aplica & -- & OBJ-02; OBJ-03 & Pasajero universitario & EV-CUE-001; EV-ENT-004; EV-ENT-011 & RF-05 & RF & CU-05 & HU-05 & CA-05.1; CA-05.2 & CMP-04; CMP-06; CMP-10 & MU-05\\
TR-006 & No aplica & -- & OBJ-02; OBJ-03 & Pasajero / conductor & EV-ENT-004; EV-ENT-011; EV-ENT-012 & RF-06 & RF & CU-06 & -- & -- & CMP-03; CMP-04; CMP-06 & --\\
TR-007 & LOPDP & 10 y 12 & OBJ-01; OBJ-05 & Usuario registrado & EV-CUE-001; EV-ENT-005 & RF-07 & RF & -- & -- & -- & CMP-02; CMP-10 & --\\
TR-008 & LOPDP & 10 & OBJ-03; OBJ-05 & Usuario registrado & EV-ENT-005; EV-ENT-009 & RF-08 & RF & -- & -- & -- & CMP-02; CMP-09; CMP-10 & --\\
TR-009 & LOPDP & 10 & OBJ-03 & Pasajero / conductor & EV-CUE-001; EV-ENT-003; EV-ENT-009 & RF-09 & RF & CU-10 & -- & -- & CMP-07; CMP-10 & MU-09\\
TR-010 & LOPDP & 10 & OBJ-04; OBJ-05 & Administrador / autoridades & EV-CUE-001; EV-DOC-003 & RF-10 & RF & -- & -- & -- & CMP-08; CMP-09 & --\\
TR-011 & LOPDP & 10 & OBJ-04; OBJ-05 & Administrador / autoridades & EV-DOC-003 & RF-11 & RF & -- & -- & -- & CMP-08; CMP-09 & --\\
TR-012 & LOPDP & 10 literal j & OBJ-02; OBJ-03 & Participante del viaje & EV-ENT-006; EV-ENT-011; EV-ENT-012 & RF-12 & RF & CU-07 & HU-06 & CA-06.1; CA-06.2 & CMP-06; CMP-10 & MU-06\\
TR-013 & LOPDP & 10 literal j & OBJ-03; OBJ-05 & Usuario / administrador & EV-ENT-006; EV-ENT-013 & RF-13 & RF & -- & -- & -- & CMP-07; CMP-09; CMP-10 & --\\
TR-014 & LOPDP & 10 & OBJ-02; OBJ-04 & Pasajero universitario & EV-CUE-001; EV-ENT-007; EV-ENT-010 & RF-14 & RF & CU-03 & -- & -- & CMP-03; CMP-05; CMP-10 & MU-03\\
TR-015 & LOPDP & 10 & OBJ-02; OBJ-03 & Conductor universitario & EV-ENT-002; EV-ENT-015 & RF-15 & RF & CU-08 & HU-07 & CA-07.1; CA-07.2 & CMP-02; CMP-10 & MU-07\\
TR-016 & LOPDP & 7, 8 y 12 & OBJ-01; OBJ-05 & Nuevo usuario & EV-CUE-001; EV-ENT-001; EV-ENT-006; EV-ENT-009 & RF-16 & RF & CU-09 & HU-08 & CA-08.1; CA-08.2 & CMP-01; CMP-11; CMP-10 & MU-08\\
TR-017 & LOPDP & 7, 8, 10 y 12 & OBJ-03; OBJ-05 & Participante del viaje & EV-ENT-003; EV-ENT-010; EV-ENT-014; EV-DOC-002 & RF-17 & RF & -- & -- & -- & CMP-03; CMP-11; CMP-10 & --\\
TR-018 & No aplica & -- & OBJ-03 & Pasajero / conductor & EV-ENT-009; EV-ENT-014 & RF-18 & RF & -- & -- & -- & CMP-03; CMP-06; CMP-10 & --\\
TR-019 & LOPDP & 7, 8 y 10 & OBJ-02; OBJ-03; OBJ-05 & Pasajero / conductor & EV-ENT-003; EV-ENT-010; EV-DOC-002 & RF-19 & RF & CU-04 & -- & -- & CMP-03; CMP-11; CMP-10 & MU-04\\
TR-020 & LOPDP & 10 literal j & OBJ-02; OBJ-03; OBJ-05 & Pasajero / conductor & EV-ENT-004; EV-ENT-011 & RF-20 & RF & -- & -- & -- & CMP-06; CMP-09; CMP-10 & --\\
TR-021 & No aplica & -- & OBJ-02; OBJ-03 & Pasajero universitario & EV-ENT-004; EV-ENT-011; EV-ENT-012 & RF-21 & RF & CU-06 & -- & -- & CMP-03; CMP-05; CMP-06 & --\\
TR-022 & LOPDP & 12 nums. 14 y 17 & OBJ-04; OBJ-05 & Pasajero universitario & EV-DOC-001; EV-DOC-002; EV-VAL-XXX & RF-22 & RF & CU-03 & HU-03 & -- & CMP-05; CMP-10 & MU-03\\
TR-023 & LOPDP & 10 y 12 & OBJ-02; OBJ-04; OBJ-05 & Pasajero universitario & EV-CUE-001; EV-ENT-007; EV-ENT-010 & RF-23 & RF & CU-03 & HU-03 & CA-03.1; CA-03.2 & CMP-02; CMP-05; CMP-10 & MU-03\\
TR-024 & LOPDP & 10 & OBJ-03; OBJ-05 & Pasajero universitario & EV-CUE-001; EV-ENT-003; EV-ENT-009; EV-ENT-014 & RF-24 & RF & CU-10 & HU-09 & CA-09.1; CA-09.2 & CMP-07; CMP-10 & MU-09\\
TR-025 & LOPDP & 10 & OBJ-04; OBJ-05 & Autoridades UTEQ / administrador & EV-DOC-003; EV-VAL-XXX & RF-25 & RF & -- & -- & -- & CMP-08; CMP-09; CMP-10 & --\\
TR-026 & No aplica & -- & OBJ-06 & Equipo del proyecto & EV-DOC-001; EV-DOC-004 & RNF-01 & RNF & CU-01; CU-03; CU-04; CU-05; CU-06; CU-10 & HU-01; HU-03; HU-04; HU-05; HU-09 & CA asociados & CMP-01 a CMP-10 & MU-01; MU-03; MU-04; MU-05; MU-09\\
TR-027 & No aplica & -- & OBJ-04; OBJ-06 & Pasajero universitario & EV-DOC-004; EV-CUE-001 & RNF-02 & RNF & CU-03 & HU-03 & CA-03.1; CA-03.2 & CMP-03; CMP-05 & MU-03\\
TR-028 & No aplica & -- & OBJ-02; OBJ-06 & Participante del viaje & EV-ENT-012; EV-DOC-004 & RNF-03 & RNF & CU-07 & HU-06 & CA-06.1; CA-06.2 & CMP-06 & MU-06\\
TR-029 & No aplica & -- & OBJ-06 & Todos los usuarios & EV-DOC-004 & RNF-04 & RNF & CU-01 a CU-10 & HU-01 a HU-09 & CA asociados & CMP-10 & MU-01 a MU-09\\
TR-030 & No aplica & -- & OBJ-06 & Usuarios técnicos y no técnicos & EV-CUE-001; EV-DOC-004; EV-VAL-XXX & RNF-05 & RNF & CU-01; CU-03; CU-04; CU-05; CU-08; CU-09; CU-10 & HU-01; HU-03; HU-04; HU-05; HU-07; HU-08; HU-09 & CA asociados & CMP-10 & MU-01; MU-03; MU-04; MU-05; MU-07; MU-08; MU-09\\
TR-031 & No aplica & -- & OBJ-06 & Todos los usuarios & EV-DOC-001; EV-DOC-004; EV-VAL-XXX & RNF-06 & RNF & CU-01 a CU-10 & HU-01 a HU-09 & CA asociados & CMP-10 & MU-01 a MU-09\\
TR-032 & No aplica & -- & OBJ-06 & Todos los usuarios & EV-DOC-004 & RNF-07 & RNF & CU-01 a CU-10 & -- & -- & CMP-01 a CMP-10 & --\\
TR-033 & LOPDP & 10 literal j & OBJ-02; OBJ-03; OBJ-05 & Pasajero / conductor & EV-DOC-004 & RNF-08 & RNF & CU-05; CU-06 & HU-05 & CA-05.1; CA-05.2 & CMP-04; CMP-09 & MU-05\\
TR-034 & LOPDP & 10 literal j & OBJ-01; OBJ-05 & Usuario / administrador & EV-CUE-001; EV-ENT-001; EV-ENT-006; EV-DOC-004 & RNF-09 & RNF & CU-01; CU-02 & HU-01; HU-02 & CA-01.1; CA-01.2; CA-02.1; CA-02.2 & CMP-01; CMP-09 & MU-01; MU-02\\
TR-035 & LOPDP & 10 literal j & OBJ-05; OBJ-06 & Todos los usuarios & EV-DOC-003; EV-DOC-004 & RNF-10 & RNF & CU-01; CU-09 & HU-01; HU-08 & CA-01.1; CA-08.1 & CMP-01; CMP-02; CMP-09 & MU-01; MU-08\\
TR-036 & LOPDP & 7, 8, 10 y 12 & OBJ-03; OBJ-05 & Participante del viaje & EV-ENT-003; EV-ENT-010; EV-ENT-014; EV-DOC-002 & RNF-11 & RNF & CU-03; CU-04 & -- & -- & CMP-03; CMP-11 & MU-03; MU-04\\
TR-037 & No aplica & -- & OBJ-06 & Equipo del proyecto & EV-DOC-004 & RNF-12 & RNF & -- & -- & -- & CMP-01 a CMP-09 & --\\
TR-038 & No aplica & -- & OBJ-06 & Equipo del proyecto & EV-DOC-001; EV-DOC-004 & RNF-13 & RNF & -- & -- & -- & CMP-01 a CMP-11 & --\\
TR-039 & No aplica & -- & OBJ-04; OBJ-06 & Equipo del proyecto / administrador & EV-DOC-004 & RNF-14 & RNF & CU-03 & HU-03 & -- & CMP-05 & MU-03\\
TR-040 & LOPDP & 12 nums. 14 y 17 & OBJ-04; OBJ-05 & Pasajero universitario & EV-DOC-001; EV-DOC-002; EV-VAL-XXX & RNF-15 & RNF & CU-03 & HU-03 & -- & CMP-05; CMP-10 & MU-03\\
TR-041 & No aplica & -- & OBJ-06 & Equipo del proyecto & EV-DOC-001 & RD-01 & RD & CU-01 a CU-10 & HU-01 a HU-09 & CA asociados & CMP-10 & MU-01 a MU-09\\
TR-042 & LOPDP & 7, 8 y 10 literal j & OBJ-01; OBJ-05 & Usuario / administrador & EV-ENT-001; EV-ENT-006 & RD-02 & RD & CU-01; CU-02; CU-09 & HU-01; HU-02; HU-08 & CA asociados & CMP-01; CMP-10 & MU-01; MU-02; MU-08\\
TR-043 & No aplica & -- & OBJ-02 & Pasajero / conductor & EV-DOC-001 & RD-03 & RD & -- & -- & -- & CMP-03; CMP-04 & --\\
TR-044 & LOPDP & 7, 8, 10 y 12 & OBJ-03; OBJ-05 & Participante del viaje & EV-DOC-002 & RD-04 & RD & CU-03; CU-04 & -- & -- & CMP-03; CMP-11 & MU-03; MU-04\\
TR-045 & LOPDP & 10 & OBJ-05; OBJ-06 & Equipo del proyecto & EV-DOC-002; EV-DOC-003 & RD-05 & RD & -- & -- & -- & CMP-09 & --\\
TR-046 & LOPDP & 7 y 8 & OBJ-05 & Participantes del estudio & EV-DOC-002 & RD-06 & RD & -- & -- & -- & -- & --\\
TR-047 & LOPDP & 12 nums. 14 y 17 & OBJ-04; OBJ-05 & Pasajero universitario & EV-DOC-001; EV-DOC-002 & RD-07 & RD & CU-03; CU-06 & HU-03 & CA-03.1 & CMP-05; CMP-10 & MU-03\\
TR-048 & No aplica & -- & OBJ-06 & Equipo del proyecto & EV-DOC-004 & RD-08 & RD & CU-01 a CU-10 & -- & -- & CMP-01 a CMP-09 & --\\
\end{longtable}
\endgroup
\end{landscape}

\subsection{Control de consistencia y actualización}

La matriz debe actualizarse cuando cambie un requisito, una obligación legal, una prioridad o un artefacto de diseño. La revisión mínima comprende los controles siguientes:

\begin{itemize}
  \item cada RF y RNF debe relacionarse con al menos un objetivo, stakeholder, evidencia y componente;
  \item todo RF Must have debe conservar CU, HU, CA y mockup, salvo que la interfaz se integre explícitamente en otro mockup trazado;
  \item las filas asociadas con la LOPDP deben utilizar la misma base legal definida en la Tabla~\ref{tab:requisitos-legales};
  \item ningún código \texttt{EV-VAL-XXX} puede permanecer en la versión entregada si la validación ya fue ejecutada;
  \item el archivo CSV debe contener la misma cantidad de filas y los mismos identificadores que la tabla del ERS;
  \item los cambios de prioridad deben reflejarse simultáneamente en el catálogo RF, las historias, el MVP y \texttt{priorizacion\_moscow\_kano.csv}.
\end{itemize}

La versión editable de los dos artefactos se almacenará en \texttt{04\_Trazabilidad/}. La tabla de priorización se exportará como \texttt{priorizacion\_moscow\_kano.csv} y la matriz como \texttt{matriz\_trazabilidad.csv}, permitiendo su revisión, filtrado y validación automática.

% =====================================================
% 6. PRODUCTO MINIMO VIABLE
% =====================================================
\section{Producto Mínimo Viable (MVP)}

\subsection{Propósito y condición de evaluación}

El Producto Mínimo Viable de RoutTech tiene como finalidad demostrar, mediante código ejecutable, que los flujos esenciales definidos en el ERS pueden integrarse en una solución coherente y verificable. El MVP no sustituye al producto completo ni representa una operación comercial de transporte; constituye una versión académica orientada a validar la coordinación básica entre pasajeros universitarios, conductores y administradores.

De acuerdo con la guía de la Entrega 3 (2A), el código del MVP debe mantenerse en un repositorio Git público separado, disponer de instrucciones de despliegue local, cubrir al menos el 60\,\% de los RF clasificados como \emph{Must have} y acompañarse de un video de demostración de duración menor o igual a tres minutos \cite{uteq2026}. El enlace exacto, la rama evaluada y el hash del commit deberán registrarse en la Tabla~\ref{tab:identificacion-mvp} cuando el repositorio quede congelado para el corte.

\begin{table}[H]
\centering
\caption{Identificación y control de versión del MVP de RoutTech}
\label{tab:identificacion-mvp}
\small
\rowcolors{2}{DocLight}{white}
\begin{tabularx}{\textwidth}{L{3.6cm} Y}
\rowcolor{DocNavy}
\textcolor{white}{\textbf{Elemento}} & \textcolor{white}{\textbf{Valor de control}}\\
Nombre del producto & RoutTech MVP -- Entrega 3 (2A).\\
Repositorio público separado & \texttt{[PENDIENTE: URL exacta del repositorio MVP]}.\\
Rama evaluada & \texttt{main}.\\
Commit evaluado & \texttt{[PENDIENTE: hash completo del commit de corte]}.\\
Licencia del código & MIT o Apache-2.0, según el archivo \texttt{LICENSE} del repositorio.\\
Instrucciones de despliegue & Archivo \texttt{README.md} del repositorio MVP.\\
Video de demostración & \nolinkurl{05_MVP/video_demo.mp4}, duración máxima de 3 minutos.\\
Responsables & Equipo CN: Córdova Carreño Mayra Lucila y Naranjo Flores Anderson Jeampiere.\\
\end{tabularx}
\end{table}

\subsection{Alcance funcional comprometido}

El catálogo de RoutTech contiene nueve RF de prioridad \emph{Must have}. Para la línea base del MVP se seleccionan siete de ellos, equivalentes al 77,8\,\% del conjunto prioritario. Esta cobertura supera el mínimo de 60\,\% exigido para la Entrega 3. La selección conserva un recorrido funcional completo: acceso, autorización, registro del vehículo, publicación de una ruta, recomendación, reserva y notificación.

\begin{table}[H]
\centering
\caption{Cobertura objetivo de requisitos Must have en el MVP}
\label{tab:cobertura-mvp-must}
\small
\rowcolors{2}{DocLight}{white}
\begin{tabularx}{\textwidth}{C{1.25cm} Y C{2.15cm} L{3.25cm}}
\rowcolor{DocNavy}
\textcolor{white}{\textbf{RF}} &
\textcolor{white}{\textbf{Funcionalidad}} &
\textcolor{white}{\textbf{MVP 2A}} &
\textcolor{white}{\textbf{Evidencia prevista}}\\
RF-01 & Autenticación de usuarios institucionales. & Incluido & Inicio de sesión válido e inválido en el video y pruebas funcionales.\\
RF-02 & Gestión de roles y permisos. & Incluido & Acceso diferenciado para pasajero, conductor y administrador.\\
RF-03 & Recomendación de rutas compatibles. & Incluido & Lista ordenada y explicación breve de los factores utilizados.\\
RF-04 & Publicación de rutas. & Incluido & Formulario validado y ruta visible para búsqueda.\\
RF-05 & Reserva de cupos. & Incluido & Reserva confirmada y actualización atómica de disponibilidad.\\
RF-12 & Notificaciones de eventos de viaje. & Incluido & Aviso interno por publicación, reserva o cambio de estado.\\
RF-15 & Registro de vehículo. & Incluido & Creación y consulta del vehículo habilitado para publicar rutas.\\
RF-16 & Verificación de identidad institucional. & Diferido & Se simulará únicamente la condición de cuenta verificada; la integración externa queda fuera del MVP.\\
RF-24 & Consulta del historial de reputación. & Diferido & Requiere viajes completados y calificaciones acumuladas; se mantiene para una iteración posterior.\\
\end{tabularx}
\end{table}

La cobertura se calcula mediante la Ecuación~\ref{eq:cobertura-mvp}:

\begin{equation}
\mathrm{Cobertura\ del\ MVP}=\frac{7}{9}\times 100=77{,}8.
\label{eq:cobertura-mvp}
\end{equation}

El porcentaje anterior expresa el alcance comprometido. Antes de la entrega deberá comprobarse contra el código del commit indicado en la Tabla~\ref{tab:identificacion-mvp}. Una función que aparezca únicamente en un mockup, una captura o una descripción, pero no pueda ejecutarse, no se contabilizará como cubierta.

\subsection{Recorrido funcional principal}

La demostración del MVP deberá seguir un recorrido reproducible que evidencie la integración entre los módulos seleccionados. El flujo principal será el siguiente:

\begin{enumerate}
  \item un usuario accede con credenciales institucionales de prueba y el sistema aplica el rol autorizado;
  \item un conductor registra los datos mínimos de su vehículo;
  \item el conductor publica una ruta con origen y destino generalizados, horario, punto de encuentro y cupos;
  \item un pasajero introduce sus criterios de búsqueda y recibe rutas compatibles ordenadas;
  \item el sistema muestra los factores básicos que justifican la recomendación presentada;
  \item el pasajero reserva un cupo y la disponibilidad se actualiza sin exceder la capacidad;
  \item el sistema registra una notificación para las partes involucradas;
  \item el administrador consulta el resultado con permisos limitados a la información necesaria.
\end{enumerate}

La Figura~\ref{fig:flujo-funcional-mvp} presenta el recorrido funcional principal previsto para el MVP. El diagrama deberá mantenerse alineado con el comportamiento observable en el video y con los RF de la Tabla~\ref{tab:cobertura-mvp-must}.

\FiguraDocumento
  {05_MVP/evidencias/FIG-07_Flujo_Funcional_MVP.jpg}
  {0.94}
  {Flujo funcional principal del MVP de RoutTech}
  {fig:flujo-funcional-mvp}

\subsection{Arquitectura mínima de implementación}

El MVP conservará una separación básica de responsabilidades para facilitar las pruebas y la evolución posterior. La interfaz web adaptable se comunicará con una API de aplicación; la API aplicará las reglas de identidad, rutas, reservas, recomendaciones y notificaciones; y la persistencia relacional almacenará datos sintéticos de usuarios, vehículos, rutas, cupos y reservas. Las integraciones externas de autenticación, mapas y mensajería se abstraerán mediante adaptadores, permitiendo reemplazarlas por servicios simulados durante la evaluación.

\begin{table}[H]
\centering
\caption{Módulos mínimos del MVP y trazabilidad asociada}
\label{tab:modulos-mvp}
\small
\rowcolors{2}{DocLight}{white}
\begin{tabularx}{\textwidth}{C{1.45cm} L{3.55cm} Y L{3.25cm}}
\rowcolor{DocNavy}
\textcolor{white}{\textbf{Código}} &
\textcolor{white}{\textbf{Módulo}} &
\textcolor{white}{\textbf{Responsabilidad en el MVP}} &
\textcolor{white}{\textbf{Trazabilidad}}\\
MVP-M01 & Identidad y autorización & Iniciar sesión, reconocer el rol y bloquear funciones no autorizadas. & RF-01, RF-02; CU-01, CU-02; MU-01, MU-02.\\
MVP-M02 & Perfiles y vehículos & Mantener la información mínima del conductor y su vehículo. & RF-15; CU-08; HU-07; MU-07.\\
MVP-M03 & Rutas & Crear, validar, consultar y mantener rutas activas con cupos. & RF-04; CU-04; HU-04; MU-04.\\
MVP-M04 & Recomendación & Filtrar y ordenar rutas mediante reglas deterministas y explicables. & RF-03; CU-03; HU-03; MU-03; RNF-15.\\
MVP-M05 & Reservas & Registrar solicitudes y actualizar la disponibilidad de manera consistente. & RF-05; CU-05; HU-05; MU-05; RNF-08.\\
MVP-M06 & Notificaciones & Mostrar avisos internos asociados con los cambios del viaje. & RF-12; CU-07; HU-06; MU-06.\\
MVP-M07 & Persistencia y auditoría básica & Conservar datos sintéticos, estados y eventos necesarios para la demostración. & CMP-09; RNF-10; RD-05.\\
\end{tabularx}
\end{table}

La Figura~\ref{fig:arquitectura-mvp} presenta la distribución mínima de estos módulos y sus conexiones. Antes del corte deberá contrastarse con la arquitectura realmente implementada y actualizarse si el repositorio evaluado difiere del esquema mostrado.

\FiguraDocumento
  {05_MVP/evidencias/FIG-08_Arquitectura_MVP_RoutTech.jpg}
  {0.94}
  {Arquitectura mínima de implementación del MVP de RoutTech}
  {fig:arquitectura-mvp}

\subsection{Criterios técnicos y no funcionales aplicables}

Aunque el alcance funcional del MVP sea reducido, no se excluyen las restricciones críticas de calidad. La versión evaluada deberá demostrar al menos los siguientes controles:

\begin{itemize}
  \item los flujos incluidos deben ejecutarse de extremo a extremo y conservar trazabilidad con RF, CU, HU, CA y mockup;
  \item la búsqueda y recomendación deberán respetar el objetivo cuantificado de RNF-02 en el entorno de prueba documentado;
  \item las reservas no deberán generar sobreasignación de cupos, conforme a RNF-08;
  \item los permisos deberán impedir que un rol acceda a funciones no autorizadas, conforme a RNF-09;
  \item las credenciales y datos personales no deberán almacenarse ni mostrarse en texto plano, conforme a RNF-10;
  \item la recomendación deberá incluir una explicación breve y comprensible, conforme a RNF-15;
  \item el despliegue deberá poder repetirse siguiendo el README, conforme a RNF-13.
\end{itemize}

Los valores medidos, tiempos de respuesta y resultados de pruebas no se incorporarán al documento hasta disponer de la ejecución real y de los registros correspondientes. Esta regla evita declarar resultados que no puedan reproducirse desde el repositorio.

\subsection{Despliegue y reproducción local}

El archivo \texttt{README.md} del repositorio MVP deberá permitir que el docente ejecute el sistema sin depender de instrucciones externas. Como mínimo, documentará:

\begin{enumerate}
  \item requisitos de software y versiones utilizadas;
  \item clonación del repositorio y selección del commit evaluado;
  \item creación del archivo de variables de entorno a partir de un ejemplo sin secretos;
  \item instalación de dependencias o ejecución mediante contenedores;
  \item inicialización de la base de datos con datos sintéticos;
  \item credenciales de prueba para los tres roles, sin reutilizar cuentas personales;
  \item comando de inicio, puerto de acceso y forma de detener el servicio;
  \item ejecución de las pruebas automatizadas y criterio para interpretar el resultado;
  \item relación exacta entre los RF cubiertos y los módulos del código.
\end{enumerate}

Cuando se utilice Docker, el procedimiento preferente será un comando equivalente a \texttt{docker compose up --build}. Si se utiliza una instalación local, las versiones del entorno deberán fijarse para reducir diferencias entre equipos.

\subsection{Pruebas y evidencia de funcionamiento}

La aceptación del MVP se apoyará en pruebas funcionales asociadas con los escenarios Gherkin de la Sección~\ref{sec:criterios-gherkin}, complementadas con pruebas negativas de autorización, validación de campos, disponibilidad de cupos y consistencia de datos. Cada ejecución deberá identificar el RF, el escenario, los datos utilizados, el resultado esperado, el resultado observado y el estado de aprobación.

La Tabla~\ref{tab:evidencias-mvp} define los artefactos mínimos que deberán quedar disponibles para la revisión.

\begin{table}[H]
\centering
\caption{Evidencias mínimas de funcionamiento del MVP}
\label{tab:evidencias-mvp}
\small
\rowcolors{2}{DocLight}{white}
\begin{tabularx}{\textwidth}{L{3.35cm} L{3.45cm} Y}
\rowcolor{DocNavy}
\textcolor{white}{\textbf{Artefacto}} &
\textcolor{white}{\textbf{Ubicación prevista}} &
\textcolor{white}{\textbf{Contenido verificable}}\\
Código fuente & Repositorio Git público separado & Módulos ejecutables, historial de commits y licencia.\\
README de despliegue & \texttt{05\_MVP/README.md} y repositorio MVP & Pasos reproducibles, dependencias, credenciales sintéticas y RF cubiertos.\\
Video de demostración & \nolinkurl{05_MVP/video_demo.mp4} & Recorrido funcional continuo de hasta tres minutos, sin datos personales reales.\\
Resultados de pruebas & Repositorio MVP, carpeta de pruebas & Casos ejecutados, resultado esperado, observado y evidencia de aprobación.\\
Capturas del recorrido & \nolinkurl{05_MVP/evidencias/} & Imágenes de los estados principales enlazadas a RF y CA.\\
Commit de corte & Tabla~\ref{tab:identificacion-mvp} & Hash completo que permite reconstruir exactamente la versión evaluada.\\
\end{tabularx}
\end{table}

La Figura~\ref{fig:evidencia-recorrido-mvp} presenta el mosaico disponible del recorrido funcional principal. Para considerarlo evidencia definitiva, las capturas deberán corresponder al mismo commit indicado en la Tabla~\ref{tab:identificacion-mvp}.

\FiguraDocumento
  {05_MVP/evidencias/FIG-09_Evidencia_Recorrido_MVP.jpg}
  {0.94}
  {Evidencia visual del recorrido funcional del MVP de RoutTech}
  {fig:evidencia-recorrido-mvp}

\subsection{Limitaciones y funciones diferidas}

El MVP utilizará datos sintéticos y no realizará pagos, seguimiento GPS permanente, atención de emergencias ni verificación directa contra sistemas institucionales sin autorización. RF-16 se representará mediante un estado de verificación simulado claramente identificado y RF-24 quedará fuera del recorrido principal hasta disponer de viajes completados y calificaciones consistentes. Estas decisiones no eliminan los requisitos del producto completo; únicamente delimitan la versión demostrable de la Entrega 3.

Las funciones diferidas deberán permanecer registradas en el backlog y en la priorización, conservando sus relaciones de trazabilidad. Cualquier sustitución de un RF seleccionado por otro RF Must deberá actualizar simultáneamente la Tabla~\ref{tab:cobertura-mvp-must}, el archivo de cobertura del repositorio, el video de demostración y la matriz de trazabilidad.

\subsection{Estado de cierre de la sección}

La estructura, el alcance y las evidencias esperadas del MVP quedan definidos en esta sección. Para cerrar el criterio de evaluación todavía deberán incorporarse los datos reales del repositorio separado, el hash del commit, la ejecución documentada, el video de demostración y las capturas correspondientes. Hasta que esos elementos existan, el porcentaje de la Ecuación~\ref{eq:cobertura-mvp} debe interpretarse como cobertura objetivo y no como resultado empírico verificado.




% =====================================================
% 7. COMPONENTE EMPIRICO Y PROTOCOLO EXPERIMENTAL
% =====================================================
\section{Componente empírico y protocolo experimental}\label{sec:componente-empirico}

\subsection{Propósito del componente empírico}

El componente empírico de RoutTech tiene como finalidad obtener evidencia primaria que permita validar la calidad y la utilidad de los requisitos de explicabilidad asociados con el módulo de recomendación de rutas. Esta sección no presenta resultados hipotéticos: define el protocolo que deberá registrarse antes de la ejecución, los participantes, los instrumentos, las variables, el procedimiento y el plan de análisis. Los resultados se incorporarán únicamente después de que existan datos verificables y scripts capaces de reproducir las tablas y figuras correspondientes.

La línea base del proyecto dispone de 16 entrevistas semiestructuradas y 43 respuestas del cuestionario de movilidad, registradas en la Sección~1. Estas fuentes sustentan la identificación inicial de necesidades, pero no reemplazan las dos rondas específicas de validación de explicabilidad previstas en este protocolo. El cumplimiento de los mínimos de la Entrega 3 se determinará mediante los archivos reales, sus fichas técnicas y sus identificadores de evidencia, no solo mediante las cantidades declaradas en el documento.

\subsection{Enfoque empírico seleccionado}

Para RoutTech se adopta el \textbf{Enfoque 3: elicitación y validación de requisitos de explicabilidad para el componente de IA}. La selección se justifica porque RF-03 incorpora la recomendación de rutas y RNF-15 exige que cada recomendación comunique los factores principales que la originan. La explicabilidad debe especificarse como una propiedad verificable desde la Ingeniería de Requisitos y evaluarse con usuarios afectados, en lugar de tratarse únicamente como una característica técnica del algoritmo \cite{chazette2020,chazette2021,chazette2022}. Esta decisión también responde al interés reciente por integrar componentes de IA en tareas de Ingeniería de Requisitos bajo procedimientos sistemáticos y verificables \cite{cheng2026,arora2024,vogelsang2024}.

El componente estudiado será el mecanismo que ordena rutas compatibles según origen y destino generalizados, coincidencia de horario, disponibilidad de cupos, proximidad del punto de encuentro y restricciones declaradas por la persona usuaria. Durante el estudio se utilizará un prototipo controlado con datos sintéticos. La explicación no revelará información personal de otros participantes ni mostrará coordenadas exactas.

\begin{table}[H]
\centering
\caption{Identificación del protocolo experimental de RoutTech}
\label{tab:identificacion-protocolo}
\small
\rowcolors{2}{DocLight}{white}
\begin{tabularx}{\textwidth}{L{4.15cm} Y}
\rowcolor{DocNavy}
\textcolor{white}{\textbf{Elemento}} &
\textcolor{white}{\textbf{Definición}}\\
Enfoque seleccionado & Enfoque 3: elicitación y validación de requisitos de explicabilidad.\\
Componente evaluado & Recomendación y ordenamiento de rutas compatibles de RoutTech.\\
Diseño & Estudio de caso único con dos rondas sucesivas de validación.\\
Unidad de análisis & Requisito de explicabilidad y percepción de la explicación por perfil de usuario.\\
Versión del protocolo & v1.0; deberá coincidir con el archivo registrado en OSF.\\
Estado de ejecución & Pendiente de registro previo y recolección; no se declaran resultados en esta versión.\\
Ubicación documental & \nolinkurl{06_Experimento/protocolo.pdf} y \nolinkurl{06_Experimento/osf_registration.pdf}.\\
\end{tabularx}
\end{table}

La Figura~\ref{fig:diseno-estudio-explicabilidad} reservará el espacio para el diseño general del estudio. La imagen definitiva deberá mostrar la relación entre evidencia inicial, prototipo de explicación, primera ronda, refinamiento de RNF, segunda ronda y análisis final.

\FiguraDocumento
  {06_Experimento/evidencias/FIG-10_Diseno_Estudio_Explicabilidad_RoutTech.jpg}
  {0.94}
  {Diseño del estudio de explicabilidad de RoutTech}
  {fig:diseno-estudio-explicabilidad}

\subsection{Preguntas de investigación y marco PICOC}

Las preguntas se formularon mediante PICOC para mantener una relación explícita entre población, intervención, comparación, resultados y contexto. El estudio se orienta a comprender qué explicaciones requieren los perfiles de RoutTech y cómo convertir esas expectativas en RNF cuantificados y verificables.

\begin{table}[H]
\centering
\caption{Aplicación del marco PICOC al estudio de RoutTech}
\label{tab:picoc-routtech}
\small
\rowcolors{2}{DocLight}{white}
\begin{tabularx}{\textwidth}{C{1.25cm} L{3.0cm} Y}
\rowcolor{DocNavy}
\textcolor{white}{\textbf{Código}} &
\textcolor{white}{\textbf{Elemento}} &
\textcolor{white}{\textbf{Aplicación en RoutTech}}\\
P & Población & Estudiantes o pasajeros, conductores universitarios y revisores técnicos adultos vinculados con el contexto UTEQ.\\
I & Intervención & Prototipo de explicación asociado con una recomendación de ruta generada a partir de datos sintéticos.\\
C & Comparación & Perfiles técnicos frente a no técnicos y primera ronda frente a segunda ronda de validación.\\
O & Resultados & Necesidad percibida, comprensión, satisfacción, cobertura del marco, acuerdo y requisitos refinados.\\
C & Contexto & Movilidad colaborativa universitaria entre sedes y sectores relacionados con la UTEQ.\\
\end{tabularx}
\end{table}

\textbf{RQ1.} ¿Qué información y qué formas de explicación consideran necesarias los usuarios técnicos y no técnicos para comprender una recomendación de ruta emitida por RoutTech?

\textbf{RQ2.} ¿En qué medida el catálogo refinado de requisitos de explicabilidad cubre las dimensiones aplicables del marco seleccionado y mantiene métricas verificables?

\textbf{RQ3.} ¿Qué nivel de comprensión y satisfacción percibida reportan los participantes después de revisar la explicación propuesta, y qué diferencias se observan entre perfiles y rondas?

\subsection{Proposiciones del estudio}

Debido a que el enfoque corresponde a un estudio de caso orientado a requisitos de explicabilidad, se plantean proposiciones de investigación y no afirmaciones de resultado:

\begin{itemize}
  \item \textbf{PR-01:} los perfiles técnicos y no técnicos requerirán explicaciones con distinto nivel de detalle, aunque compartirán la necesidad de conocer los factores principales de compatibilidad.
  \item \textbf{PR-02:} la segunda ronda de validación mostrará una mayor cobertura de dimensiones de explicabilidad que la primera, después del refinamiento de RNF y prototipos.
  \item \textbf{PR-03:} una explicación breve, visible antes de confirmar la reserva y vinculada con factores concretos permitirá alcanzar una valoración media de comprensión de al menos 4 sobre 5.
\end{itemize}

El umbral de PR-03 funciona como criterio de aceptación del diseño; no constituye un resultado hasta que se calcule con los datos reales.

\subsection{Participantes y estrategia de selección}

El protocolo contempla dos rondas de validación con un mínimo de seis participantes por ronda: tres personas con perfil técnico y tres con perfil no técnico. El perfil técnico podrá estar integrado por estudiantes avanzados, profesionales o docentes con conocimientos de software, datos o evaluación de sistemas. El perfil no técnico estará compuesto por estudiantes o conductores universitarios que representen a los usuarios finales y que no desempeñen una función técnica en el estudio.

\begin{table}[H]
\centering
\caption{Criterios de participación en el estudio de explicabilidad}
\label{tab:criterios-participantes}
\small
\rowcolors{2}{DocLight}{white}
\begin{tabularx}{\textwidth}{L{3.45cm} Y Y}
\rowcolor{DocNavy}
\textcolor{white}{\textbf{Perfil}} &
\textcolor{white}{\textbf{Criterios de inclusión}} &
\textcolor{white}{\textbf{Criterios de exclusión}}\\
Técnico & Ser mayor de edad; comprender conceptos básicos de sistemas o requisitos; aceptar el consentimiento actualizado; no haber diseñado directamente el prototipo evaluado. & Participar como autor del requisito evaluado; no completar el instrumento; revocar el consentimiento.\\
No técnico & Ser mayor de edad; pertenecer o estar vinculado con la comunidad universitaria; tener experiencia real o potencial como pasajero o conductor; aceptar el consentimiento actualizado. & No corresponder al perfil de uso; no completar el recorrido; revocar el consentimiento.\\
\end{tabularx}
\end{table}

La selección será intencional y documentada por código de participante. La muestra no se presentará como representativa de toda la población UTEQ. Cada intervención deberá quedar respaldada por consentimiento, acta, ficha técnica y evidencia audiovisual en la zona correspondiente. Los tamaños mínimos y la separación entre perfiles siguen el diseño indicado para la validación de explicabilidad y las recomendaciones de estudios de caso en Ingeniería de Software \cite{runeson2009,wohlin2012}.

\subsection{Variables e indicadores}

La Tabla~\ref{tab:variables-estudio} operacionaliza las variables que se medirán durante las dos rondas. Las escalas, fórmulas y fuentes deberán coincidir con la matriz editable almacenada en \nolinkurl{06_Experimento/matriz_operacionalizacion.csv}.

\begin{longtable}{L{2.45cm} L{3.0cm} L{4.45cm} L{4.15cm}}
\caption{Variables e indicadores del estudio empírico}\label{tab:variables-estudio}\\
\rowcolor{DocNavy}
\textcolor{white}{\textbf{Tipo}} &
\textcolor{white}{\textbf{Variable}} &
\textcolor{white}{\textbf{Indicador o medición}} &
\textcolor{white}{\textbf{Fuente}}\\
\endfirsthead
\rowcolor{DocNavy}
\textcolor{white}{\textbf{Tipo}} &
\textcolor{white}{\textbf{Variable}} &
\textcolor{white}{\textbf{Indicador o medición}} &
\textcolor{white}{\textbf{Fuente}}\\
\endhead
Independiente & Perfil de usuario & Técnico o no técnico, registrado mediante código anónimo. & Ficha de caracterización.\\
Independiente & Ronda de validación & Primera o segunda ronda. & Registro de sesión.\\
Independiente & Tipo de explicación & Factores, contraste, ejemplo o combinación definida en el prototipo. & Versión del mockup o prototipo.\\
Dependiente & Necesidad percibida & Selección y prioridad de elementos explicativos requeridos. & Cuestionario y entrevista breve.\\
Dependiente & Comprensión & Puntuación Likert de 1 a 5 y respuesta correcta a preguntas de comprobación. & Instrumento de evaluación.\\
Dependiente & Satisfacción & Media de ítems Likert de 1 a 5 sobre utilidad, claridad y oportunidad. & Instrumento de evaluación.\\
Dependiente & Cobertura del marco & Porcentaje de dimensiones aplicables cubiertas por al menos un RNF verificable. & Catálogo de RNF y matriz de cobertura.\\
Dependiente & Acuerdo & Coeficiente de acuerdo entre revisores para la clasificación de necesidades y cobertura. & Matriz de codificación.\\
Control & Escenario y datos & Mismo conjunto de rutas sintéticas, instrucciones y orden de tareas por ronda. & Guion del estudio.\\
Control & Entorno de evaluación & Dispositivo, navegador, versión del prototipo y tiempo máximo documentados. & Ficha técnica de sesión.\\
\end{longtable}

La cobertura del marco se calculará mediante la Ecuación~\ref{eq:cobertura-explicabilidad}:

\begin{equation}
C_{exp} = \frac{D_c}{D_a} \times 100,
\label{eq:cobertura-explicabilidad}
\end{equation}

donde $D_c$ representa las dimensiones aplicables cubiertas por al menos un RNF y $D_a$ el total de dimensiones aplicables del marco.

\subsection{Instrumentos de recolección y validación}

Los instrumentos se versionarán en \nolinkurl{06_Experimento/instrumentos/}. Cada archivo deberá indicar código, versión, fecha, responsable y relación con las RQ.

\begin{table}[H]
\centering
\caption{Instrumentos previstos para el componente empírico}
\label{tab:instrumentos-experimento}
\small
\rowcolors{2}{DocLight}{white}
\begin{tabularx}{\textwidth}{C{1.35cm} L{3.35cm} Y L{3.45cm}}
\rowcolor{DocNavy}
\textcolor{white}{\textbf{ID}} &
\textcolor{white}{\textbf{Instrumento}} &
\textcolor{white}{\textbf{Finalidad}} &
\textcolor{white}{\textbf{Archivo previsto}}\\
INS-01 & Guion de validación v2.0 & Presentar el caso, orientar las tareas y formular preguntas sin inducir respuestas. & \nolinkurl{guion_validacion_v2.pdf}.\\
INS-02 & Ficha de caracterización & Registrar perfil, experiencia y rol mediante código anónimo. & \nolinkurl{ficha_caracterizacion.csv}.\\
INS-03 & Prototipo de explicación & Mostrar la recomendación, sus factores y la acción de ampliar detalles. & \nolinkurl{prototipo_explicacion_v1.png}.\\
INS-04 & Cuestionario de evaluación & Medir necesidad, comprensión, claridad, utilidad y satisfacción en escala Likert. & \nolinkurl{cuestionario_explicabilidad_v2.pdf}.\\
INS-05 & Rúbrica de cobertura & Verificar que cada RNF sea necesario, claro, cuantificado y trazable a una dimensión. & \nolinkurl{rubrica_cobertura_RNF.xlsx}.\\
INS-06 & Acta de walkthrough & Registrar requisitos revisados, observaciones, decisiones y firmas enmascaradas. & \nolinkurl{acta_walkthrough.pdf}.\\
INS-07 & Consentimiento actualizado & Autorizar participación, grabación, anonimización y uso científico conforme a la LOPDP. & \nolinkurl{consentimiento_v2.pdf}.\\
\end{tabularx}
\end{table}

Los ítems del cuestionario deberán evitar conceptos dobles, términos vagos y preguntas que sugieran la respuesta, siguiendo criterios de calidad para encuestas en Ingeniería de Software \cite{molleri2020}. Antes de la primera ronda se realizará una prueba piloto de comprensión del instrumento; cualquier cambio posterior se registrará en el historial de versiones y, cuando afecte el protocolo, en la adenda ética y en el registro OSF.

\subsection{Procedimiento de ejecución}

El procedimiento deberá aplicarse en el orden indicado para conservar comparabilidad entre sesiones:

\begin{enumerate}
  \item finalizar el protocolo, la adenda ética y los instrumentos v2.0;
  \item registrar previamente el protocolo en OSF antes de recolectar datos del estudio \cite{nosek2018};
  \item reclutar participantes conforme a los criterios de inclusión y asignar un código anónimo;
  \item obtener el consentimiento informado actualizado antes de grabar o aplicar instrumentos;
  \item presentar el escenario de movilidad y permitir que la persona examine la recomendación;
  \item solicitar que explique con sus palabras por qué el sistema recomendó esa ruta;
  \item aplicar el cuestionario de comprensión, satisfacción y necesidad de información;
  \item realizar una entrevista breve o walkthrough para recoger observaciones cualitativas;
  \item codificar la primera ronda, revisar los RNF y actualizar el prototipo sin alterar los datos originales;
  \item ejecutar la segunda ronda con el mismo protocolo base y la versión refinada;
  \item analizar los datos con scripts versionados y vincular cada hallazgo con su evidencia;
  \item redactar resultados solo después de verificar archivos, cálculos y trazabilidad.
\end{enumerate}

La Figura~\ref{fig:procedimiento-validacion-explicabilidad} reservará el espacio del procedimiento completo. La versión definitiva deberá mostrar los puntos de control ético, la separación de rondas y el ciclo de refinamiento.

\FiguraDocumento
  {06_Experimento/evidencias/FIG-11_Procedimiento_Validacion_Explicabilidad.jpg}
  {0.94}
  {Procedimiento de validación de requisitos de explicabilidad}
  {fig:procedimiento-validacion-explicabilidad}

\subsection{Plan de análisis de datos}

El análisis integrará evidencia cuantitativa y cualitativa y se reportará con información suficiente para su reproducción \cite{jedlitschka2008}. Los scripts se almacenarán en \nolinkurl{06_Experimento/scripts_analisis/} y deberán generar directamente las tablas y figuras incluidas en la versión de resultados.

\begin{itemize}
  \item \textbf{Análisis descriptivo:} frecuencias, porcentajes, mediana, media y dispersión para los ítems de necesidad, comprensión y satisfacción.
  \item \textbf{Consistencia de escalas:} revisión de respuestas incompletas y cálculo de consistencia interna cuando el instrumento contenga varios ítems para el mismo constructo.
  \item \textbf{Comparación entre perfiles:} prueba de normalidad Shapiro--Wilk; prueba t para grupos independientes cuando se cumplan los supuestos o prueba U de Mann--Whitney cuando no se cumplan.
  \item \textbf{Comparación entre rondas:} prueba apropiada según si las personas son independientes o repetidas; la decisión se registrará antes del análisis.
  \item \textbf{Tamaño del efecto:} Cohen $d$ para comparaciones paramétricas o una medida no paramétrica equivalente.
  \item \textbf{Acuerdo entre revisores:} coeficiente $\kappa$ de Cohen para dos revisores y $\kappa$ de Fleiss cuando intervengan tres o más \cite{cohen1960,fleiss1971}.
  \item \textbf{Análisis cualitativo:} codificación temática de comentarios, agrupación por tipo de explicación, perfil y ronda, con vínculo al fragmento anonimizado.
  \item \textbf{Saturación:} registro acumulado de códigos nuevos por sesión para identificar cuándo dejan de emerger categorías relevantes \cite{hennink2017,guest2006}.
\end{itemize}

La Figura~\ref{fig:curva-saturacion-explicabilidad} queda reservada para la curva de saturación real. No deberá insertarse una curva simulada; la imagen se generará desde la tabla de codificación mediante el script correspondiente.

\FiguraDocumento
  {06_Experimento/resultados/FIG-12_Curva_Saturacion_Tematica.jpg}
  {0.92}
  {Curva de saturación temática del estudio de explicabilidad}
  {fig:curva-saturacion-explicabilidad}

\textit{Nota de control:} los valores representados en la Figura~\ref{fig:curva-saturacion-explicabilidad} deberán sustituirse o confirmarse con la tabla real de codificación temática. La figura no debe presentarse como resultado empírico definitivo mientras no exista el archivo de datos y el script reproducible que la sustenten.

\subsection{Criterios de aceptación del estudio}

El componente empírico se considerará ejecutado únicamente cuando se cumplan simultáneamente los siguientes criterios:

\begin{itemize}
  \item el protocolo y los instrumentos utilizados coinciden con la versión registrada previamente en OSF o sus cambios están justificados;
  \item se completan dos rondas con el mínimo de perfiles definido y consentimientos válidos;
  \item el 100\% de las sesiones declaradas tiene evidencia verificable, ficha técnica y hash;
  \item los datos crudos, datos procesados y scripts permiten reproducir las tablas y figuras;
  \item cada requisito refinado conserva trazabilidad hacia observaciones, códigos y participantes anonimizados;
  \item los resultados de comprensión, satisfacción, cobertura y acuerdo se reportan con su método de cálculo;
  \item las limitaciones y desacuerdos se documentan sin eliminar respuestas desfavorables.
\end{itemize}

\subsection{Ética, privacidad y gestión de datos}

La participación será voluntaria y estará condicionada a un consentimiento informado actualizado que incluya grabación, anonimización, uso académico y eventual depósito de datos anonimizados. El tratamiento se sujetará a la Ley Orgánica de Protección de Datos Personales del Ecuador \cite{lopdp2021}. La persona participante podrá retirarse antes de la anonimización definitiva sin que esto produzca consecuencias académicas o institucionales.

Los consentimientos originales, videos, audios y actas con firmas visibles se almacenarán dentro de \nolinkurl{02_Evidencias/00_Restringido/} en un contenedor cifrado con AES-256. La zona pública contendrá exclusivamente transcripciones anonimizadas, copias enmascaradas, fotografías sin rostros identificables, respuestas sin columnas personales y fichas técnicas con hash SHA-256. La contraseña del contenedor no se incorporará al repositorio.

El conjunto de datos abierto incluirá únicamente información anonimizada y se preparará conforme a los principios FAIR \cite{wilkinson2016}. Los datos crudos restringidos se conservarán durante el período definido en el plan ético y posteriormente se eliminarán mediante el procedimiento documentado. Ningún archivo restringido se distribuirá bajo la licencia del dataset.

\subsection{Registro previo y reproducibilidad}

Antes de la primera sesión del estudio, el equipo deberá registrar en OSF una versión congelada del protocolo con RQ, proposiciones, variables, participantes, instrumentos y plan de análisis. El comprobante se almacenará en \nolinkurl{06_Experimento/osf_registration.pdf}. La URL persistente se añadirá a la Tabla~\ref{tab:identificacion-protocolo}, al README del repositorio y al manuscrito paralelo.

La carpeta experimental deberá conservar, como mínimo, la estructura indicada en la Tabla~\ref{tab:estructura-experimento}.

\begin{table}[H]
\centering
\caption{Estructura de archivos del componente experimental}
\label{tab:estructura-experimento}
\small
\rowcolors{2}{DocLight}{white}
\begin{tabularx}{\textwidth}{L{4.2cm} Y}
\rowcolor{DocNavy}
\textcolor{white}{\textbf{Ruta}} &
\textcolor{white}{\textbf{Contenido verificable}}\\
\nolinkurl{06_Experimento/protocolo.pdf} & Protocolo completo con RQ, proposiciones, variables, muestra, procedimiento y análisis.\\
\nolinkurl{06_Experimento/osf_registration.pdf} & Comprobante de registro previo con URL y fecha anterior a la ejecución.\\
\nolinkurl{06_Experimento/instrumentos/} & Guion, ficha, cuestionario, rúbrica, actas y consentimiento utilizados.\\
\nolinkurl{06_Experimento/resultados/} & Datos crudos anonimizados, datos procesados, tablas y figuras reproducibles.\\
\nolinkurl{06_Experimento/scripts_analisis/} & Scripts en Python o R y archivo de dependencias.\\
\nolinkurl{06_Experimento/evidencias/} & Diagramas del protocolo y material público no identificable.\\
\end{tabularx}
\end{table}

\subsection{Amenazas a la validez previstas}

Antes de ejecutar el estudio se reconocen las siguientes amenazas y medidas de mitigación, siguiendo la clasificación habitual de investigación empírica en Ingeniería de Software \cite{wohlin2012,runeson2009}:

\begin{table}[H]
\centering
\caption{Amenazas a la validez y mitigaciones previstas}
\label{tab:amenazas-validez}
\small
\rowcolors{2}{DocLight}{white}
\begin{tabularx}{\textwidth}{L{2.65cm} Y Y}
\rowcolor{DocNavy}
\textcolor{white}{\textbf{Categoría}} &
\textcolor{white}{\textbf{Amenaza}} &
\textcolor{white}{\textbf{Mitigación prevista}}\\
Interna & Influencia del moderador, aprendizaje entre tareas o diferencias en las instrucciones. & Guion único, orden controlado, moderador capacitado y registro de desviaciones.\\
Externa & Muestra intencional y localizada en un único contexto universitario. & Describir perfiles y contexto; evitar generalizar los resultados a todas las universidades.\\
Constructo & Confundir satisfacción con comprensión real de la recomendación. & Combinar escala Likert con preguntas de comprobación y observación del recorrido.\\
Conclusión & Tamaño pequeño, supuestos estadísticos no cumplidos o múltiples comparaciones. & Reportar tamaño del efecto, usar pruebas acordes a los datos y declarar incertidumbre.\\
\end{tabularx}
\end{table}

\subsection{Estado de cumplimiento y evidencias pendientes}

La Tabla~\ref{tab:estado-componente-empirico} distingue entre la evidencia de campo ya declarada y los elementos que todavía deben verificarse para cerrar el protocolo. Esta distinción evita considerar automáticamente las 16 entrevistas y las 43 respuestas como cumplimiento de todas las exigencias de la segunda ronda.

\begin{table}[H]
\centering
\caption{Estado documental del componente empírico}
\label{tab:estado-componente-empirico}
\small
\rowcolors{2}{DocLight}{white}
\begin{tabularx}{\textwidth}{L{4.2cm} L{3.0cm} Y}
\rowcolor{DocNavy}
\textcolor{white}{\textbf{Elemento}} &
\textcolor{white}{\textbf{Estado actual}} &
\textcolor{white}{\textbf{Acción para cierre}}\\
16 entrevistas iniciales & Declaradas en la línea base & Verificar archivos, duración, consentimiento, transcripción, ficha y hash.\\
43 respuestas del cuestionario & Declaradas en la línea base & Conservar exportación anonimizada y fotografías de aplicación verificables.\\
Adenda ética de segunda ronda & Pendiente de comprobación & Incorporar versión aprobada anterior a la nueva recolección.\\
Registro previo OSF & Pendiente & Registrar protocolo y guardar comprobante antes de ejecutar el estudio.\\
Dos rondas de explicabilidad & Pendiente & Ejecutar mínimo seis participantes por ronda y documentar perfiles.\\
Scripts y resultados & Pendiente & Procesar únicamente datos reales y reproducir cada tabla y figura.\\
\end{tabularx}
\end{table}

\subsection{Cierre de la sección}

El protocolo define un estudio reproducible para convertir las necesidades de explicación del módulo de recomendación en RNF verificables y validados con perfiles técnicos y no técnicos. La sección queda metodológicamente estructurada, pero su cierre empírico dependerá del registro previo, de la ejecución real de las dos rondas, de la evidencia ética y audiovisual y de los resultados reproducidos mediante scripts. Hasta entonces, cualquier métrica indicada debe interpretarse como objetivo o criterio de aceptación, no como hallazgo del proyecto.



% =====================================================
% 8. PREPARACION DE PUBLICACION Y DATOS FAIR
% =====================================================
\section{Preparación de publicación y conjunto de datos FAIR}

\subsection{Propósito y alcance de la preparación editorial}

Esta sección define los artefactos que permitirán transformar la documentación y la evidencia verificable de RoutTech en un paquete de publicación reproducible. La preparación editorial no implica afirmar resultados que todavía no han sido obtenidos; su propósito es establecer la estructura del conjunto de datos, del manuscrito, de los scripts, de la selección de revista y de la declaración de disponibilidad de materiales. Los resultados y la discusión solo se completarán después de ejecutar el protocolo experimental descrito en la Sección~\ref{sec:componente-empirico}, procesar los datos reales y verificar que las tablas y figuras puedan reproducirse desde los scripts versionados.

El paquete se organizará en \nolinkurl{07_Publicacion/} y se vinculará con las evidencias anonimizadas, la matriz de trazabilidad, el protocolo registrado en OSF y el código del MVP. La preparación seguirá los principios FAIR de localización, accesibilidad, interoperabilidad y reutilización de datos \cite{wilkinson2016}, así como los principios de citación de software y artefactos digitales \cite{smith2016,druskat2021}.

\begin{table}[H]
\centering
\caption{Artefactos de preparación para publicación}
\label{tab:artefactos-publicacion}
\small
\rowcolors{2}{DocLight}{white}
\begin{tabularx}{\textwidth}{L{4.25cm} Y L{2.55cm}}
\rowcolor{DocNavy}
\textcolor{white}{\textbf{Ruta}} &
\textcolor{white}{\textbf{Contenido esperado}} &
\textcolor{white}{\textbf{Estado}}\\
\nolinkurl{07_Publicacion/manuscrito_borrador/} & Fuente LaTeX, PDF preliminar y archivos de la plantilla de la revista objetivo. & En preparación\\
\nolinkurl{07_Publicacion/analisis_revistas.md} & Comparación de candidatas con APC y por suscripción o modalidad híbrida. & Plantilla creada\\
\nolinkurl{07_Publicacion/dataset_zenodo/} & Dataset anonimizado, diccionario, scripts, licencia e instrucciones de citación. & Estructura creada\\
\nolinkurl{06_Experimento/scripts_analisis/} & Código que genera todas las tablas, estadísticas y figuras del manuscrito. & Pendiente de ejecución\\
\nolinkurl{06_Experimento/osf_registration.pdf} & Comprobante del registro previo y URL persistente del protocolo. & Pendiente de verificación\\
\nolinkurl{CITATION.cff} & Metadatos normalizados para citar el repositorio y sus autores. & Pendiente de completar\\
\end{tabularx}
\end{table}

\subsection{Paquete de datos abierto conforme a FAIR}

El depósito abierto contendrá únicamente datos anonimizados y materiales de replicación que puedan compartirse sin revelar identidad, voz, rostro, firma, cédula, correo, teléfono, dirección o coordenadas exactas. Los archivos originales restringidos no formarán parte de Zenodo ni de OSF y permanecerán bajo el régimen ético y de retención definido para el proyecto. Esta separación es coherente con la Ley Orgánica de Protección de Datos Personales del Ecuador \cite{lopdp2021}.

La estructura propuesta del paquete se muestra en la Tabla~\ref{tab:estructura-dataset-zenodo}. Los formatos se seleccionan para favorecer la lectura por personas y máquinas, evitando que la reutilización dependa de una aplicación propietaria.

\begin{table}[H]
\centering
\caption{Estructura propuesta del paquete para Zenodo}
\label{tab:estructura-dataset-zenodo}
\small
\rowcolors{2}{DocLight}{white}
\begin{tabularx}{\textwidth}{L{4.25cm} L{2.0cm} Y}
\rowcolor{DocNavy}
\textcolor{white}{\textbf{Archivo o carpeta}} &
\textcolor{white}{\textbf{Formato}} &
\textcolor{white}{\textbf{Contenido y control}}\\
\nolinkurl{README_dataset.md} & Markdown & Descripción, alcance, licencia, estructura, requisitos de uso y forma de citación.\\
\nolinkurl{data/transcripciones/} & JSON & Transcripciones anonimizadas con código de participante, turno, fragmento e ID de evidencia.\\
\nolinkurl{data/cuestionario/} & CSV & Respuestas sin columnas identificativas y con diccionario de variables.\\
\nolinkurl{data/requisitos/} & JSON & Corpus de RF, RNF y RD con prioridad, fuente, evidencia y estado de validación.\\
\nolinkurl{data/trazabilidad/} & CSV & Matriz extendida desde Ley, objetivo y stakeholder hasta mockup y componente.\\
\nolinkurl{data/codificacion/} & CSV & Fragmentos anonimizados, códigos, categorías, requisito derivado y analista.\\
\nolinkurl{scripts/} & Python o R & Procesamiento, análisis estadístico y generación de tablas y figuras.\\
\nolinkurl{outputs/} & CSV, PNG, SVG & Resultados reproducidos automáticamente; no se editarán manualmente.\\
\nolinkurl{LICENSE_DATASET.txt} & Texto & Licencia CC BY 4.0 y exclusión expresa de la zona restringida.\\
\nolinkurl{CITATION.cff} & YAML & Metadatos de citación del dataset y autores.\\
\end{tabularx}
\end{table}

La Figura~\ref{fig:estructura-paquete-fair} reservará el esquema visual del paquete FAIR. La imagen final deberá mostrar la relación entre fuentes anonimizadas, scripts, salidas reproducibles, metadatos y depósito con DOI.

\FiguraDocumento
  {07_Publicacion/figuras/FIG-13_Estructura_Paquete_FAIR.jpg}
  {0.94}
  {Estructura del paquete de datos FAIR de RoutTech}
  {fig:estructura-paquete-fair}

\subsection{Diccionario de datos y metadatos}

Cada archivo tabular deberá disponer de un diccionario que describa el nombre de la variable, definición, tipo de dato, valores permitidos, unidad, regla de anonimización, valores perdidos y procedencia. Las columnas identificativas se eliminarán antes de publicar y los códigos de participante no permitirán reconstruir la identidad sin la tabla de correspondencia restringida.

\begin{table}[H]
\centering
\caption{Campos mínimos del diccionario de datos}
\label{tab:diccionario-datos}
\small
\rowcolors{2}{DocLight}{white}
\begin{tabularx}{\textwidth}{L{3.2cm} Y Y}
\rowcolor{DocNavy}
\textcolor{white}{\textbf{Campo}} &
\textcolor{white}{\textbf{Descripción}} &
\textcolor{white}{\textbf{Ejemplo de aplicación}}\\
Nombre de variable & Identificador ASCII, sin espacios ni acentos. & \nolinkurl{comprension_explicacion}\\
Definición & Significado operativo y relación con la RQ o el requisito. & Valoración de comprensión de la explicación recibida.\\
Tipo y escala & Texto, entero, decimal, booleano, nominal, ordinal o razón. & Entero ordinal de 1 a 5.\\
Valores permitidos & Dominio válido y codificación de categorías. & 1 = muy baja; 5 = muy alta.\\
Valor perdido & Código utilizado cuando no existe una respuesta válida. & Vacío o \nolinkurl{NA}; nunca cero por defecto.\\
Anonimización & Transformación aplicada para prevenir identificación. & Eliminación de correo y sustitución por código.\\
Procedencia & Archivo, técnica e ID de evidencia de origen. & \nolinkurl{EV-CUE-001}.\\
\end{tabularx}
\end{table}

Los metadatos del depósito incluirán título, autores, afiliación, versión, fecha, descripción, palabras clave, licencia, financiador cuando corresponda, relación con el repositorio GitHub y relación con el registro OSF. Una vez creado el depósito, el DOI se incorporará al README, al archivo \nolinkurl{CITATION.cff}, al manuscrito y a la portada o sección de disponibilidad del ERS.

\subsection{Anonimización, licencia y retención}

Antes de publicar cualquier conjunto se realizará una revisión en dos niveles: el analista que prepara el archivo y un segundo integrante que verifica que no existan identificadores directos o combinaciones que permitan identificar razonablemente a una persona. La revisión abarcará contenido textual, nombres de archivos, propiedades del documento, metadatos EXIF y comentarios internos de hojas de cálculo.

\begin{itemize}
  \item \textbf{Material abierto:} transcripciones anonimizadas, respuestas depuradas, corpus de requisitos, matriz de trazabilidad, scripts, figuras, diccionarios y documentación metodológica.
  \item \textbf{Material excluido:} consentimientos originales, videos, audios, actas con firmas visibles, tabla de correspondencia de participantes y documentos originales de la organización.
  \item \textbf{Licencia del dataset y documentos:} Creative Commons Atribución 4.0 Internacional (CC BY 4.0).
  \item \textbf{Licencia del código:} MIT o Apache-2.0, conforme a la decisión declarada en el repositorio del MVP.
  \item \textbf{Retención restringida:} conservación durante el plazo definido en el plan de gestión de datos y eliminación segura mediante acta al finalizar dicho período.
\end{itemize}

La publicación bajo una licencia abierta no modifica las restricciones éticas: solo se licenciarán los archivos que hayan superado el proceso de anonimización y que estén incluidos explícitamente en el inventario del depósito.

\subsection{Manuscrito paralelo}

El manuscrito paralelo documentará el estudio de requisitos de explicabilidad aplicado a RoutTech. Como título preliminar se propone: \textit{Requisitos de explicabilidad para un sistema universitario de movilidad colaborativa: un estudio empírico}. Este título podrá ajustarse cuando exista una revista objetivo y se conozcan los resultados, sin alterar la pregunta de investigación registrada.

La Tabla~\ref{tab:estructura-manuscrito} muestra las partes del manuscrito y la condición necesaria para cerrarlas. Los apartados de resultados, discusión y conclusiones permanecerán marcados como pendientes hasta que existan datos reales y salidas reproducidas.

\begin{table}[H]
\centering
\caption{Estructura y estado del manuscrito paralelo}
\label{tab:estructura-manuscrito}
\small
\rowcolors{2}{DocLight}{white}
\begin{tabularx}{\textwidth}{L{3.2cm} Y L{3.0cm}}
\rowcolor{DocNavy}
\textcolor{white}{\textbf{Sección}} &
\textcolor{white}{\textbf{Contenido verificable}} &
\textcolor{white}{\textbf{Condición de cierre}}\\
Título y metadatos & Título, autoría, filiación, correos institucionales y palabras clave. & Puede prepararse ahora\\
Resumen preliminar & Contexto, objetivo y método, sin inventar resultados ni conclusiones. & Antes de seleccionar revista\\
Introducción & Problema concreto, brecha, contribuciones previstas y RQ. & Con referencias verificadas\\
Trabajo relacionado & Búsqueda abreviada, criterios y tabla comparativa de estudios cercanos. & Búsqueda documentada\\
Metodología & PICOC, participantes, instrumentos, procedimiento, ética y análisis. & Protocolo OSF registrado\\
Resultados & Tablas y figuras generadas desde scripts, organizadas por RQ. & Después del experimento\\
Discusión & Interpretación vinculada a datos y literatura; implicaciones y hallazgos. & Después de resultados\\
Amenazas a la validez & Amenazas internas, externas, de constructo y de conclusión. & Con mitigaciones reales\\
Conclusiones & Respuesta por RQ, contribuciones demostradas y trabajo futuro. & Al finalizar el análisis\\
Disponibilidad de datos & DOI de Zenodo, URL OSF, repositorio y licencias. & Después del depósito\\
\end{tabularx}
\end{table}

La redacción de metodología seguirá una estructura reproducible y coherente con las guías de investigación empírica en Ingeniería de Software \cite{jedlitschka2008,runeson2009,wohlin2012}. El trabajo relacionado priorizará estudios primarios recientes sobre Ingeniería de Requisitos, modelos generativos y explicabilidad, incluyendo los trabajos que fundamentan el protocolo \cite{cheng2026,arora2024,vogelsang2024,chazette2020,chazette2022}.

\subsection{Flujo de preparación y publicación}

El proceso editorial seguirá una secuencia que preserve la trazabilidad y evite publicar resultados no reproducibles:

\begin{enumerate}
  \item cerrar y registrar el protocolo en OSF antes de ejecutar el estudio;
  \item recolectar y verificar las evidencias de campo y validación;
  \item anonimizar los datos y conservar los originales en la zona restringida;
  \item ejecutar los scripts desde un entorno limpio y generar tablas y figuras;
  \item redactar resultados y discusión únicamente a partir de esas salidas;
  \item completar el paquete Zenodo, el diccionario y la licencia;
  \item comparar candidatas editoriales con las herramientas oficiales;
  \item adaptar el manuscrito a la plantilla de la revista seleccionada;
  \item depositar el dataset, incorporar el DOI y congelar la versión de envío.
\end{enumerate}

La Figura~\ref{fig:flujo-preparacion-publicacion} presenta el flujo desde la consolidación de resultados hasta el depósito final del conjunto de datos. Su aplicación deberá respetar la secuencia ética y reproducible descrita en esta sección.

\FiguraDocumento
  {07_Publicacion/figuras/FIG-14_Flujo_Preparacion_Publicacion.jpg}
  {0.94}
  {Flujo de preparación del manuscrito y del conjunto de datos}
  {fig:flujo-preparacion-publicacion}

\subsection{Selección de la revista objetivo}

La revista no se fijará por preferencia personal. La selección se realizará después de disponer de un título y un resumen preliminar coherentes, utilizando las herramientas oficiales de Springer Nature, Elsevier e IEEE indicadas en la guía. Para cada editorial se registrará al menos una candidata de acceso abierto con APC y otra por suscripción o híbrida sin cargo obligatorio para las personas autoras.

\begin{table}[H]
\centering
\caption{Herramientas y criterios para la búsqueda de revista}
\label{tab:herramientas-revista}
\small
\rowcolors{2}{DocLight}{white}
\begin{tabularx}{\textwidth}{L{3.0cm} L{4.25cm} Y}
\rowcolor{DocNavy}
\textcolor{white}{\textbf{Editorial}} &
\textcolor{white}{\textbf{Herramienta oficial}} &
\textcolor{white}{\textbf{Información que deberá registrarse}}\\
Springer Nature & \nolinkurl{journalsuggester.springer.com} & Ajuste temático, JCR, cuartil, modelo de acceso, APC y tiempo a primera decisión.\\
Elsevier & \nolinkurl{journalfinder.elsevier.com} & Coincidencia con resumen, CiteScore, factor de impacto, acceso, APC y tiempo editorial.\\
IEEE & \nolinkurl{publication-recommender.ieee.org} & Coincidencia temática, tipo de publicación, factor de impacto, acceso abierto y enlace de envío.\\
\end{tabularx}
\end{table}

La comparación definitiva se registrará en \nolinkurl{07_Publicacion/analisis_revistas.md}. La Figura~\ref{fig:comparacion-revistas} incorpora una comparación preliminar de candidatas; sus nombres, indexación, cuartiles, modalidad de acceso y costos deberán verificarse nuevamente en las herramientas oficiales al momento de realizar la búsqueda, porque pueden cambiar.

\begin{table}[H]
\centering
\caption{Matriz de evaluación de candidatas editoriales}
\label{tab:matriz-revistas}
\scriptsize
\rowcolors{2}{DocLight}{white}
\begin{tabularx}{\textwidth}{L{2.45cm} L{1.45cm} L{1.35cm} L{1.35cm} L{1.55cm} L{1.35cm} Y}
\rowcolor{DocNavy}
\textcolor{white}{\textbf{Revista}} &
\textcolor{white}{\textbf{Editorial}} &
\textcolor{white}{\textbf{JCR / Q}} &
\textcolor{white}{\textbf{Modelo}} &
\textcolor{white}{\textbf{APC USD}} &
\textcolor{white}{\textbf{1.ª decisión}} &
\textcolor{white}{\textbf{Justificación de ajuste}}\\
Candidata OA 1 & Springer & Pendiente & OA & Pendiente & Pendiente & Se completará con la coincidencia oficial del título y resumen.\\
Candidata sin APC 1 & Springer & Pendiente & Suscripción / híbrida & Pendiente & Pendiente & Se verificará alcance, tiempo y viabilidad económica.\\
Candidata OA 2 & Elsevier & Pendiente & OA & Pendiente & Pendiente & Se completará desde Journal Finder.\\
Candidata sin APC 2 & Elsevier & Pendiente & Suscripción / híbrida & Pendiente & Pendiente & Se verificará la política editorial vigente.\\
Candidata OA 3 & IEEE & Pendiente & OA & Pendiente & Pendiente & Se completará desde Publication Recommender.\\
Candidata sin APC 3 & IEEE & Pendiente & Suscripción / híbrida & Pendiente & Pendiente & Se comparará con las demás candidatas.\\
\end{tabularx}
\end{table}

La Figura~\ref{fig:comparacion-revistas} muestra la comparación preliminar disponible. Esta visualización no sustituye la verificación editorial exigida ni convierte en definitivos los valores pendientes de la Tabla~\ref{tab:matriz-revistas}.

\FiguraDocumento
  {07_Publicacion/figuras/FIG-15_Comparacion_Revistas_Candidatas.jpg}
  {0.94}
  {Comparación de revistas candidatas para el manuscrito de RoutTech}
  {fig:comparacion-revistas}

\subsection{Reproducibilidad y control de versiones}

Cada tabla o figura del manuscrito deberá tener una ruta de reproducción documentada. El README del experimento indicará el sistema operativo, versión de Python o R, dependencias, comando de ejecución, archivos de entrada y salidas esperadas. Los scripts no sobrescribirán los datos crudos; producirán un directorio separado de datos procesados y resultados.

La versión enviada del manuscrito se asociará a un commit y a una versión etiquetada del repositorio. El archivo \nolinkurl{CHANGELOG.md} registrará los cambios entre versiones, mientras que \nolinkurl{checksums.sha256} permitirá verificar la integridad de los archivos multimedia. El archivo \nolinkurl{CITATION.cff} incluirá los nombres de los autores, título, versión, fecha y DOI cuando este exista, de acuerdo con el formato de citación de software \cite{druskat2021}.

\begin{table}[H]
\centering
\caption{Controles mínimos de reproducibilidad}
\label{tab:controles-reproducibilidad}
\small
\rowcolors{2}{DocLight}{white}
\begin{tabularx}{\textwidth}{L{4.0cm} Y}
\rowcolor{DocNavy}
\textcolor{white}{\textbf{Control}} &
\textcolor{white}{\textbf{Criterio verificable}}\\
Entorno & Archivo de dependencias y versión exacta del lenguaje o contenedor utilizado.\\
Datos de entrada & Rutas relativas, archivos anonimizados y hash de la versión analizada.\\
Ejecución & Un comando o script principal reproduce el análisis completo sin pasos manuales ocultos.\\
Salidas & Cada tabla y figura tiene nombre estable y se genera en formato abierto o vectorial.\\
Trazabilidad & La salida identifica RQ, variable, script y archivo fuente utilizado.\\
Control de versiones & Commit, etiqueta de versión y CHANGELOG asociados a la entrega o envío.\\
Verificación & Segundo integrante ejecuta el análisis desde una copia limpia y registra el resultado.\\
\end{tabularx}
\end{table}

\subsection{Declaración de disponibilidad de datos y materiales}

Cuando los depósitos estén disponibles, el manuscrito incluirá una declaración equivalente a la siguiente, sustituyendo los campos pendientes por identificadores verificables:

\begin{tcolorbox}[
  colback=DocLighter,
  colframe=DocRule,
  boxrule=0.6pt,
  arc=1mm,
  title={Texto previsto para disponibilidad de datos y materiales},
  colbacktitle=DocLight,
  coltitle=DocText,
  fonttitle=\bfseries
]
El conjunto de datos anonimizado, el diccionario de datos y los scripts de análisis asociados a este estudio se encuentran disponibles en Zenodo bajo licencia CC BY 4.0, DOI: \texttt{PENDIENTE}. El protocolo fue registrado previamente en OSF en la dirección persistente \texttt{PENDIENTE}. El código del MVP y la versión reproducible de los artefactos se encuentran en el repositorio GitHub de RoutTech, versión \texttt{PENDIENTE}. Los consentimientos originales, videos, audios y demás datos personales identificables no se publican y permanecen en la zona restringida conforme al protocolo ético.
\end{tcolorbox}

Esta declaración no deberá presentarse como definitiva hasta comprobar que los enlaces funcionan, el DOI resuelve correctamente, las licencias coinciden y el contenido depositado es exactamente el descrito.

\subsection{Estado de preparación para publicación}

La Tabla~\ref{tab:estado-publicacion} resume el estado documental sin confundir una plantilla preparada con un artefacto terminado.

\begin{table}[H]
\centering
\caption{Estado de los artefactos de publicación}
\label{tab:estado-publicacion}
\small
\rowcolors{2}{DocLight}{white}
\begin{tabularx}{\textwidth}{L{4.25cm} L{3.0cm} Y}
\rowcolor{DocNavy}
\textcolor{white}{\textbf{Artefacto}} &
\textcolor{white}{\textbf{Estado actual}} &
\textcolor{white}{\textbf{Acción necesaria}}\\
Estructura de dataset & Preparada & Incorporar únicamente datos anonimizados y validados.\\
Diccionario de datos & Plantilla preparada & Completar variables después de congelar los instrumentos.\\
Scripts de análisis & Pendientes & Implementar, versionar y probar desde un entorno limpio.\\
Registro OSF & Pendiente de verificación & Registrar antes de ejecutar el estudio y guardar el comprobante.\\
Manuscrito preliminar & Estructurado & Adaptar a la plantilla oficial y cerrar introducción, relacionados y metodología.\\
Análisis de revistas & Plantilla preparada & Ejecutar búsquedas oficiales con título y resumen; registrar evidencias.\\
Depósito Zenodo & Pendiente & Crear después de anonimizar y validar el paquete.\\
DOI y declaración de datos & Pendientes & Incorporar únicamente después de publicar el depósito.\\
\end{tabularx}
\end{table}

\subsection{Cierre de la sección}

La preparación editorial de RoutTech queda definida mediante una estructura de datos FAIR, un manuscrito paralelo trazable, un procedimiento de selección de revista y controles de reproducibilidad. La sección no declara que el dataset, el DOI, el registro OSF ni la selección editorial estén terminados; identifica exactamente qué debe producirse, dónde deberá almacenarse y qué verificaciones se requieren antes de presentar esos elementos como concluidos.



% =====================================================
% 9. CONCLUSIONES Y PLAN DE CIERRE
% =====================================================
\section{Conclusiones y plan de cierre de la Entrega 3 (2A)}

\subsection{Síntesis de la especificación consolidada}

La presente versión consolida la especificación de RoutTech desde la definición del problema y del límite del producto hasta la planificación del MVP, el componente empírico y la preparación de publicación. El documento establece una cadena de trazabilidad que conecta las fuentes de evidencia con requisitos funcionales, requisitos no funcionales, obligaciones legales, restricciones de diseño, historias de usuario, criterios de aceptación, casos de uso, componentes y mockups. Esta organización permite que cada decisión futura de diseño o implementación pueda contrastarse con una necesidad identificada y con un criterio de verificación.

La especificación contiene 25 requisitos funcionales con atributos de identificación, descripción, actor, entradas, salidas, precondiciones, postcondiciones, prioridad, verificación y evidencia; 15 requisitos no funcionales cuantificados distribuidos sobre las nueve características de calidad de ISO/IEC 25010:2023; requisitos legales relacionados con la protección de datos personales; y ocho restricciones de diseño. También se incorporaron historias de usuario Connextra para los requisitos clasificados como \textit{Must have}, revisión INVEST, escenarios Gherkin y la estructura de los modelos UML requeridos. Estas cantidades satisfacen la cobertura documental mínima definida para la Entrega 3, aunque su aceptación definitiva depende de la existencia y verificación de los artefactos asociados en el repositorio.

\subsection{Cobertura documental alcanzada}

La Tabla~\ref{tab:resumen-cobertura-ers} resume la cobertura lograda dentro del ERS/SRS. La condición ``estructurado'' significa que el contenido y su ubicación están definidos; no implica que una evidencia externa, un diagrama, un mockup o un repositorio pendiente ya exista o haya sido validado.

\begin{table}[H]
\centering
\caption{Resumen de cobertura documental de la Entrega 3 (2A)}
\label{tab:resumen-cobertura-ers}
\small
\rowcolors{2}{DocLight}{white}
\begin{tabularx}{\textwidth}{L{4.1cm} L{2.75cm} Y}
\rowcolor{DocNavy}
\textcolor{white}{\textbf{Elemento}} &
\textcolor{white}{\textbf{Estado documental}} &
\textcolor{white}{\textbf{Observación de cierre}}\\
Introducción y alcance & Consolidado & Verificar que las fuentes empíricas citadas existan en el repositorio y conserven su código.\\
Contexto y stakeholders & Estructurado & Insertar los diagramas definitivos en los espacios reservados y conservar sus archivos editables.\\
Requisitos funcionales & 25 RF documentados & Confirmar cada relación RF--EV contra transcripciones, cuestionario, documentos o actas reales.\\
Requisitos no funcionales & 15 RNF cuantificados & Ejecutar las pruebas o validaciones definidas para comprobar sus valores objetivo.\\
Requisitos legales y RD & Documentados & Revisar el mapeo legal con la documentación ética vigente y la política real de tratamiento.\\
Historias y aceptación & Estructuradas & Validar con stakeholders y registrar decisiones en las actas de walkthrough.\\
Modelado UML & Catálogo y especificaciones preparados & Insertar fuentes editables y exportaciones PNG/SVG de todos los diagramas obligatorios.\\
Priorización y trazabilidad & 48 filas preparadas & Revisar consistencia de identificadores y exportar los CSV finales desde una única fuente maestra.\\
MVP & Alcance y cobertura previstos & Sustituir la planificación por un repositorio ejecutable, commit verificable y video de demostración.\\
Componente empírico & Protocolo estructurado & Registrar en OSF antes de ejecutar las actividades que todavía estén pendientes.\\
Preparación de publicación & Plantillas preparadas & Completar datos, scripts, análisis editorial, Zenodo y DOI únicamente después de la verificación.\\
\end{tabularx}
\end{table}

\subsection{Estado de los gatekeepers}

La Tabla~\ref{tab:estado-gatekeepers} funciona como control preventivo. Se utiliza una evaluación conservadora: cuando la existencia o validez de un artefacto no puede demostrarse únicamente desde el ERS, se mantiene como pendiente de verificación.

\begin{table}[H]
\centering
\caption{Estado preventivo de los gatekeepers de la Entrega 3 (2A)}
\label{tab:estado-gatekeepers}
\scriptsize
\rowcolors{2}{DocLight}{white}
\begin{tabularx}{\textwidth}{C{0.75cm} L{3.4cm} L{2.65cm} Y}
\rowcolor{DocNavy}
\textcolor{white}{\textbf{ID}} &
\textcolor{white}{\textbf{Condición}} &
\textcolor{white}{\textbf{Estado actual}} &
\textcolor{white}{\textbf{Acción antes del corte}}\\
G1 & Entrega oportuna y repositorio verificable & Pendiente de comprobación final & Probar el enlace desde una sesión sin autenticación y mantenerlo en una sola línea de la portada.\\
G2 & PDF único y reproducible desde LaTeX & En desarrollo avanzado & Compilar desde una copia limpia mediante el README y verificar que no existan referencias rotas.\\
G3 & Sistema real y aval institucional & Pendiente de verificación & Confirmar que el aval y el seudónimo institucional estén archivados en \nolinkurl{08_Etica/}.\\
G4 & Evidencia audiovisual real, íntegra y verificable & Pendiente de verificación técnica & Contrastar fichas, hashes y contenido del contenedor cifrado antes de declarar cumplimiento.\\
G5 & Mínimos de RF, RNF, CU, trazabilidad y campo & Cumplimiento documental parcial & Los conteos del ERS están cubiertos; los mínimos de campo deben verificarse archivo por archivo.\\
G6 & Registro previo del protocolo en OSF & Pendiente & Registrar el protocolo antes de cualquier actividad experimental aún no ejecutada y guardar el comprobante.\\
G7 & Explicabilidad para componentes de IA/AA & Abordado en la especificación & Mantener el RNF de explicabilidad, sus métricas y su validación con usuarios.\\
G8 & Documentación ética vigente y adenda previa & Pendiente de verificación & Revisar paquete base, categoría de riesgo, adenda y fechas de consentimientos en \nolinkurl{08_Etica/}.\\
\end{tabularx}
\end{table}

La Figura~\ref{fig:mapa-cierre-2a} presenta un mapa preliminar de cierre y verificación. Los estados mostrados deberán contrastarse con el checklist final, los gatekeepers y las evidencias reales del repositorio antes de declararlos como definitivos.

\FiguraDocumento
  {09_Cierre/figuras/FIG-16_Mapa_Cierre_Entrega_2A.jpg}
  {0.94}
  {Mapa de cierre y verificación de la Entrega 3 (2A)}
  {fig:mapa-cierre-2a}

\textit{Nota de control:} las marcas ``Completo'' y ``En revisión'' de la Figura~\ref{fig:mapa-cierre-2a} son un resumen visual preliminar. Prevalece el estado conservador documentado en la Tabla~\ref{tab:estado-gatekeepers} hasta concluir la verificación técnica y ética de cada artefacto.

\subsection{Limitaciones de la versión actual}

Esta versión no constituye por sí sola evidencia de que todas las actividades de campo, validación, desarrollo o publicación hayan concluido. Permanecen como limitaciones principales:

\begin{itemize}
  \item varios diagramas UML todavía se representan mediante espacios reservados y deberán sustituirse por sus archivos definitivos; las figuras organizacionales, los mockups y las visualizaciones ya incorporadas deberán contrastarse con los artefactos editables y la evidencia real;
  \item la trazabilidad a evidencia requiere una comprobación final contra los archivos reales, su zona de almacenamiento y sus hashes;
  \item la cobertura del MVP corresponde a una planificación y debe verificarse sobre código ejecutable y un commit estable;
  \item los resultados del componente empírico no pueden redactarse hasta completar el protocolo, la ejecución y el análisis reproducible;
  \item el registro OSF, el depósito Zenodo, el DOI y la revista objetivo no deben presentarse como concluidos mientras permanezcan pendientes;
  \item las cifras de entrevistas y cuestionarios registradas en la línea base no sustituyen la verificación de los mínimos técnicos, éticos y de diversidad de perfiles.
\end{itemize}

Estas limitaciones no invalidan la estructura del ERS/SRS, pero determinan las tareas necesarias para convertir la especificación en una entrega verificable y reproducible.

\subsection{Plan de cierre documental y técnico}

El cierre de la Entrega 3 seguirá el siguiente orden para evitar inconsistencias entre el PDF y el repositorio:

\begin{enumerate}
  \item congelar el catálogo de evidencias y validar los códigos utilizados en RF, RNF, CU, HU y matriz de trazabilidad;
  \item verificar consentimientos, archivos multimedia, transcripciones, fotografías, cuestionarios, documentos y walkthroughs según su zona pública o restringida;
  \item completar la codificación temática y contrastar que cada requisito sustantivo tenga respaldo verificable;
  \item producir e insertar todos los diagramas UML, diagramas organizacionales y mockups en los espacios reservados;
  \item actualizar la matriz de trazabilidad y la priorización desde los archivos CSV definitivos;
  \item ejecutar el MVP desde una copia limpia, registrar el commit evaluado y producir el video de demostración;
  \item finalizar y registrar el protocolo en OSF antes de ejecutar cualquier fase experimental pendiente;
  \item compilar el documento completo, revisar tablas, figuras, citas, enlaces, numeración y referencias cruzadas;
  \item ejecutar la autocertificación del repositorio y corregir cualquier gatekeeper pendiente antes del corte.
\end{enumerate}

\subsection{Conclusiones}

RoutTech queda definido como un sistema de movilidad colaborativa universitaria cuya especificación prioriza la verificación de identidad, la coordinación de rutas y cupos, la protección de datos, la reputación y la explicación de recomendaciones. La estructura desarrollada permite evaluar el producto no solo por las funciones que ofrece, sino también por atributos medibles de calidad, restricciones legales y criterios de aceptación observables.

La contribución principal de esta versión es la integración de requisitos, modelado, priorización, trazabilidad, MVP y protocolo empírico dentro de un mismo documento reproducible. La especificación establece qué debe construir el sistema, cómo deberá comprobarse y qué evidencia debe conservarse. No obstante, el cumplimiento final dependerá de que los artefactos declarados estén presentes, sean auténticos, mantengan su integridad y correspondan exactamente con los identificadores utilizados en el ERS.

Por tanto, el siguiente trabajo no consiste en ampliar el documento con afirmaciones nuevas, sino en cerrar y verificar los elementos pendientes: evidencias, imágenes, modelos editables, mockups, código ejecutable, registro OSF y controles de publicación. Una vez completados esos elementos, el ERS/SRS podrá funcionar como línea base formal para la validación final, el desarrollo del MVP y la preparación de la Entrega 4.



% =====================================================
% APENDICES
% =====================================================
\clearpage
\clearpage
\phantomsection
\printbibliography[heading=bibintoc,title={Referencias}]

\clearpage
\appendix
\phantomsection
\addcontentsline{toc}{section}{APÉNDICES}
\begin{center}
  {\Large\bfseries\color{DocNavy} APÉNDICES}
\end{center}
\vspace{0.5em}

Los apéndices reúnen los índices, controles y formatos de apoyo que permiten comprobar la relación entre el ERS/SRS y los artefactos del repositorio. No sustituyen los archivos originales ni las evidencias primarias: cada elemento citado debe existir en la ruta indicada, conservar su identificador y corresponder con la zona pública o restringida definida para el proyecto.

% -----------------------------------------------------
\section{Catálogo maestro de evidencias}
\label{ap:catalogo-evidencias}

El catálogo maestro se utiliza como punto único de control para los identificadores de evidencia empleados en requisitos, historias de usuario, casos de uso, trazabilidad y validación. El registro definitivo deberá mantenerse en un archivo editable dentro de \nolinkurl{02_Evidencias/} y concordar con los nombres físicos del repositorio.

\begin{longtable}{L{2.15cm} L{2.40cm} L{3.25cm} L{3.20cm} L{2.20cm}}
\caption{Catálogo maestro preliminar de evidencias de RoutTech}\label{tab:catalogo-maestro-evidencias}\\
\rowcolor{DocNavy}
\textcolor{white}{\textbf{ID}} &
\textcolor{white}{\textbf{Tipo}} &
\textcolor{white}{\textbf{Contenido}} &
\textcolor{white}{\textbf{Ruta de control}} &
\textcolor{white}{\textbf{Estado}}\\
\endfirsthead
\rowcolor{DocNavy}
\textcolor{white}{\textbf{ID}} &
\textcolor{white}{\textbf{Tipo}} &
\textcolor{white}{\textbf{Contenido}} &
\textcolor{white}{\textbf{Ruta de control}} &
\textcolor{white}{\textbf{Estado}}\\
\endhead
EV-ENT-001 a EV-ENT-016 & Entrevistas & Dieciséis entrevistas y sus transcripciones anonimizadas, con audio o video verificable según corresponda. & \nolinkurl{02_Evidencias/Transcripciones/} y contenedor restringido & Verificar archivo por archivo\\
EV-CUE-001 & Cuestionario & Instrumento aplicado y exportación depurada de 43 respuestas declaradas en la línea base. & \nolinkurl{02_Evidencias/Cuestionario/Respuestas/} & Verificar exportación y fotografías\\
EV-DOC-001 a EV-DOC-004 & Documentos & Documentos de apoyo institucional u organizacional empleados para contextualizar procesos, políticas o restricciones. & \nolinkurl{02_Evidencias/Documentos_Organizacion/} & Verificar anonimización y tipo\\
EV-VAL-001 a EV-VAL-003 & Walkthrough & Actas de las tres sesiones mínimas de validación: dos con usuarios técnicos y una con usuario no técnico. & \nolinkurl{02_Evidencias/Validacion_Walkthrough/} & Pendiente de consolidación\\
EV-OBS-001 & Observación & Registro de observación del entorno y notas de campo asociadas. & \nolinkurl{02_Evidencias/Fotos_Entorno/} & Pendiente de catálogo final\\
EV-CONS-001 a EV-CONS-008 & Consentimientos & Copias públicas enmascaradas y originales cifrados de participantes distintos. & \nolinkurl{02_Evidencias/Consentimientos/} y contenedor restringido & Verificar correspondencia\\
EV-FOT-001 a EV-FOT-015 & Fotografías & Fotografías fechadas del entorno, sin rostros identificables y sin coordenadas GPS. & \nolinkurl{02_Evidencias/Fotos_Entorno/} & Pendiente de numeración final\\
EV-APL-001 a EV-APL-005 & Aplicación de cuestionario & Fotografías de aplicación real del cuestionario con anonimización visual. & \nolinkurl{02_Evidencias/Cuestionario/Fotos_Aplicacion/} & Pendiente de numeración final\\
\end{longtable}

\subsection{Reglas de identificación}

\begin{itemize}
  \item Cada identificador debe aparecer una sola vez en el catálogo y conservarse durante todo el proyecto.
  \item El código del participante reemplaza al nombre propio en archivos públicos, transcripciones, actas y tablas.
  \item Las evidencias audiovisuales deben estar vinculadas a una ficha técnica pública con duración, códec, tamaño y hash SHA-256.
  \item Una evidencia solo puede marcarse como verificada cuando el archivo abre correctamente, coincide con la descripción y se encuentra en la zona que le corresponde.
  \item Los rangos reservados en este apéndice deberán ajustarse al catálogo definitivo sin modificar los identificadores ya citados por los requisitos.
\end{itemize}

% -----------------------------------------------------
\section{Instrumentos de elicitación y validación}
\label{ap:instrumentos}

Este apéndice identifica los instrumentos que respaldan la obtención y validación de requisitos. Las versiones completas deberán permanecer en \nolinkurl{06_Experimento/instrumentos/} y las versiones aplicadas o firmadas en la zona de evidencia correspondiente.

\begin{table}[H]
\centering
\caption{Índice de instrumentos del proyecto}
\label{tab:indice-instrumentos}
\small
\rowcolors{2}{DocLight}{white}
\begin{tabularx}{\textwidth}{L{1.45cm} L{3.10cm} Y L{3.25cm}}
\rowcolor{DocNavy}
\textcolor{white}{\textbf{Código}} &
\textcolor{white}{\textbf{Instrumento}} &
\textcolor{white}{\textbf{Finalidad}} &
\textcolor{white}{\textbf{Archivo esperado}}\\
INS-01 & Guion de entrevista semiestructurada & Identificar problemas de movilidad, prácticas actuales, necesidades, riesgos y expectativas de pasajeros y conductores. & \nolinkurl{06_Experimento/instrumentos/Guion_Entrevista_v2.0.pdf}\\
INS-02 & Cuestionario estructurado & Medir frecuencia de movilidad, dificultades, disposición de uso, prioridades y percepción de seguridad. & \nolinkurl{06_Experimento/instrumentos/Cuestionario_v2.0.pdf}\\
INS-03 & Guía de observación & Registrar condiciones del entorno, puntos de encuentro, conectividad y restricciones operativas sin identificar personas. & \nolinkurl{06_Experimento/instrumentos/Guia_Observacion.pdf}\\
INS-04 & Acta de walkthrough & Documentar requisitos revisados, acuerdos, cambios solicitados, responsables y decisión de aceptación. & \nolinkurl{06_Experimento/instrumentos/Plantilla_Acta_Walkthrough.pdf}\\
INS-05 & Consentimiento informado & Registrar participación voluntaria, tratamiento de datos, uso académico y condiciones de publicación anonimizada. & \nolinkurl{06_Experimento/instrumentos/Consentimiento_Informado_v2.0.pdf}\\
INS-06 & Rúbrica de explicabilidad & Valorar claridad, utilidad, oportunidad, extensión y comprensión de las explicaciones del recomendador. & \nolinkurl{06_Experimento/instrumentos/Rubrica_Explicabilidad.pdf}\\
\end{tabularx}
\end{table}

\subsection{Control de versiones de instrumentos}

\begin{table}[H]
\centering
\caption{Control de cambios de los instrumentos}
\label{tab:versiones-instrumentos}
\small
\rowcolors{2}{DocLight}{white}
\begin{tabularx}{\textwidth}{L{1.70cm} L{1.60cm} L{2.55cm} Y L{2.40cm}}
\rowcolor{DocNavy}
\textcolor{white}{\textbf{Instrumento}} &
\textcolor{white}{\textbf{Versión}} &
\textcolor{white}{\textbf{Fecha}} &
\textcolor{white}{\textbf{Cambio principal}} &
\textcolor{white}{\textbf{Aprobación}}\\
INS-01 & v2.0 & Pendiente de registrar & Ajuste de preguntas para la segunda ronda y separación por perfil. & Adenda ética\\
INS-02 & v2.0 & Pendiente de registrar & Inclusión de variables del perfil dominante y explicabilidad. & Adenda ética\\
INS-03 & v1.0 & Pendiente de registrar & Formato de observación y control de privacidad. & Equipo del proyecto\\
INS-04 & v1.0 & Pendiente de registrar & Estructura de acuerdos, cambios y firmas. & Equipo del proyecto\\
INS-05 & v2.0 & Pendiente de registrar & Cláusula de uso anonimizado en publicación y dataset. & Adenda ética\\
INS-06 & v1.0 & Pendiente de registrar & Métricas de claridad y comprensión de explicaciones. & Protocolo OSF\\
\end{tabularx}
\end{table}

% -----------------------------------------------------
\section{Índice multimedia y fichas técnicas}
\label{ap:multimedia}

Los archivos audiovisuales originales permanecen dentro del contenedor cifrado. En la zona pública se conserva la información técnica necesaria para comprobar la integridad de cada archivo. La Tabla~\ref{tab:plantilla-ficha-tecnica} registra las 16 entrevistas realizadas y coincide con el archivo editable \nolinkurl{02_Evidencias/fichas_tecnicas.csv}. Los valores de fecha, duración, códec, tamaño y hash SHA-256 fueron incorporados desde la ficha técnica entregada por el equipo.

\begin{landscape}
\scriptsize
\setlength{\tabcolsep}{3.5pt}
\begin{longtable}{L{3.35cm} L{1.35cm} L{1.90cm} L{1.55cm} L{1.65cm} L{1.45cm} L{1.75cm} L{6.50cm}}
\caption{Ficha técnica consolidada de las entrevistas multimedia de RoutTech}\label{tab:plantilla-ficha-tecnica}\\
\rowcolor{DocNavy}
\textcolor{white}{\textbf{Archivo registrado}} &
\textcolor{white}{\textbf{Tipo}} &
\textcolor{white}{\textbf{Fecha}} &
\textcolor{white}{\textbf{Código}} &
\textcolor{white}{\textbf{Duración}} &
\textcolor{white}{\textbf{Códec}} &
\textcolor{white}{\textbf{Tamaño}} &
\textcolor{white}{\textbf{SHA-256}}\\
\endfirsthead
\rowcolor{DocNavy}
\textcolor{white}{\textbf{Archivo registrado}} &
\textcolor{white}{\textbf{Tipo}} &
\textcolor{white}{\textbf{Fecha}} &
\textcolor{white}{\textbf{Código}} &
\textcolor{white}{\textbf{Duración}} &
\textcolor{white}{\textbf{Códec}} &
\textcolor{white}{\textbf{Tamaño}} &
\textcolor{white}{\textbf{SHA-256}}\\
\endhead
Video de entrevista (MP4) & Video & 2026-07-21 & ENT-001 & 00:07:16 & H.264 & 15,9 MB & \seqsplit{72A535129323FAF1F145467A5DC25E0F4075DD37F9B0FA260C5E664377D03F51}\\
Video de entrevista (MP4) & Video & 2026-07-21 & ENT-002 & 00:06:38 & H.264 & 14,6 MB & \seqsplit{346426548E674D6F47BB63D09957906D7AB342DD2A58B42BE17BF7DF688053CD}\\
Video de entrevista (MP4) & Video & 2026-07-21 & ENT-003 & 00:07:01 & H.264 & 17,9 MB & \seqsplit{675D6DCE2BBCB6913377725C3781118105A8B4E7E57F05A2BDFAEEC8F88AAD96}\\
Video de entrevista (MP4) & Video & 2026-07-21 & ENT-004 & 00:06:10 & H.264 & 11,9 MB & \seqsplit{BCFBD0E623BFFFDA8ED5228D388EE3FF536A36487CB31AE5D3B55F2CD00D8737}\\
Video de entrevista (MP4) & Video & 2026-07-22 & ENT-005 & 00:06:29 & H.264 & 23,4 MB & \seqsplit{348F14B76E3F999A634ABAF73FC48235507A214EBE8244EED1E37FF22970E70F}\\
Video de entrevista (MP4) & Video & 2026-07-22 & ENT-006 & 00:10:59 & H.264 & 21,8 MB & \seqsplit{7BDE0450FDB59C8A68CFC10BFB9ECD4B0B7373C662815B03D1965FA098B05400}\\
Video de entrevista (MP4) & Video & 2026-07-22 & ENT-007 & 00:08:17 & H.264 & 24,5 MB & \seqsplit{F734FCC1B450C382EC72563F1766A8331B4043A3E9012CF7FF308BF8859F224E}\\
Video de entrevista (MP4) & Video & 2026-07-22 & ENT-008 & 00:08:33 & H.264 & 23,0 MB & \seqsplit{3D655F6C5421CE2DEE2C2F2E1896B430E421EDA3DDF8477D4CE9D54D523A8C99}\\
Video de entrevista (MP4) & Video & 2026-07-22 & ENT-009 & 00:13:16 & H.264 & 24,9 MB & \seqsplit{BBDB760D72B28E7D5C1D22C3D17D80A919E557E1D8F8E6128C2E9638B14923DC}\\
Video de entrevista (MP4) & Video & 2026-07-22 & ENT-010 & 00:05:49 & H.264 & 12,0 MB & \seqsplit{7D359FF94D5C8A2BB46911E7070193EF877E816D0CE2014D4AD9D9C1A628B35A}\\
Video de entrevista (MP4) & Video & 2026-07-22 & ENT-011 & 00:06:53 & H.264 & 23,2 MB & \seqsplit{E582DCF8A8247352349496835B0AE93570318BBB8E3A7F324DC2EDCE0E21761D}\\
Video de entrevista (MP4) & Video & 2026-07-22 & ENT-012 & 00:06:12 & H.264 & 14,5 MB & \seqsplit{53163BDCA6953EF10811A9FB48CE0904FD284A78F327BBEDCD961127C5D873D0}\\
Video de entrevista (MP4) & Video & 2026-07-23 & ENT-013 & 00:05:21 & H.264 & 15,1 MB & \seqsplit{34C7D4FB885951AA856094D67D8C2C1BBB4F7F9397598F797F0E817DDF37FFF8}\\
Video de entrevista (MP4) & Video & 2026-07-23 & ENT-014 & 00:05:09 & H.264 & 15,5 MB & \seqsplit{E1FF5379D17BFF33592358F7F8A3A60F58E74E3CCA212E6853863F83ED7D5D39}\\
Video de entrevista (MP4) & Video & 2026-07-25 & ENT-015 & 00:06:16 & H.264 & 16,9 MB & \seqsplit{DCF2C12A17D0D8B7FEA80EDB15091D6FA8EB1C457C5D219BB4E3A2082EBDC1C3}\\
Video de entrevista (MP4) & Video & 2026-07-25 & ENT-016 & 00:07:04 & H.264 & 16,1 MB & \seqsplit{4B4C6B4ADD075C734E1B199D0F3032FF53664C81C34545D912858D6DA603342B}\\
\end{longtable}
\end{landscape}

La nomenclatura física que deberá utilizarse es \nolinkurl{AAAA-MM-DD_TipoParticipante_CodigoParticipante_Tecnica.ext}. El nombre exacto de cada archivo se registra en el CSV; la tabla anterior muestra únicamente los campos obligatorios para evitar cortes tipográficos dentro del PDF.

\subsection{Comprobaciones técnicas mínimas}

\begin{enumerate}
  \item comprobar con \texttt{ffprobe} que el archivo sea multimedia real y que su duración sea mayor que cero;
  \item calcular el hash SHA-256 antes de cifrar y registrar exactamente el mismo valor en la ficha pública;
  \item verificar que el nombre use código de participante y no contenga datos personales;
  \item contrastar que el archivo figure dentro del contenedor restringido y no exista una copia sin cifrar en la zona pública;
  \item conservar una fila por cada audio, video o grabación de walkthrough declarado en el ERS.
\end{enumerate}

% -----------------------------------------------------
\section{Trazabilidad y priorización}
\label{ap:trazabilidad-priorizacion}

La matriz completa se mantiene en \nolinkurl{04_Trazabilidad/matriz_trazabilidad.csv}; la priorización se mantiene en \nolinkurl{04_Trazabilidad/priorizacion_moscow_kano.csv}. Este apéndice define la estructura que debe conservarse durante la revisión final.

\begin{table}[H]
\centering
\caption{Campos obligatorios de la matriz de trazabilidad extendida}
\label{tab:campos-trazabilidad-anexo}
\small
\rowcolors{2}{DocLight}{white}
\begin{tabularx}{\textwidth}{L{3.25cm} Y}
\rowcolor{DocNavy}
\textcolor{white}{\textbf{Campo}} &
\textcolor{white}{\textbf{Uso}}\\
ID & Identificador único de la fila de trazabilidad.\\
Ley y Artículo & Base legal aplicable cuando el requisito involucra datos personales, seguridad o derechos del titular.\\
Objetivo & Objetivo del sistema al que contribuye el requisito.\\
Stakeholder & Parte interesada que origina, utiliza, valida o supervisa la necesidad.\\
ID-EV & Evidencia primaria o documental que respalda la afirmación.\\
ID-RF y Tipo & Identificador del RF, RNF o RD y clasificación correspondiente.\\
ID-CU & Caso de uso que operacionaliza el comportamiento.\\
ID-HU e ID-CA & Historia de usuario y criterio de aceptación relacionado.\\
ID-Componente & Módulo o servicio responsable de implementar el requisito.\\
ID-Mockup & Interfaz que permite observar o validar el comportamiento.\\
\end{tabularx}
\end{table}

\begin{table}[H]
\centering
\caption{Control resumido de la priorización}
\label{tab:control-priorizacion-anexo}
\small
\rowcolors{2}{DocLight}{white}
\begin{tabularx}{\textwidth}{L{2.35cm} L{2.35cm} L{2.35cm} L{2.35cm} Y}
\rowcolor{DocNavy}
\textcolor{white}{\textbf{ID-RF}} &
\textcolor{white}{\textbf{MoSCoW}} &
\textcolor{white}{\textbf{Kano}} &
\textcolor{white}{\textbf{WSJF}} &
\textcolor{white}{\textbf{Justificación}}\\
RF-XX & Must / Should / Could / Won't & Básica / desempeño / deleite / indiferente / rechazo & Valor calculado & Explicación breve sustentada en riesgo, valor y dependencia.\\
\end{tabularx}
\end{table}

% -----------------------------------------------------
\section{Índice de modelado UML y mockups}
\label{ap:modelado-mockups}

Los modelos deberán publicarse en formato editable y como exportaciones PNG y SVG. La referencia incluida en el cuerpo del ERS debe coincidir con el identificador, el nombre y la ruta de este índice.

\begin{longtable}{L{1.55cm} L{3.25cm} L{4.00cm} L{4.35cm}}
\caption{Índice de diagramas y prototipos}\label{tab:indice-modelos-anexo}\\
\rowcolor{DocNavy}
\textcolor{white}{\textbf{Código}} &
\textcolor{white}{\textbf{Artefacto}} &
\textcolor{white}{\textbf{Cobertura}} &
\textcolor{white}{\textbf{Ruta esperada}}\\
\endfirsthead
\rowcolor{DocNavy}
\textcolor{white}{\textbf{Código}} &
\textcolor{white}{\textbf{Artefacto}} &
\textcolor{white}{\textbf{Cobertura}} &
\textcolor{white}{\textbf{Ruta esperada}}\\
\endhead
MOD-01 & Diagrama general de casos de uso & Actores y funciones principales del sistema. & \nolinkurl{03_Modelado/Diagramas_UML/Casos_Uso/}\\
MOD-02 & Especificaciones de casos de uso & Mínimo diez CU detallados asociados a RF Must. & \nolinkurl{03_Modelado/Diagramas_UML/Casos_Uso/Especificaciones/}\\
MOD-03 & Diagrama de clases & Clases, atributos, operaciones, asociaciones y multiplicidades. & \nolinkurl{03_Modelado/Diagramas_UML/Clases/}\\
MOD-04 & Diagramas de secuencia & Interacciones temporales de los CU Must. & \nolinkurl{03_Modelado/Diagramas_UML/Secuencia/}\\
MOD-05 & Diagramas de actividad & Flujos principales, decisiones y excepciones. & \nolinkurl{03_Modelado/Diagramas_UML/Actividad/}\\
MOD-06 & Diagramas de estados & Ciclos de vida de reserva, ruta, solicitud, usuario y reporte. & \nolinkurl{03_Modelado/Diagramas_UML/Estados/}\\
MOD-07 & Diagrama de componentes & Módulos, servicios, interfaces y dependencias técnicas. & \nolinkurl{03_Modelado/Diagramas_UML/Componentes/}\\
MOD-08 & Diagrama de despliegue & Nodos, aplicaciones, servicios y base de datos. & \nolinkurl{03_Modelado/Diagramas_UML/Despliegue/}\\
MU-01 a MU-09 & Mockups de alta fidelidad & Pantallas relacionadas con los RF Must, HU y criterios de aceptación. & \nolinkurl{03_Modelado/Mockups/}\\
\end{longtable}

\subsection{Lista de control de cada modelo}

\begin{itemize}
  \item archivo fuente editable en XMI, PlantUML, QEA o formato equivalente;
  \item exportación PNG con resolución mínima suficiente para lectura en el PDF;
  \item exportación SVG cuando la herramienta lo permita;
  \item identificador visible y nombre coherente con la referencia del ERS;
  \item ausencia de datos personales, capturas externas no autorizadas o marcas de agua ajenas;
  \item correspondencia verificable con RF, CU, HU, componentes o mockups según aplique.
\end{itemize}

% -----------------------------------------------------
\section{Verificación del Producto Mínimo Viable}
\label{ap:verificacion-mvp}

Este apéndice registra los datos que deben completarse cuando el MVP sea ejecutable. La cobertura planificada no se considera validada hasta reproducir el despliegue desde una copia limpia y comprobar el comportamiento descrito.

\begin{table}[H]
\centering
\caption{Ficha de verificación del MVP}
\label{tab:ficha-verificacion-mvp}
\small
\rowcolors{2}{DocLight}{white}
\begin{tabularx}{\textwidth}{L{4.10cm} Y}
\rowcolor{DocNavy}
\textcolor{white}{\textbf{Elemento}} &
\textcolor{white}{\textbf{Dato o criterio de comprobación}}\\
Repositorio público & URL exacta del repositorio Git separado o submódulo asociado.\\
Versión evaluada & Rama, etiqueta y hash completo del commit.\\
Entorno & Sistema operativo, runtime, gestor de paquetes, base de datos y versiones.\\
Comando de despliegue & Instrucción reproducible indicada en el README.\\
Cobertura Must & Número de RF Must implementados dividido para el total de RF Must.\\
Pruebas & Casos ejecutados, resultado esperado, resultado obtenido y evidencia.\\
Video de demostración & Ruta \nolinkurl{05_MVP/video_demo.mp4}, duración menor o igual a tres minutos.\\
Licencia & MIT o Apache-2.0 aplicada al código.\\
Limitaciones & Funciones simuladas, datos sintéticos y dependencias no disponibles.\\
\end{tabularx}
\end{table}

\begin{table}[H]
\centering
\caption{Acta resumida de prueba del MVP}
\label{tab:acta-prueba-mvp}
\small
\rowcolors{2}{DocLight}{white}
\begin{tabularx}{\textwidth}{L{1.65cm} L{2.50cm} Y Y L{1.75cm}}
\rowcolor{DocNavy}
\textcolor{white}{\textbf{Prueba}} &
\textcolor{white}{\textbf{RF asociado}} &
\textcolor{white}{\textbf{Resultado esperado}} &
\textcolor{white}{\textbf{Resultado obtenido}} &
\textcolor{white}{\textbf{Estado}}\\
PR-MVP-01 & Completar & Completar a partir del criterio de aceptación. & Registrar únicamente después de ejecutar. & Pendiente\\
PR-MVP-02 & Completar & Completar a partir del criterio de aceptación. & Registrar únicamente después de ejecutar. & Pendiente\\
PR-MVP-03 & Completar & Completar a partir del criterio de aceptación. & Registrar únicamente después de ejecutar. & Pendiente\\
\end{tabularx}
\end{table}

% -----------------------------------------------------
\section{Control del protocolo experimental y la publicación}
\label{ap:control-experimento-publicacion}

La Tabla~\ref{tab:control-experimental-anexo} diferencia los documentos que pueden prepararse antes del estudio de aquellos que solo pueden cerrarse después de recolectar y analizar datos. Esta separación evita que se presenten como resultados las métricas objetivo o las estructuras todavía pendientes.

\begin{table}[H]
\centering
\caption{Control de artefactos experimentales y de publicación}
\label{tab:control-experimental-anexo}
\small
\rowcolors{2}{DocLight}{white}
\begin{tabularx}{\textwidth}{L{3.50cm} L{3.40cm} Y L{2.40cm}}
\rowcolor{DocNavy}
\textcolor{white}{\textbf{Artefacto}} &
\textcolor{white}{\textbf{Ruta}} &
\textcolor{white}{\textbf{Criterio de cierre}} &
\textcolor{white}{\textbf{Estado}}\\
Protocolo experimental & \nolinkurl{06_Experimento/protocolo.pdf} & PICOC, proposiciones, variables, participantes, instrumentos y plan de análisis completos. & Estructurado\\
Registro OSF & \nolinkurl{06_Experimento/osf_registration.pdf} & Registro previo con URL persistente y fecha anterior a la ejecución pendiente. & Pendiente\\
Instrumentos & \nolinkurl{06_Experimento/instrumentos/} & Versiones aplicadas coherentes con la adenda ética. & Verificar\\
Datos crudos y procesados & \nolinkurl{06_Experimento/resultados/} & Separación de entradas originales y salidas reproducibles. & Pendiente\\
Scripts & \nolinkurl{06_Experimento/scripts_analisis/} & Reproducen todas las tablas y figuras desde los datos crudos. & Pendiente\\
Borrador del manuscrito & \nolinkurl{07_Publicacion/manuscrito_borrador.pdf} & Introducción, trabajo relacionado y metodología verificables; resultados solo después del análisis. & En desarrollo\\
Análisis de revistas & \nolinkurl{07_Publicacion/analisis_revistas.md} & Candidatas verificadas con las herramientas oficiales y datos vigentes. & Pendiente\\
Paquete Zenodo & \nolinkurl{07_Publicacion/dataset_zenodo/} & Dataset anonimizado, diccionario, licencia, scripts y metadatos. & En preparación\\
DOI y disponibilidad & Manuscrito y \nolinkurl{CITATION.cff} & DOI incorporado únicamente después del depósito efectivo. & No disponible aún\\
\end{tabularx}
\end{table}

% -----------------------------------------------------
\section{Checklist final del repositorio}
\label{ap:checklist-repositorio}

La revisión final debe realizarse desde una copia nueva del repositorio y no desde el entorno de trabajo habitual. Cada fila se marcará como cumplida únicamente después de ejecutar la comprobación indicada.

\begin{longtable}{C{0.75cm} L{8.30cm} L{4.45cm}}
\caption{Checklist de aceptación del repositorio}\label{tab:checklist-anexo}\\
\rowcolor{DocNavy}
\textcolor{white}{\textbf{N.º}} &
\textcolor{white}{\textbf{Comprobación}} &
\textcolor{white}{\textbf{Resultado / observación}}\\
\endfirsthead
\rowcolor{DocNavy}
\textcolor{white}{\textbf{N.º}} &
\textcolor{white}{\textbf{Comprobación}} &
\textcolor{white}{\textbf{Resultado / observación}}\\
\endhead
1 & La estructura de carpetas coincide con el árbol obligatorio y contiene \nolinkurl{08_Etica/}. & Pendiente de autocertificación\\
2 & Existen y están completos \nolinkurl{README.md}, \nolinkurl{LICENSE}, \nolinkurl{CITATION.cff}, \nolinkurl{CHANGELOG.md} y \nolinkurl{checksums.sha256}. & Pendiente de autocertificación\\
3 & El PDF se compila desde el archivo LaTeX siguiendo el README y no contiene referencias rotas. & Verificar en copia limpia\\
4 & El enlace de GitHub de la portada abre sin autenticación y apunta al repositorio correcto. & Verificar antes del corte\\
5 & La documentación ética y la adenda están vigentes y sus fechas son coherentes con la recolección. & Verificar en \nolinkurl{08_Etica/}\\
6 & El contenedor cifrado abre con la contraseña entregada por el SGA y coincide con las fichas técnicas. & Verificación docente/equipo\\
7 & Todos los videos y audios pasan \texttt{ffprobe} y tienen duración mayor que cero. & Pendiente de prueba técnica\\
8 & Los hashes SHA-256 coinciden con los archivos originales antes del cifrado. & Pendiente de prueba técnica\\
9 & Los consentimientos públicos tienen firma y cédula enmascaradas; los originales solo están cifrados. & Pendiente de revisión visual\\
10 & Ninguna fotografía pública conserva rostros identificables ni coordenadas GPS. & Pendiente de revisión EXIF\\
11 & El cuestionario conserva respuestas literales y elimina columnas identificativas. & Pendiente de revisión\\
12 & Las transcripciones cubren el 100\% de entrevistas declaradas y no incluyen nombres propios. & Pendiente de cotejo\\
13 & Existen al menos tres documentos organizacionales de tipos distintos, debidamente anonimizados. & Pendiente de cotejo\\
14 & Existen tres walkthroughs verificables, con actas enmascaradas y grabaciones restringidas. & Pendiente de cotejo\\
15 & La codificación temática cubre todas las transcripciones y permite reconstruir la derivación de requisitos. & Pendiente de revisión\\
16 & Los 25 RF, 15 RNF y 10 CU detallados se relacionan con evidencia y artefactos reales. & Verificar trazabilidad\\
17 & La matriz de trazabilidad contiene al menos 40 filas y todas las columnas exigidas. & 48 filas preparadas; revisar\\
18 & Los diagramas UML y mockups tienen fuente editable y exportaciones PNG/SVG. & Pendiente de inserción\\
19 & El MVP se ejecuta desde una copia limpia y cubre al menos el 60\% de RF Must. & Pendiente de ejecución\\
20 & El protocolo OSF tiene fecha anterior a las actividades experimentales pendientes. & Pendiente\\
21 & Los scripts reproducen las tablas y figuras sin alterar los datos crudos. & Pendiente\\
22 & El paquete Zenodo contiene únicamente información anonimizada y tiene diccionario de datos. & Pendiente\\
23 & La versión final está asociada a un commit, una etiqueta y el historial correspondiente. & Pendiente de cierre\\
\end{longtable}

\subsection{Constancia de revisión}

\begin{table}[H]
\centering
\caption{Registro de revisión final del ERS/SRS y repositorio}
\label{tab:constancia-revision}
\small
\begin{tabularx}{\textwidth}{L{3.60cm} Y}
\toprule
\textbf{Dato} & \textbf{Registro}\\
\midrule
Versión del ERS/SRS & \VersionDocumento\\
Commit evaluado & \rule{0.68\linewidth}{0.4pt}\\
Fecha de revisión & \rule{0.68\linewidth}{0.4pt}\\
Analista líder & \Anderson\\
Revisión del equipo & \Mayra\\
Resultado general & \rule{0.68\linewidth}{0.4pt}\\
Observaciones & \vspace{1.10cm}\\
\bottomrule
\end{tabularx}
\end{table}

\end{document}
